\documentclass[12pt, oneside]{book}
\usepackage[margin=1in]{geometry}
\usepackage{microtype}
\usepackage{amsmath, amssymb, amsthm}
\usepackage[authoryear, round]{natbib}
\usepackage{hyperref}
\hypersetup{colorlinks=true, linkcolor=black, citecolor=black, urlcolor=black}
\usepackage{booktabs}
\usepackage{array}
\usepackage{tikz-cd}
\usepackage{fancyhdr}
\setlength{\headheight}{26pt}
\addtolength{\topmargin}{-12pt}
\pagestyle{fancy}
\fancyhf{}
\fancyhead[L]{\small\itshape Continuations Before Objects}
\fancyhead[R]{\small\itshape\leftmark}
\fancyfoot[C]{\thepage}
\renewcommand{\headrulewidth}{0.4pt}

\newtheoremstyle{flyxion}
  {9pt}{9pt}{\itshape}{0pt}{\bfseries}{.}{0.5em}{}
\theoremstyle{flyxion}
\newtheorem{theorem}{Theorem}[chapter]
\newtheorem{lemma}[theorem]{Lemma}
\newtheorem{proposition}[theorem]{Proposition}
\newtheorem{corollary}[theorem]{Corollary}
\theoremstyle{definition}
\newtheorem{definition}[theorem]{Definition}
\newtheorem{hypothesis}[theorem]{Hypothesis}
\newtheorem{conjecture}[theorem]{Conjecture}
\newtheorem{principle}[theorem]{Principle}
\newtheorem{remark}[theorem]{Remark}
\newtheorem{openproblem}[theorem]{Open Problem}

\title{\textbf{Continuations Before Objects}\\[0.6em]
       \large A Geometric Theory of Intelligent Persistence}
\author{Flyxion\\[0.4em]\small Independent Researcher}
\date{2026}

\begin{document}

\maketitle

\chapter*{Abstract}
\addcontentsline{toc}{chapter}{Abstract}

Classical theories of intelligence begin with states and derive relations
between them. This monograph begins from the opposite end. The primary
objects are continuations: admissible future trajectories through
constrained possibility spaces. States, representations, meanings,
memories, and agents are all derived from continuation structure rather
than presupposed by it.

The argument proceeds in three movements. The first establishes that
prediction-primary architectures systematically fail because they optimize
representations rather than preserve reachable futures, and introduces
HYDRA (Hybrid Dynamic Reasoning Architecture) as the first concrete
realization of the alternative ontology. The second introduces the
Generative Compression Hypothesis and its formal centerpiece, the
Generative Sufficiency Principle: the minimal generator $G^*$ satisfying
$R(G^*) = R(I)$ is the correct preservation target, not the inventory
$I$ itself. The third formalizes both claims through RSVP field theory,
stratified semantic manifolds, sheaf-theoretic semantics, and the
category of admissible states, showing that HYDRA, restricted
mind-change learning, memory dynamics, social trajectory geometry, and
cosmological field structure are all specializations of the quartet
$(\mathcal{X}, \Pi, \mathcal{A}, P)$.

The Continuation Principle, stated in Chapter~3, is the philosophical
center: two states are semantically equivalent if and only if they induce
identical admissible future structure. From this principle the Minimal
Projection Theorem, the hallucination obstruction, the projection
collapse theorems, and the Structural Universality Conjecture all
descend. The open problems are collected and ordered by dependency in
Part~VI. The monograph closes with the Reachability Thesis: reality is
organized by continuation structure, and what must survive for the future
to remain reconstructible is not an inventory but a generator.

\chapter*{Preface}
\addcontentsline{toc}{chapter}{Preface}

This monograph is the synthetic statement of a research program whose
components have accumulated across many documents, frameworks, and
working papers. The program is unified by one philosophical commitment:
reachability is more fundamental than representation, and what a system
can still become matters more than what it currently is.

HYDRA serves throughout as an operational case study rather than as the
primary subject. The primary subject is the Admissibility Program
itself, of which HYDRA is the most fully developed computational
instantiation. Readers already familiar with the corpus will recognize
material from RSVP, CLIO, Spherepop, MEM\textbar{}8, Repair Theory, the
Eight-Letter Keyboard, and the broader reachability framework. The
purpose of this monograph is to show that those frameworks are not
adjacent projects but instances of a single abstract schema organized by
the question: \emph{what generators are worth preserving?}

A statement classification convention is used throughout. Results drawn
from the established mathematical literature are identified as such.
Framework-specific definitions and interpretations are distinguished from
independent results. Claims that are conjectural---falsifiable in
principle but not yet proved---are explicitly marked as conjectures. The
convention is a philosophical commitment as much as an editorial one: a
framework that cannot distinguish its theorems from its conjectures is
not yet a theory.

The monograph is organized as a sequence of nested questions. Part~I
asks what is wrong with prediction. Part~II asks what survives
compression. Part~III asks what mathematical structures describe
survival. Part~IV asks what domains instantiate those structures.
Part~V asks whether the theory survives implementation. Part~VI asks
what has actually been proved. Part~VII asks what follows if the
framework is correct. Each part compresses its predecessor; the book is
itself an instance of the compression principle it describes.

\tableofcontents

%% ==============================================================
\part{The Failure of Prediction-Primary Architectures}
%% ==============================================================

%% --------------------------------------------------------------
\chapter{The Problem with Prediction}
%% --------------------------------------------------------------

The question most theories of intelligence set out to answer is: given
what has been observed, what is most likely to be true, and what action
does that make optimal? This question is natural. It is also, this
monograph argues, the wrong question to make primary. Not because
prediction is useless---prediction serves important purposes and appears
as a component in the framework developed here---but because making it
primary produces an architecture whose failures are systematic and
principled rather than incidental and correctible. Understanding why
requires looking carefully at what prediction-primary architectures
implicitly assume about the relationship between states, representations,
and futures.

\section{The Dominant Paradigm}

A prediction architecture, in the sense intended here, is any system
whose design criterion is the minimization of prediction loss. The
system receives inputs, constructs a representation of the current state,
and uses that representation to select actions or generate outputs. The
quality of the representation is measured by its predictive accuracy: a
representation is better to the extent that it supports more accurate
predictions about what will happen next or what actions will prove
correct.

\begin{definition}[Prediction Architecture]
A prediction architecture is a tuple $P = (X, Y, \pi, L)$ where $X$ is
a state space, $Y$ is an action or output space, $\pi : X \to Y$ is a
prediction map, and $L : Y \times Y \to \mathbb{R}_{\geq 0}$ is a
prediction loss measuring deviation from the target output.
\end{definition}

This formalization is deliberately general. It covers statistical
learning systems whose prediction map is estimated from data, planning
systems whose prediction map encodes transition models, and large
language models whose prediction map assigns probabilities to continuations.
In each case the operative criterion is the same: minimize $L$ over the
space of possible prediction maps $\pi$.

The paradigm is not merely a theoretical idealization. It governs the
design of most contemporary AI systems, from recommender systems
optimizing click-through rates, to language models trained on next-token
prediction, to reinforcement learning agents maximizing expected reward.
The differences between these systems are differences in the specific
form of $X$, $Y$, $\pi$, and $L$. The shared structure is the
prediction-first criterion.

\section{What Prediction-Primary Architectures Cannot See}

The central limitation of prediction-primary architectures follows
directly from their definition. If the quality of a representation is
measured by its predictive accuracy, then any distinction that does not
affect prediction quality is invisible to the system's design criterion.
The system is, by construction, blind to whatever it cannot predict from.

\begin{theorem}[Prediction Collapse]
\label{thm:pred_collapse}
If two states satisfy $\pi(x_1) = \pi(x_2)$---that is, if they produce
identical optimal actions---then every prediction-optimal learner
identifies them: $x_1 \sim_\pi x_2$ under optimal prediction.
\end{theorem}

\begin{proof}
The prediction loss $L(\pi(x), y^*)$ depends on the predicted output
$\pi(x)$ and the target $y^*$, not on $x$ directly. If
$\pi(x_1) = \pi(x_2)$, then $L(\pi(x_1), y^*) = L(\pi(x_2), y^*)$ for
every target $y^*$. Any distinction between $x_1$ and $x_2$ that does
not alter the prediction contributes nothing to the loss and is therefore
projected away by the optimal $\pi$. Hence $x_1 \sim_\pi x_2$.
\end{proof}

Theorem~\ref{thm:pred_collapse} is elementary, but its consequences are
not. The states a prediction-primary system distinguishes are exactly the
states whose differences affect its current prediction task. Everything
else collapses. This is appropriate when the prediction task is a
complete proxy for what the system needs to do. It becomes a liability
when the prediction task leaves out something important---specifically,
when it leaves out the structure of futures that are not yet being
predicted.

The problem is not that prediction-primary systems make incorrect
predictions. The problem is that the distinctions they preserve are
determined by the predictions they currently make, not by the futures
they might need to navigate. Two states that support identical current
predictions may have radically different future continuation structures:
one may lead to a rich range of admissible futures while the other
leads to a narrow corridor or a dead end. Prediction collapse identifies
them. Any information about the difference in their future potential is
lost.

This is not a fixable bug. It is a consequence of the design criterion.
To fix it would require changing what the system is optimizing for---
which is exactly the move made in Chapter~3.

\section{Four Manifestations of the Failure}

The theoretical argument of the previous section can be made concrete
by observing that it reappears, independently formulated, across four
distinct research programs. Each program identifies the same gap from
a different direction.

Personalized recommendation systems, formalized in the PERSCEN
framework, discover that optimizing for predicted preference across
scenarios collapses distinctions between users who have similar average
preferences but different preference dynamics. The prediction surface
is the same; the trajectory through preference space is not. Systems
optimized on the prediction surface inherit the collapse.

Embodied cognition research, formalized in Relevance Activation Theory,
finds that cue-driven gradient flows over relevance fields produce
behavior that cannot be captured by any prediction map defined over
static states, because the relevant object is the flow itself: the
direction and curvature of the trajectory through the cue space, not
the position at any given moment.

Causal memory research, formalized in the Chain of Memory framework,
identifies the limitation at the level of reasoning transparency.
Prediction-primary systems can be accurate without being causally
interpretable, because causal structure and predictive structure come
apart. A system can predict correctly for the wrong reasons, and a
system that predicts correctly for the wrong reasons will fail when
those reasons cease to hold---precisely when interpretability and causal
traceability would matter most.

Field-theoretic semantic modeling, formalized in RSVP and TARTAN,
encounters the limitation at the representational level. Semantic
content, on the field-theoretic account, is not a property of a
representation but a property of the relationship between a
representation and the space of future representations it can reach.
Any architecture that evaluates representations by their predictive
accuracy rather than by their future-preserving capacity is, on this
account, measuring the wrong thing.

These four research programs are not in communication with each other
in any direct sense. They arrive at the same diagnosis independently.
The convergence is evidence that the diagnosis is tracking something
real rather than a parochial concern of any single framework.

\section{The Replacement Question}

The failure of prediction-primary architectures, understood precisely,
points directly toward the replacement question. If the problem is that
prediction collapse destroys future-relevant distinctions, the solution
is an architecture whose design criterion explicitly preserves
future-relevant structure.

The replacement question is therefore not: \emph{what is most likely?}
The replacement question is: \emph{what futures remain reachable?}

The shift appears subtle. It is not. A system that asks what is most
likely optimizes a probability distribution over outcomes. A system that
asks what futures remain reachable maintains an admissibility structure
over trajectories. These are different mathematical objects, and they
produce different architectures, not better versions of the same
architecture.

Prediction becomes, in the reachability-first framework, a useful tool
for identifying which futures are more or less likely to be reachable.
It is not the criterion by which the quality of a representation is
judged. The criterion is whether the representation preserves the
structure of admissible futures: whether what can still be reached from
the represented state is the same as what could be reached from the
original state.

\section{Why the Shift Produces a Different Architecture}

The practical consequence of the shift from prediction to reachability
can be stated simply. In a prediction-primary architecture, the question
asked at every design decision is: does this change improve predictive
accuracy? In a reachability-primary architecture, the question asked at
every design decision is: does this change preserve admissible future
structure?

These criteria pull in different directions. A prediction-primary
system discards features that do not predict well; a
reachability-primary system retains features whose loss would eliminate
admissible futures even if those features are not currently predictive.
A prediction-primary system represents the current state as accurately
as possible; a reachability-primary system represents the geometry of
the space surrounding the current state, because that geometry
determines what can still be reached. A prediction-primary system
improves by acquiring better probability estimates; a
reachability-primary system improves by acquiring better models of which
transitions are admissible.

The von Neumann architecture is the canonical engineering instantiation
of prediction-primary design. Its fetch-execute cycle retrieves a static
state from a fixed memory address, processes it under a globally
synchronized clock, and returns a result. The global clock imposes a
temporal projection $\pi_t : \mathcal{X} \to \mathcal{X}_t$ that
identifies all physical states within a clock cycle, discarding
mid-cycle dynamics as irrelevant. The bottleneck this creates---the
constant shuttling of static snapshots between processor and memory---is
the physical cost of treating memory as indexed storage rather than as
a navigable field. The thermodynamic inefficiency is not incidental; it
is a structural consequence of the snapshot ontology.

The reachability-primary alternative, instantiated in HYDRA at the
cognitive level and in continuation-first hardware at the physical level
(Chapter~20), replaces indexed retrieval with field navigation: a query
perturbs the field, which resonates toward stored configurations rather
than scanning through them. The architecture is different not because
it is more efficient at prediction but because it is organized around
a different primitive.

The HYDRA architecture, introduced in Chapter~3, is the operational
realization of the reachability-primary criterion. It is presented there
not as a response to the engineering limitations of specific existing
systems, but as the natural computational architecture for a system
whose design criterion is the preservation of admissible future
structure. The four source frameworks whose shared critique occupies
the next chapter are the theoretical landscape HYDRA inhabits, not
the historical origin of its design.

%% --------------------------------------------------------------
\chapter{Four Source Frameworks and Their Shared Critique}
%% --------------------------------------------------------------

Four research programs, developed with different applications and
different formal vocabularies, converge on the same structural
diagnosis: prediction-primary architectures preserve the wrong things.
This chapter presents each framework in sufficient detail to make the
convergence visible, then formalizes it as a shared quotient
structure. The frameworks are not components of HYDRA; they are the
theoretical landscape that makes HYDRA's design decisions intelligible.

\section{PERSCEN: Personalization Across Scenarios}

The PERSCEN framework addresses the problem of multi-scenario matching
in large-scale recommendation: a system must recommend appropriately
across heterogeneous contexts---homepage feeds, search results, product
pages---while preserving something of the individual user's preference
structure across all of them. The naive approach, training separate
models for each scenario, fails to transfer what is stable about a
user's preferences; the naive unified approach averages across scenarios
and loses what is specific.

PERSCEN's response is to construct a user-specific feature graph,
encoding not a static preference vector but a relational structure over
features whose edges represent interaction patterns. The graph is not
a knowledge graph in the usual sense: it does not record facts about
the user. It records the adjacency structure of the user's preference
space---which features co-activate, which transitions are typical,
which regions of preference space are reachable from which others.

The PERSCEN projection $\Pi_P : X \to M_P$ maps from the full state
space of user-item interactions to a personalized preference manifold.
The manifold is not defined by what the user currently prefers. It is
defined by the geometry of the user's preference dynamics: the
directions along which preferences are stable, the directions along
which they vary across scenarios, and the boundaries beyond which
preference predictions become unreliable.

\section{Relevance Activation Theory: Embodied Cue Dynamics}

Relevance Activation Theory (RAT) begins from the observation that
behavior is not well described as the selection of an action from a
menu of options, evaluated by a utility function defined over outcomes.
Behavior is better described as a flow: a trajectory through a space
of possible states, guided by the gradient of a relevance field that
is itself modified by the trajectory.

In RAT, each cue $c \in C$ induces a relevance field $\rho_c : X \to
\mathbb{R}_{\geq 0}$ over the agent's state space. The field is not a
static evaluation function; it is a dynamical landscape whose shape
depends on context, history, and the agent's current position. Behavior
emerges from gradient flow over this landscape:
\[
\frac{dx}{dt} = \nabla \rho_c(x).
\]
The trajectory, not the position, is the primary object. Two agents at
the same position who have arrived there by different paths may face
different local gradient structures---because the cue fields they have
activated, and the fields that have been modified by their passage,
differ.

The RAT projection $\Pi_R : X \to M_R$ maps from the full state space
to a cue activation manifold. The manifold encodes not the current
activation level of each cue but the gradient structure of the
relevance field at the current position: what the agent is moving
toward, at what rate, and with what curvature.

\section{Chain of Memory: Causal Traceability}

The Chain of Memory framework is motivated by a different failure mode:
systems that predict correctly but cannot be audited. A memory system
is causally faithful if the influence of past states on current outputs
can be traced through an explicit causal chain. Most deep learning
systems are not causally faithful in this sense: their outputs depend
on their internal states in ways that cannot be decomposed into
interpretable causal pathways.

Chain of Memory formalizes causal faithfulness as a constraint on the
memory stack. Memory states $M_i \in \mathbb{R}^d$ evolve according to
a differentiable operator $M_{i+1} = \phi(M_i, u_i, c_i)$, and causal
influence is measured by
\[
I(M_i \to y) = \left\|\frac{\partial y}{\partial M_i}\right\|.
\]
A memory state is causally influential if small changes to it produce
measurable changes in the output. A system is causally faithful if
every output-influencing memory state can be traced through the update
chain to specific past inputs.

The CoM projection $\Pi_C : X \to M_C$ maps from the full state space
to a causal trace manifold. The manifold encodes not what the system
currently believes but the history of influences that produced the
current state: the causal pathway, not the endpoint.

\section{RSVP and TARTAN: Field-Theoretic Semantic Recursion}

The RSVP framework begins from a different starting point entirely.
Rather than asking how to improve a specific architectural component,
it asks what the correct ontology for semantic content is. Its answer:
semantic content is not a property of representations but a property
of the relationship between representations and the spaces of future
representations they can reach.

This is formalized through the RSVP field triple $(\Phi, v, S)$ over
a smooth manifold $M$, where $\Phi : M \times \mathbb{R} \to \mathbb{R}$
is the scalar accessibility potential, $v : M \times \mathbb{R} \to TM$
is the admissible flow field, and $S(x,t) = \log|\mathcal{A}(x,t)|$ is
the entropic accessibility measure encoding the volume of admissible
futures available from state $x$ at time $t$. The full field equations
and their properties occupy Chapter~11; here it suffices to note that
the RSVP field is not a static assignment of meaning to representations
but a dynamical structure encoding how meaning flows, accumulates, and
dissipates through semantic space.

TARTAN extends RSVP by providing the stratification structure:
the semantic manifold is decomposed into tiles, each approximately
uniform in its RSVP field values, whose recursive refinements capture
hierarchical semantic organization. The TARTAN projection
$\Pi_S : X \to M_S$ maps from the full state space to the RSVP field
manifold, encoding the local accessibility geometry rather than the
current representational content.

\section{The Convergence Argument}

The four projections $\Pi_P$, $\Pi_R$, $\Pi_C$, $\Pi_S$ are defined
differently, formalized in different mathematical vocabularies, and
motivated by different application domains. They nevertheless share
a single structural property: each is a quotient construction that
identifies states indistinguishable with respect to the framework's
operative criterion.

\begin{definition}[Projection System]
A projection system over $X$ is a family $\Pi = \{\pi_i : X \to M_i\}$
of continuous surjections onto lower-dimensional manifolds satisfying
mutual consistency on overlapping domains.
\end{definition}

\begin{theorem}[Shared Quotient Structure]
\label{thm:shared_quotient}
Each of the four source projections is a quotient construction
identifying states that are indistinguishable with respect to the
framework's operative criterion.
\end{theorem}

\begin{proof}
For $\Pi_P$: PERSCEN identifies users whose preference dynamics are
identical, that is, whose personalized feature graphs have the same
adjacency structure. Two users who arrive at the same preference
position by different trajectories are distinguished if and only if
their graphs differ. The projection $\Pi_P$ is the quotient by the
equivalence relation of identical preference dynamics.

For $\Pi_R$: RAT identifies states that produce identical cue
activation patterns and identical gradient structures at the current
position. The projection $\Pi_R$ is the quotient by the equivalence
relation of identical relevance field geometry.

For $\Pi_C$: Chain of Memory identifies states whose causal traces
are identical, that is, states whose history of memory updates has
produced the same pattern of causal influences on the current output.
The projection $\Pi_C$ is the quotient by the equivalence relation of
identical causal pathways.

For $\Pi_S$: RSVP identifies states with identical accessibility
potential, flow direction, and entropic accessibility---the same local
RSVP field triple. The projection $\Pi_S$ is the quotient by the
equivalence relation of identical field geometry.

In each case the projection is the canonical surjection to the
equivalence class space under the operative criterion.
\end{proof}

What Theorem~\ref{thm:shared_quotient} reveals is that the four
frameworks are not pursuing different goals. They are constructing
different quotient spaces over the same underlying phenomenon: the
gap between what a state currently is and what it can still become.
PERSCEN's preference dynamics, RAT's gradient structure, CoM's causal
trace, and RSVP's field geometry are four different attempts to capture
the same thing---the trajectory structure that survives in the
representation after projection.

The operative criterion all four frameworks are converging on is
admissibility: the structure of transitions a state permits. Chapter~3
introduces this criterion directly and builds the architecture that
makes it primary.

%% --------------------------------------------------------------
\chapter{The Reachability Turn}
%% --------------------------------------------------------------

The previous two chapters have established a diagnosis and a
convergence. The diagnosis: prediction-primary architectures preserve
the wrong distinctions, collapsing states that differ in their future
potential. The convergence: four independent research programs identify
this failure from different directions, each constructing a quotient
space that attempts to preserve something of trajectory structure
rather than merely current state.

This chapter introduces the reachability-primary alternative directly.
It does not trace a historical development; it presents the
architecture in its current geometric form and explains what
design commitments that form embodies.

\section{The Ontological Inversion}

The shift from prediction-primary to reachability-primary design can be
stated as a sequence of three inversions, each more radical than the
last.

The first inversion is from states to trajectories. In a
prediction-primary architecture, the fundamental object is the current
state $x \in X$. Trajectories are sequences of states, useful for
defining temporal prediction tasks but not ontologically primary.
In the reachability framework, the fundamental object is the trajectory
$\gamma : [0,1] \to X$, and states are positions along trajectories.
What a state is, is determined by where it can go next---by its
position in trajectory space, not by its intrinsic properties.

The second inversion is from representations to continuations. A
prediction-primary architecture evaluates representations by their
accuracy: how well does the representation encode the current state?
The reachability framework evaluates representations by their
continuation-preserving capacity: does the representation retain the
information needed to determine what admissible continuations remain
available from the current position? A representation that is accurate
but continuation-destroying is worse, by the reachability criterion,
than a representation that is less accurate but continuation-preserving.

The third inversion is from optimization to admissibility. A
prediction-primary architecture improves by optimizing: finding the
prediction map $\pi$ that minimizes loss $L$. The reachability
framework imposes admissibility as a prior constraint on the space of
possible trajectories and then, within that constrained space, may
apply optimization criteria. Admissibility is not a constraint added
to an optimizer; it is the primary structure within which everything
else is defined.

\section{Two Design Assumptions Contrasted}

The six-component HYDRA architecture can be read under two different
design assumptions, which produce different accounts of what it does
and why it works.

Under the engineering design assumption, the six components are modules
that process information in sequence. Each module receives the output
of the previous one, transforms it, and passes it on. The architecture
is a pipeline, evaluated by whether the final output is correct. The
categorical language in which the components are described is
organizational convenience.

Under the geometric design assumption, the six components are
admissibility-preserving projections. Each projection maps from a
larger space to a smaller one, eliminating distinctions that make no
difference to admissible future continuations while attempting to
preserve those that do. The architecture is a chain of quotient
constructions, and its correctness condition is not output accuracy
but continuation preservation. The categorical language is not
organizational convenience; it is the natural description of the
objects in question.

This monograph operates under the geometric design assumption. The
engineering description is not wrong, but it is less informative. The
geometric description reveals properties of the architecture that are
invisible at the engineering level---in particular, it reveals the
conditions under which two very different implementations are
equivalent (they realize the same functor), and the conditions under
which an implementation has failed (it has not preserved admissibility
through some stage of the projection chain).

\section{The HYDRA Pipeline as a Chain of Projections}

The HYDRA architecture is:
\[
H = \mathrm{GLU}_\mathrm{RSVP} \circ M \circ T \circ F_a \circ G_a
    \circ R.
\]
Under the geometric design assumption, this is a sequence of six
admissibility-preserving projections through a sequence of spaces:
\[
X \;\xrightarrow{R}\; X_R \;\xrightarrow{G_a}\; \mathcal{G}
  \;\xrightarrow{F_a}\; \mathcal{G}_a \;\xrightarrow{T}\; \Gamma
  \;\xrightarrow{M}\; \Gamma_M \;\xrightarrow{\mathrm{GLU}}\; Y.
\]
Each stage narrows the possibility space by a different admissibility
criterion. Together they constitute a machine for repeatedly computing
the minimal admissible quotient of reality while preserving the largest
reachable future.

\subsection{Recognition ($R$): Which Distinctions Deserve Attention?}

The raw world state $x \in X$ contains vastly more information than
can be processed or is relevant to any particular cognitive task. The
recognition functor $R : X \to X_R$ constructs an activated subspace
$X_R \subset X$ by identifying which features participate in
potentially important trajectories.

In the original formulation, each cue $c \in C$ induces a Gaussian
relevance field:
\[
\rho_c(x) = \exp\!\left(-\tfrac{1}{2}(x - \mu_c)^\top \Sigma_c^{-1}
            (x - \mu_c)\right),
\]
and the gradient flow $\dot{x} = \nabla\rho_c(x)$ governs attention.
The functor $R : \mathbf{Cue} \to \mathbf{Field}$ maps cues to
relevance fields, and the activated subspace consists of those features
participating in the gradient flow. Features that participate in no
admissible trajectory under any relevant cue are projected away.

The admissibility criterion at this stage is participation: a
distinction is worth preserving if its presence or absence affects
which trajectories are available.

\subsection{Graph Construction ($G_a$): Where Can I Still Go?}

The activated subspace $X_R$ contains distinguished features but no
relational structure. The graph construction functor
$G_a : X_R \to \mathcal{G}$ organizes those features into a
reachability graph $\mathcal{G} = (V, E)$ whose nodes are activated
features and whose edges represent possible transitions. This is not a
knowledge graph encoding what is true; it is a reachability graph
encoding what is reachable.

The user-specific feature graph, drawn from PERSCEN, provides the
concrete construction: the adjacency matrix $[A_a^{(1)}]_{m,:} =
\mathrm{MLP}_m([e_{a,1}, \ldots, e_{a,N_f}, \mathrm{one\text{-}hot}(m)])$
encodes the personalized neighborhood structure of the feature space.
A graph neural network computes higher-order interactions, and the
resulting representation $h_a^{(L)}$ encodes the connectivity structure
of the agent's accessible feature space rather than the features
themselves.

The admissibility criterion at this stage is connectivity: a transition
is represented if it is possible given the current agent and context.

\subsection{Admissibility Filtering ($F_a$):
            Which Continuations Survive?}

The reachability graph $\mathcal{G}$ contains many nodes that, while
reachable, lead nowhere useful---dead ends, cycles that exhaust
admissibility, or regions whose future potential is beneath a useful
threshold. The admissibility filtering functor
$F_a : \mathcal{G} \to \mathcal{G}_a$ removes inadmissible regions:
\[
\mathcal{G}_a = \bigl\{v \in \mathcal{G} \mid
\mathrm{Vol}(\mathcal{A}(v)) > \varepsilon\bigr\}.
\]
Nodes from which the volume of admissible future continuations is
insufficient are removed. This is where CLIO-style reasoning appears
in the architecture: rather than maximizing prediction, the stage
minimizes irreversible collapse, retaining only those regions that
preserve sufficient future reachability.

The admissibility criterion at this stage is future volume: a node
survives if the space of admissible continuations from it is large
enough to be worth preserving.

\subsection{Trajectory Formation ($T$): What Paths Remain?}

The filtered graph $\mathcal{G}_a$ is still a static structure. The
trajectory formation functor $T : \mathcal{G}_a \to \Gamma$ lifts
the static graph to a space of trajectories $\Gamma$, where a
trajectory is a path:
\[
\gamma : [0,1] \to \mathcal{G}_a.
\]
This is the stage at which the ontological shift from states to
trajectories becomes operational. The primitive object of all subsequent
processing is no longer a node $v$ but a path $\gamma$ through the
admissible graph. The TARTAN tiling, introduced formally in Chapter~12,
provides the stratification structure within which these paths are
defined: trajectories within a tile are uniform in RSVP field values,
and admissible transitions between tiles are governed by the
tangent-cone condition at stratum boundaries.

The admissibility criterion at this stage is path admissibility: a
trajectory survives if it respects the transition constraints of the
stratified manifold at every step.

\subsection{Memory ($M$): Which Paths Have Persisted Before?}

Memory is commonly understood as storage: a record of past states that
can be retrieved and used. Within HYDRA, memory is trajectory
stabilization: the modification of the trajectory space $\Gamma$ by
residue from prior trajectories, producing a residue-weighted space
$\Gamma_M$.

The memory functor $M : \Gamma \to \Gamma_M$ implements this through
a latent stack $\{M_i\}$ whose states evolve according to
$M_{i+1} = \phi(M_i, u_i, c_i)$, where $u_i$ is the current
user or agent state and $c_i$ is the current cue embedding. Past
trajectories leave persistent curvature in $\Gamma_M$: repeated paths
become attractors, rarely traversed regions become energetically
costly, and the metric on trajectory space is modified by the
accumulated residue of prior motion.

In RSVP language, memory is field residue: past field configurations
leave persistent modifications to $\Phi$, $v$, and $S$ that bias
future trajectories toward familiar regions. In geometric language,
memory modifies the metric on $\Gamma$, making some paths shorter
(more available) and others longer (more costly) based on history.

The admissibility criterion at this stage is historical weighting:
paths that have been traversed and confirmed admissible have higher
effective availability; paths that have been attempted and found
inadmissible have their availability reduced.

\subsection{RSVP Coupling ($\mathrm{GLU}_\mathrm{RSVP}$): Which Path
             Best Preserves Future Structure?}

The final stage selects a trajectory from $\Gamma_M$ by evaluating
each candidate against the RSVP field triple $(\Phi, v, S)$:
\[
\gamma^* = \arg\max_{\gamma \in \Gamma_M}
\int_\gamma \bigl(\alpha\,\Phi + \beta\,\|v\| + \delta\,S\bigr)\,dt.
\]
This is not prediction: the trajectory is not selected because it is
most probable. It is navigation: the trajectory is selected because
it occupies the most favorable position in the admissibility field,
passing through regions of high accessibility potential $\Phi$,
strong admissible flow $v$, and high entropic accessibility $S$.

The selection criterion integrates three components of the RSVP field
along the entire trajectory, not just at the current position. A
trajectory that passes through a favorable region briefly but then
enters an admissibility basin is evaluated differently from one that
maintains favorable field values throughout. The architecture is
sensitive to the global structure of the path, not merely its
instantaneous properties.

\section{Six Successive Restrictions on a Possibility Space}

The architecture as a whole performs a progressive narrowing of
the possibility space through six admissibility criteria applied in
sequence:

Recognition asks which distinctions exist. Graph construction asks
how they are connected. Admissibility filtering asks which
continuations survive with sufficient future volume. Trajectory
formation asks what paths remain through the constrained graph.
Memory asks which paths have been historically confirmed. RSVP
coupling asks which remaining path best preserves future structure
according to the field geometry.

Each stage reduces the space. The reduction at each stage is governed
by a different admissibility criterion, but all six criteria are
instances of the same underlying question: what can be eliminated
without losing information relevant to admissible future continuations?
The answer at each stage is different because the space being
filtered and the criterion being applied are different. But the
question is the same.

\section{The Strong Interpretation and the Continuation Principle}

The geometric reading of HYDRA supports a strong interpretation
that connects it directly to the Minimal Projection Theorem
(Chapter~10) and the broader Admissibility Program.

Under the strong interpretation, the chain
\[
X \to X_R \to \mathcal{G} \to \mathcal{G}_a \to \Gamma \to \Gamma_M
  \to Y
\]
is a sequence of admissibility-preserving projections, and the
Minimal Projection Theorem guarantees that each projection is, in
the appropriate sense, the minimal one compatible with preserving
the admissibility structure relevant at that stage. HYDRA is not
merely a system that happens to use reachability-friendly components.
It is a machine for repeatedly computing the minimal admissible
quotient of reality while preserving the largest reachable future.

This interpretation motivates and is formalized by the Continuation
Principle, which is stated here for the first time and will be the
reference point for all subsequent chapters.

\begin{principle}[Continuation Principle]
\label{prin:continuation}
Two states are semantically equivalent if and only if they induce
identical admissible future structure:
\[
x_1 \equiv_{\mathcal{A}} x_2 \iff \mathcal{A}(x_1) = \mathcal{A}(x_2).
\]
Meaning is not a property of states. Meaning is a property of the
admissible continuation structure states induce.
\end{principle}

\subsection*{States as Coordinates, Not Ontological Primitives}

A natural objection arises: to define $\mathcal{A}(x)$, one must
already have a state $x$ from which admissible futures are measured.
Does the framework therefore presuppose states as primitives,
contradicting its claim to derive objects from continuations?

The answer is that states function here as \emph{coordinates on an
admissibility manifold}, not as ontological primitives. The distinction
is the same one that differential geometry draws between a coordinate
chart and the manifold it charts. A coordinate system exists, is
useful, and is necessary for local computation---but it is not the
manifold. Different coordinate charts cover the same manifold;
coordinates are conventional and the manifold is not. The framework
uses state labels to compute admissibility structures; it is those
structures, not the labels, that constitute the object.

\begin{proposition}[State Reconstruction Up to Continuation
Equivalence]
\label{prop:state_recon}
States are identifiable only up to continuation equivalence: if
$\mathcal{A}(x) = \mathcal{A}(y)$, then $x$ and $y$ are
indistinguishable within the framework and are identified as the same
point in the semantic space $M^* = \mathcal{X}/{\sim_{\mathcal{A}}}$.
\end{proposition}

\begin{proof}
By the Minimal Projection Theorem (Theorem~10.1), $\pi^*(x) =
\pi^*(y)$ whenever $\mathcal{A}(x) = \mathcal{A}(y)$.
\end{proof}

The framework does not deny that states exist. It denies that states
are fundamental. Two states with identical continuation structures are
the same object within the framework regardless of any other
properties they may share or differ in. The coordinate $x$ is a handle
for computation; the continuation structure $\mathcal{A}(x)$ is the
object. This is the operational content of the title \emph{Continuations
Before Objects}: not that states are illusions, but that they have
the status of coordinate charts rather than that of primary entities.

The Continuation Principle is the philosophical center of this
monograph. The Minimal Projection Theorem (Chapter~10) is its
mathematical formalization. Every chapter of Part~III establishes
a piece of the mathematical infrastructure needed to work with it
precisely. Every chapter of Part~IV is a consequence or application.
The Reachability Thesis (Chapter~21) is its philosophical
culmination. The title \emph{Continuations Before Objects} is its
compressed statement.

\bigskip
\noindent\emph{What follows.} Part~II develops the Generative
Compression Hypothesis: the claim that what must be preserved to
maintain admissible future structure is not an inventory of states
but a minimal generator from which the inventory can be reconstructed.
The argument begins with the empirical observation that this pattern
recurs across a wide range of apparently unrelated frameworks, then
abstracts from that observation to a formal principle, and then shows
that HYDRA is its computational realization.



%% ==============================================================
\part{The Geometry of Continuations}
%% ==============================================================

\noindent\textit{Part~I established the failure of prediction-primary
architectures and introduced HYDRA as the reachability-primary
alternative, organized by the Continuation Principle: two states are
semantically equivalent if and only if they induce identical admissible
future structure. Part~II asks the next question. If the correct
primitive is the admissible continuation structure rather than the
state or the representation, what must be preserved to maintain that
structure across transformations, compressions, and transmissions? The
answer is not the inventory of states but the generator from which
the inventory can be reconstructed. This part establishes that claim
empirically, through convergent evidence from independent frameworks,
and then formalizes it as the Generative Compression Hypothesis and
its operational principle.}

\bigskip

%% --------------------------------------------------------------
\chapter{Convergence Across Domains}
%% --------------------------------------------------------------

The Continuation Principle, stated at the close of Chapter~3, is
a philosophical claim about what meaning is. Claims of this scope
invite the objection that they are artifacts of a particular
theoretical vocabulary rather than discoveries about the phenomenon
they purport to describe. The strongest response to that objection
is not further argumentation within the vocabulary but evidence from
outside it: independent research programs, working in different
domains with different formalisms and different applications, that
arrive at structurally identical conclusions.

This chapter presents six such programs. They span cosmological field
theory, cognitive architecture, memory dynamics, symbolic systems,
computation theory, and repair epistemology. None of them was
developed in response to the others. Each uses its own vocabulary.
Each is motivated by a different problem. And each independently
arrives at the same structural question: what is the smallest
structure capable of preserving the largest space of future
distinctions?

The convergence is the argument. When the same question recurs
across domains of this variety, it is evidence that the question
is tracking something real---a geometric fact about how systems
maintain the capacity for continuation under transformation---rather
than a parochial concern of any single framework.

\section{RSVP and the Geometry of Future Fields}

The RSVP framework approaches the question of what must be preserved
from the direction of cosmological and cognitive field theory. Its
answer is: the field triple $(\Phi, v, S)$. This is not preserved
because it is large, or because it contains all the information in
the system. It is preserved because it is the generator of future
geometry.

The scalar accessibility potential $\Phi$ encodes where the system
can go with high probability of remaining admissible. The vector
flow field $v$ encodes the preferred directions of admissible
motion. The entropic accessibility measure $S = \log|\mathcal{A}(x,t)|$
encodes how many admissible futures remain available from each
position. Together they do not describe the current state; they
describe the geometry surrounding the current state---the shape of
the space of reachable futures.

The insight of RSVP is that this geometry is the right primitive.
A complete description of the current state is useful only insofar
as it encodes the geometry of what comes next. The field triple is
precisely the minimal encoding of that geometry: the smallest
structure from which the admissibility structure of the immediate
future can be derived. Everything else about the current state that
does not contribute to $(\Phi, v, S)$ is, from the perspective of
future continuation, redundant.

The RSVP question is therefore: given that $(\Phi, v, S)$ is the
correct preservation target, what dynamics govern its evolution,
and what conditions ensure that an admissibility basin---a region
of the field in which future continuations remain stable---persists
over time? Those questions are addressed formally in Chapter~11.
The point here is that RSVP identifies the same abstract pattern
that will recur across the five frameworks that follow: a large
inventory of possible states is generated from a much smaller
generative structure, and what must be preserved is the structure,
not the inventory.

\section{CLIO and Projection Preservation}

The CLIO framework addresses the question of what an epistemically
adequate projection is. Its starting point is the observation that
every representation is a projection: a map from a larger space to
a smaller one that necessarily destroys some distinctions. The
question is which distinctions can be safely destroyed and which
cannot.

CLIO's answer is that the distinctions that cannot be safely destroyed
are those that affect future continuation structure. A projection
$\pi : X \to M$ is epistemically adequate if and only if it preserves
the admissibility structure of $X$: if $\mathcal{A}(x_1) \neq
\mathcal{A}(x_2)$ in $X$, then $\pi(x_1) \neq \pi(x_2)$ in $M$.
A projection that collapses states with different future potentials
is epistemically defective regardless of how accurate it is with
respect to current content.

This is the Minimal Projection Theorem stated as an epistemic
constraint rather than a mathematical fact, and it connects directly
to the question of what must be preserved. The answer CLIO gives
is: the projection. Not the content of the projected space, not
the inventory of states in $M$, but the projection structure itself
---the map that encodes which distinctions in $X$ are
continuation-relevant and which are not. Preserving the projection
is equivalent to preserving the generator of the epistemically
relevant inventory.

\section{Spherepop and Collapse Dynamics}

The Spherepop framework models cognitive structures as bounded regions
---bubbles---with an associated repertoire of operations: pop, refuse,
collapse, bind, repair. The framework is concerned with the dynamics
of cognitive collapse: the conditions under which a cognitive structure
loses the capacity for future continuation, and the operations that
can restore it.

The central observation of Spherepop, relevant here, is that collapse
is not merely the loss of content. A cognitive structure can lose
large amounts of content---many beliefs, many representations, many
activated features---without collapsing, provided the operations that
generate future continuation remain intact. What collapses is not the
inventory but the generative capacity: the ability to pop, bind, and
repair in ways that produce new admissible continuations.

Spherepop therefore identifies the preservation target as the
operation vocabulary, not the content it operates on. The smallest
structure capable of preserving the largest space of cognitive futures
is the minimal set of operations sufficient to generate the repair,
binding, and continuation dynamics that keep the cognitive structure
admissible. The content is derived; the operations are the generator.

This observation connects Spherepop to the collapse risk metric
$C(\tau)$ used in HYDRA's admissibility filtering stage. A
trajectory's collapse risk is not a function of how much content it
carries but of whether it preserves the generative capacity for
future continuation. High content, high collapse risk is worse
than low content, low collapse risk: the generator matters more
than the inventory it has so far produced.

\section{MEM\textbar{}8 and Residue-Based Identity}

The MEM\textbar{}8 framework addresses the question of what memory
is and what identity consists in. Its answer is distinctive and
consequential: identity is not persistence of a substrate but
continuity through residue. A system persists not because the same
material is present at successive times but because the residue
of past trajectories is present in the field that generates future
ones.

In MEM\textbar{}8, a memory state is a wave packet parameterized by
amplitude $\Phi_m$ (encoding scalar accessibility), frequency
(encoding semantic content), phase $v_m$ (encoding associative flow
direction), decay factor $e^{-S_m}$ (encoding entropy: high-entropy
memories decay faster), and an interference term governing how the
packet interacts with others during retrieval. Memory is not stored
as a discrete object; it is encoded as a deformation of the field
that future trajectories navigate.

The consequence is that retrieval and learning are the same
operation. In a system that stores discrete objects, retrieving
a memory and storing a new one are distinct operations with distinct
effects. In a residue field, retrieving a memory by traversing the
path associated with it deepens the deformation: the groove becomes
more pronounced, the attractor stronger, the path more available
to future trajectories. The system learns by remembering and
remembers by learning. Retrieval is reinforcement.

This has an immediate consequence for what must be preserved. The
smallest structure capable of maintaining a rich identity over time
is not the inventory of stored memories but the field deformation
structure: the minimal set of residue patterns sufficient to
reconstruct the memory's continuation structure. The memories
themselves are derived objects; the residue field is the generator.

\section{Repair Theory and the Restoration of Admissibility}

Repair Theory characterizes intelligence not as problem-solving or
prediction but as the ongoing restoration of admissible continuation
structure. A system is intelligent to the extent that it can detect
when its continuation structure has been damaged---when some set of
previously admissible futures has become inadmissible---and restore
it through targeted repair operations.

The repair operation $R(x, t)$ restores admissibility by modifying
the system's state or its environment in ways that reopen closed
futures. The operation is not undoing what happened; in an
irreversible system, that is generally impossible. It is finding
a new path to the future from the current position---a path that
maintains admissibility even though the path that was damaged is
no longer available.

Repair Theory identifies the preservation target as the admissibility
relation itself: the structure of which transitions are possible from
which states. This is the generator of all future navigation.
Preserve the admissibility relation and futures can be reconstructed,
even from a damaged or reduced state inventory. Lose the admissibility
relation and the system cannot navigate even if the full inventory
of past states is intact. The relation is the generator; the
inventory of navigated states is derived.

\section{The Eight-Letter Keyboard and Motor Compression}

The Eight-Letter Keyboard proposes a symbolic system generated from
a small number of motor primitives: finger identities crossed with
directional distinctions. The full alphabet---and in principle any
symbolic inventory of comparable size---can be reconstructed from
this generative substrate. The keyboard is not a convenient encoding
of an alphabet. It is a demonstration that the alphabet is a derived
object, generated from something considerably smaller.

The same principle extends to skilled typists and, more strikingly,
to gesture-based text input. A skilled practitioner of gesture
typing does not remember letters; they remember paths. The hand
traces a trajectory through the keyboard manifold, and the letter
sequence is the result of parsing that trajectory against the
lexical graph. The user's knowledge is not an inventory of letter
positions but a repertoire of path patterns: the generator. The
letters are reconstructed from the paths, not stored alongside them.

This is the $G \to R(G)$ schema in its most concrete and accessible
form. The generator is the motor program: a trajectory through a
constrained manifold. The inventory is the symbol sequence: the
equivalence class of trajectories that produce that sequence under
the lexical admissibility relation. Two different hand trajectories
that produce the same word are equivalent under the relation. The
word is the equivalence class; the trajectory is the primitive.

The generalization is immediate. A Swype user who loses access to
the specific letters they know how to reach is not helpless; they
still have the path patterns, and from the path patterns the letters
can be reconstructed. A user who retains the letter inventory but
loses the path patterns is in a more serious situation, because the
path patterns are the generator. Preserving the motor program
preserves more than preserving the symbol set it produces.

\section{The Pattern Before the Abstraction}

The six frameworks present the same structure from six different
directions. RSVP identifies the field triple as the generator of
future geometry. CLIO identifies the projection as the generator
of epistemically adequate representations. Spherepop identifies
the operation vocabulary as the generator of cognitive continuation.
MEM\textbar{}8 identifies the residue field as the generator of
identity over time. Repair Theory identifies the admissibility
relation as the generator of navigable futures. The Eight-Letter
Keyboard identifies the motor program as the generator of the
symbol inventory.

In every case the same asymmetry holds: the generator is small,
the inventory is large, and the generator suffices to reconstruct
the inventory. Preserving the generator is therefore strictly
more efficient than preserving the inventory, and---crucially---
strictly more important: the inventory can be lost and recovered
from the generator, but the generator cannot be recovered from
a damaged inventory.

\begin{proposition}[Generative Reconstruction]
\label{prop:gen_recon}
In each source domain there exists a surjective reconstruction map
$\rho : R(G) \to I$ from the reachable space generated by $G$ onto
the inventory $I$, with $|G| \ll |I|$.
\end{proposition}

\begin{proof}
By explicit construction for each domain. In RSVP: $(\Phi, v, S)$
generates the full space of admissible field trajectories via the
field equations; the inventory of accessible states is the image.
In CLIO: the projection structure generates the equivalence class
decomposition; the inventory of projected representations is the
image. In Spherepop: the operation vocabulary generates the space
of admissible cognitive continuations; the inventory of cognitive
states visited is the image. In MEM\textbar{}8: the residue field
generates the space of retrievable memory trajectories; the inventory
of stored memories is the image. In Repair Theory: the admissibility
relation generates the space of possible repairs; the inventory of
repaired states is the image. In the Eight-Letter Keyboard: the
motor program generates the space of producible symbol sequences;
the inventory of symbols is the image. In each case surjectivity
follows from the construction of $G$ as the generating structure
for $I$, and the size inequality $|G| \ll |I|$ follows from the
compression ratio of the generating map.
\end{proof}

The pattern is established. Chapter~5 asks why it matters---why
naïve preservation instincts target the inventory rather than the
generator, and what the cost of that mistake is.

%% --------------------------------------------------------------
\chapter{The Asymmetry of Preservation}
%% --------------------------------------------------------------

The six frameworks surveyed in Chapter~4 share a common structure:
a small generator $G$ produces a large inventory $I$ via a
reconstruction map, and $|G| \ll |I|$ while $R(G) \approx R(I)$.

A seventh case is worth noting before the asymmetry is made precise.
The PERSCEN recommendation framework \citep{Du2025} constructs a
user-specific feature graph for each agent, where the adjacency
matrix $A_u^{(1)}$ encodes which feature interactions are active for
that user. Different users inhabit different local geometries: the
graph is not a fixed shared space but a personalized continuation
structure. Viewed through the admissibility lens, $A_u \approx
\mathcal{A}(u)$---the adjacency matrix is a discrete approximation to
the local admissibility manifold of user $u$. The codebook learned by
vector quantization across scenarios is a generative minimum:
preference structures shared across contexts that can reconstruct
scenario-specific behavior through projection. The scenario-aware
GLU gate ($g_u^{(l)} = (\text{shared}) \otimes \sigma(\text{scenario})$)
is operationally an admissibility filter: it gates which shared
representations survive under a particular scenario context. PERSCEN
does not formulate these operations in terms of reachability, and its
ontology remains representational rather than continuation-first. But
at the level of mechanism, it is another instance of the pattern:
generators (the graph structure, the codebook) are smaller and more
robust than the inventories they produce (the scenario-specific
embeddings), and the gating operation selects admissible continuations
from the generated space.

This asymmetry has a practical consequence that is easily stated
but non-obvious: the correct preservation target is the generator,
not the inventory. The instinct to preserve inventories---to maintain
the full set of stored objects, recorded facts, or accumulated
states---is not wrong in general, but it is wrong about where
value resides in generative systems.

This chapter makes the asymmetry precise, identifies why the naïve
instinct fails, and extends the analysis from symbolic systems to
the full range of domains in which the pattern appears.

\section{Inventories Versus Generators}

Consider two strategies for maintaining a system across a
perturbation. The inventory strategy preserves as many elements
of $I$ as possible: it backs up states, stores representations,
maintains databases of past content. The generator strategy
preserves $G$: it maintains the structure from which $I$ can
be reconstructed. For a non-generative system, where no small
generator exists, the inventory strategy is the only option. For
a generative system, the two strategies are not equivalent.

The generator strategy is superior along two independent dimensions.
First, it is more efficient: $|G| \ll |I|$, so preserving the
generator requires far less storage, transmission bandwidth, and
maintenance effort than preserving the inventory. Second, and more
importantly, it is more robust: the inventory can be lost and
recovered from the generator, but a damaged inventory cannot
reconstruct a lost generator.

The second point requires emphasis because it is counterintuitive.
Suppose ninety percent of the inventory $I$ is intact and the
generator $G$ is lost. Can the generator be recovered from the
surviving inventory? In general, no. The map $\rho : R(G) \to I$
is surjective but not injective: many elements of $R(G)$ may map
to the same element of $I$, and the information needed to invert
the map is exactly the information encoded in $G$. A partial
inventory without $G$ is like a large collection of words without
knowledge of the phonological or morphological rules that generated
them: useful as far as it goes, but incapable of generating what
is missing.

Conversely, suppose ninety percent of the inventory is lost and
the generator $G$ is intact. The lost inventory can be recovered
by applying $\rho$ to $R(G)$. The system can regenerate what was
lost. The asymmetry is stark: losing the generator is catastrophic
and irreversible; losing the inventory is costly but recoverable.

\begin{definition}[Inventory Entropy and Generator Entropy]
For a generative system with inventory $I$ and generator $G$:
\[
S_I = \log|I|, \qquad S_G = \log|G|.
\]
The compression ratio of the system is $S_I / S_G \gg 1$.
\end{definition}

\begin{theorem}[Asymmetry of Preservation]
\label{thm:asymmetry}
For a generative system satisfying Proposition~\ref{prop:gen_recon},
the generator $G$ is a strictly more robust preservation target
than the inventory $I$: loss of $I$ given $G$ is recoverable;
loss of $G$ given $I$ is not recoverable in general.
\end{theorem}

\begin{proof}
Recovery from inventory loss: given $G$ intact and $I$ partially
or wholly lost, apply $\rho \circ R$ to recover $I$ from $G$.
Recoverability follows from the surjectivity of $\rho$ established
in Proposition~\ref{prop:gen_recon}.

Non-recovery from generator loss: given $I$ intact and $G$ lost,
recovery of $G$ requires inverting $\rho \circ R$. This inverse
exists only if $\rho \circ R$ is injective, which requires the
generator $G$ to be uniquely determined by the inventory $I$. But
if $G$ were uniquely determined by $I$, the system would not be
generative in the relevant sense: the generator would be redundant,
and the system would reduce to its inventory. In any properly
generative system the map $\rho \circ R$ is not injective, and
recovery of $G$ from $I$ alone is not possible in general.
\end{proof}

\section{The Eight-Letter Keyboard Argument}

The Eight-Letter Keyboard provides the clearest possible illustration
of Theorem~\ref{thm:asymmetry} because both $G$ and $I$ can be
described concretely and the asymmetry can be verified directly.

The generator $G$ is the set of motor primitives:
$G = \mathrm{Finger} \times \mathrm{Direction}$, with $|\mathrm{Finger}|
= 8$ and $|\mathrm{Direction}|$ determined by the directional
vocabulary of the system (typically four to eight directions per
finger). The inventory $I$ is the alphabet, with $|I| = 26$ for
the Latin alphabet or larger for scripts with more characters.
The reconstruction map $\rho : R(G) \to I$ assigns each
finger-direction combination to the letter it produces.

The asymmetry is immediate. If the full alphabet $I$ is intact
but the motor program $G$ is lost---if a user knows what letters
exist but has forgotten the gestures that produce them---the
motor program cannot be recovered from the alphabet alone. Many
different assignments of gestures to letters are consistent with
knowing the inventory. The specific assignment encoded in $G$ is
additional information.

If the motor program $G$ is intact but several letters of the
alphabet $I$ are forgotten---if a user knows the gestures but
cannot immediately recall which letter each one produces---the
missing letters can be recovered by executing the gestures and
reading off the results. The generator produces the inventory
on demand.

The deeper point is that a skilled user does not primarily
possess the inventory. They possess the generator. A skilled
typist or gesture-input user does not consciously recall letter
positions; they execute motor programs. When they produce a
letter, they are not looking up a stored location but
following a learned trajectory. The letter is the output
of the path, not the object of the memory. The path is
the primitive; the letter is derived.

\section{Script Projection: Motor Programs Surviving Domain Transfer}

Script projection extends the Eight-Letter Keyboard argument
to demonstrate that generators survive domain transfer in a
way that inventories do not.

Consider a user who has learned a gesture-input vocabulary for
one script---say, the Latin alphabet. The motor programs
constitute a generator $G_L$ for the inventory $I_L$ of Latin
characters. Now suppose the user needs to input a different script
---say, Arabic or Greek. The inventory changes entirely:
$I_L \cap I_A = \emptyset$ for Latin and Arabic inventories.
Under the inventory-preservation strategy, the user must start
over: the Latin inventory is useless for Arabic input.

Under the generator-preservation strategy, the situation is
different. The motor program $G_L$ encodes not specific
letter-gesture associations but a motor manifold: a space of
executable trajectories. The script projection assigns a new
inventory $I_A$ to the same generator $G$, defining a new
reconstruction map $\rho_A : R(G) \to I_A$ that maps the
same trajectories to Arabic characters. The user's motor
programs survive the transition; only the assignment changes.

This is not merely a convenience. It reveals something
structural about generators: they are domain-transcendent
in a way that inventories are not. The motor manifold $G$
is not defined relative to any particular inventory. It is
a structure in its own right---a space of executable
trajectories with its own topology and admissibility
structure---and different inventories are different
projections of that same structure onto different symbolic
domains.

The generalization: a generator that supports script projection
preserves more future potential than the inventory of any
particular script, because it remains productive under
domain changes that would render any specific inventory
useless. The generator outlives its current inventory.

\section{Why Naïve Preservation Targets the Wrong Layer}

Given the asymmetry demonstrated above, why does the
naïve preservation instinct consistently target the
inventory rather than the generator? Three factors
contribute.

First, the inventory is visible and the generator is not.
The letters of an alphabet can be written down; the motor
programs that generate them cannot be easily externalized.
The states that have been visited can be recorded; the
admissibility relation that governs which states can be
reached from which others requires a more abstract description.
Preservation instincts target what can be directly observed
and copied, which is the inventory.

Second, the inventory feels complete and the generator feels
partial. An inventory of twenty-six letters covers the
alphabet; a set of eight finger positions does not look like
it covers twenty-six things. The generator's completeness
is conditional on the reconstruction map, which is an
additional structure that must itself be preserved. This
makes the generator feel less self-sufficient, and
therefore a less reliable preservation target---even though
the inverse is true.

Third, the cost of inventory loss is immediately salient
and the cost of generator loss is deferred. When inventory
items are lost, their absence is immediately noticeable:
a letter is missing, a memory cannot be retrieved, a state
cannot be reached. When the generator is degraded, the
consequences are less immediately visible: the inventory
that exists is still intact, and only the capacity to
generate new inventory is diminished. The deferred cost
of generator loss makes it systematically underweighted
by systems optimizing for immediate performance.

All three factors push naïve preservation toward the
inventory. All three are, in generative systems, systematically
wrong. The correct preservation target is the generator,
even though it is less visible, feels less complete, and
whose loss has deferred costs.

\section{The General Principle}

The asymmetry of preservation, established for the
Eight-Letter Keyboard and extended to script projection,
generalizes across all the domains surveyed in Chapter~4.

In RSVP: the field triple $(\Phi, v, S)$ is the generator;
the inventory of accessible states is what it generates.
Preserving the field triple under perturbation is strictly
more robust than preserving the state inventory, because
the state inventory can be regenerated from the field but
the field cannot in general be recovered from the state
inventory alone.

In MEM\textbar{}8: the residue field deformation is the
generator; the inventory of stored memories is what it
generates. A memory system that preserves residue field
structure under perturbation can regenerate memories that
were lost; a memory system that preserves stored objects
but loses residue field structure cannot navigate to the
memories it still holds.

In Repair Theory: the admissibility relation is the
generator; the inventory of reachable states is what it
generates. A system that preserves its admissibility
relation under damage can find new paths to futures whose
original paths have been destroyed; a system that preserves
its state history but loses its admissibility relation
cannot navigate forward regardless of how much it remembers.

The pattern is uniform. The conclusion is the same in
every case: value resides in the generator, not the
inventory. The next chapter formalizes this conclusion
as the Generative Compression Hypothesis and its
operational principle.

%% --------------------------------------------------------------
\chapter{The Generative Compression Hypothesis}
%% --------------------------------------------------------------

The evidence assembled in Chapters~4 and~5 establishes a
recurring pattern: in generative systems, the correct
preservation target is the generator rather than the
inventory, and preserving the generator is both more
efficient and more robust. This chapter converts the
pattern into a formal hypothesis and derives its
operational principle.

\section{Formal Statement}

The Generative Compression Hypothesis is a claim about what
intelligent persistence consists in. It is not a claim about
any particular system or domain; it is a claim about the
general structure shared by systems that maintain productive
capacity over time.

\begin{hypothesis}[Generative Compression]
\label{hyp:gen_compression}
A system exhibits intelligent persistence when it preserves
the minimal generative structure necessary to reconstruct
a rich space of admissible future continuations. The measure
of a system's intelligent persistence is accordingly the
richness of the admissible future reachable from its current
generative state, not its predictive accuracy, its utility
score, or the breadth of its stored representations.
\end{hypothesis}

Three aspects of this hypothesis require emphasis. First,
it identifies the correct preservation target as \emph{minimal}
generative structure. Not all generative structure, but the
smallest structure sufficient to reconstruct a rich space
of admissible continuations. This is the compression claim:
the system should not preserve more than is necessary.

Second, it identifies the measure of persistence as the
\emph{richness of the admissible future}, not the accuracy
of current representation. A system that maintains a highly
accurate representation of the current state but has narrow
future potential is less persistent, by this measure, than
a system with a less accurate current representation but
rich future potential. The hypothesis inverts the standard
priority.

Third, it makes intelligent persistence a relational concept:
a system is persistent \emph{relative to} its capacity to
generate admissible continuations. The same physical substrate
may be highly persistent in one context and non-persistent
in another, depending on whether its generative capacity
is being maintained.

\section{The Minimal Generator}

The hypothesis references the minimal generative structure
sufficient to reconstruct a rich space of admissible
continuations. This object can be defined precisely.

\begin{definition}[Minimal Generator]
For a system with inventory $I$ and reconstruction map
$\rho : R(G) \to I$, the minimal generator is:
\[
G^* = \arg\min_{G} |G| \quad \text{subject to} \quad
R(G^*) \supseteq R(I),
\]
where $R(I)$ denotes the set of admissible continuations
reachable from the current inventory.
\end{definition}

The constraint $R(G^*) \supseteq R(I)$ ensures that the
minimal generator preserves at least as rich a space of
admissible futures as the full inventory. The minimization
ensures that no unnecessary structure is retained. $G^*$
is the tightest possible generator: the most compressed
structure from which the full future potential can be
reconstructed.

In practice, exact minimization over $G$ is rarely
tractable. The definition is normative: it specifies what
an ideally persistent system would preserve. Approximate
minimization---preserving a generator that is small relative
to the inventory while maintaining most of the relevant
future potential---is the operational target.

\section{The Generative Sufficiency Principle}

The hypothesis and the minimal generator definition together
imply a principle that is the operational center of this part
of the monograph. It is the claim that the hypothesis makes
actionable: given that $G^*$ exists and a reconstruction map
is available, $G^*$ is not merely a more efficient alternative
to $I$ but the correct preservation target.

\begin{principle}[Generative Sufficiency Principle]
\label{prin:sufficiency}
If a reconstruction map $\rho : R(G^*) \to I$ exists, then
$G^*$ is the correct preservation target: preserving $G^*$
is necessary and sufficient for future recovery of the full
admissible continuation structure.
\end{principle}

Sufficiency follows directly from the constraint in the
definition of $G^*$: if $R(G^*) \supseteq R(I)$, then
the full future potential is recoverable from $G^*$.
Necessity follows from Theorem~\ref{thm:asymmetry}: in a
genuinely generative system, $I$ cannot recover $G^*$, so
any preservation strategy that discards $G^*$ in favor of
$I$ is not sufficient for future recovery.

The principle is the formal counterpart of the Eight-Letter
Keyboard argument. The keyboard argument showed, for a
specific domain, that the motor program is the correct
preservation target. The Generative Sufficiency Principle
states that this is true whenever a reconstruction map
exists---which is whenever the system is genuinely generative.

\section{Verification Across Source Domains}

The Generative Sufficiency Principle can be verified in each
of the six source domains surveyed in Chapter~4.

In RSVP cosmology and cognitive field theory: the field
triple $(\Phi, v, S)$ is $G^*$. The reconstruction map
is the RSVP field evolution: given $(\Phi, v, S)$ at time
$t$, the field equations generate the full space of
admissible future field configurations. Preserving $(\Phi,
v, S)$ is necessary and sufficient for future reconstruction
of the accessible state space.

In CLIO projection theory: the projection structure
$\pi : X \to M$ together with its admissibility-preservation
condition is $G^*$. The reconstruction map is the minimal
admissibility projection established by the Minimal
Projection Theorem (Chapter~10). Preserving the projection
is necessary and sufficient for future reconstruction of
the epistemically adequate representation space.

In Spherepop cognitive dynamics: the operation vocabulary
(pop, refuse, collapse, bind, repair) is $G^*$. The
reconstruction map is the operational dynamics: given the
vocabulary, the full space of admissible cognitive
continuations can be generated. Preserving the vocabulary
is necessary and sufficient for future reconstruction of
the cognitive continuation space.

In MEM\textbar{}8 memory dynamics: the residue field
deformation structure is $G^*$. The reconstruction map is
the resonance retrieval process: given the residue field,
the full space of retrievable memory trajectories can be
generated by perturbation and attractor navigation.
Preserving the residue field is necessary and sufficient
for future reconstruction of the memory trajectory space.

In Repair Theory: the admissibility relation $\mathcal{A}$
is $G^*$. The reconstruction map is the repair operation:
given $\mathcal{A}$, the full space of admissible
continuations from any current state can be determined.
Preserving $\mathcal{A}$ is necessary and sufficient for
future reconstruction of the navigable future space.

In the Eight-Letter Keyboard: the motor program is $G^*$.
The reconstruction map is the gesture-to-symbol assignment.
Preserving the motor program is necessary and sufficient
for future reconstruction of the producible symbol inventory.

The hypothesis and principle are therefore not merely
consistent with the six source frameworks. They unify them:
each framework is a domain-specific instance of the same
abstract claim.

\section{The Hypothesis as Organizing Principle}

The Generative Compression Hypothesis is the organizing
principle of the Admissibility Program. It does not follow
from the Continuation Principle of Chapter~3 but is
complementary to it: the Continuation Principle tells us
what meaning is (admissible future structure), while the
Generative Compression Hypothesis tells us what must be
preserved to maintain it (the minimal generator of that
structure).

Together they define the research program. Every project
within the program can be understood as addressing one of
two questions: what is the admissible continuation structure
of this domain? And what is the minimal generator of that
structure? The first question is addressed by the Continuation
Principle and its mathematical formalization in Part~III.
The second is addressed by the Generative Compression
Hypothesis and the Generative Sufficiency Principle.

%% --------------------------------------------------------------
\chapter{Generators, Reachability, and Preservation}
%% --------------------------------------------------------------

The preceding chapters have established the Generative
Compression Hypothesis and the Generative Sufficiency
Principle. This chapter does something different: it makes
the abstract schema $G \to R(G)$ visible across every
major framework of the Admissibility Program simultaneously,
and asks what the recurrence of this schema implies about
the nature of the program itself.

The chapter is the conceptual center of the monograph. It
is not a technical chapter; it contains one theorem and
one definition, both brief. Its function is to let the
reader see that what appeared to be seven different
investigations---cosmological field theory, epistemological
projection, cognitive dynamics, memory theory, repair
epistemology, motor compression, computation architecture,
and political economy---are instances of a single question:
what generators are worth preserving?

\section{The Abstract Schema}

Let $G$ be a generative structure and $R(G)$ the reachable
distinction space it produces. The abstract schema is:
\[
G \;\longrightarrow\; R(G).
\]
A system governed by this schema is characterized by three
properties. Its inventory $I$ is a quotient of $R(G)$: every
element of $I$ is an equivalence class of elements of $R(G)$
under the admissibility relation. Its generative capacity is
encoded in $G$, not in $I$. And its persistence under
transformation is determined by whether $G$ survives the
transformation, not whether $I$ does.

The schema is abstract. Its power comes from the variety
of its instantiations.

\section{Instantiations Across Frameworks}

\begin{center}
\renewcommand{\arraystretch}{1.5}
\begin{tabular}{lll}
\toprule
Framework & Generator $G$ & Reachable Space $R(G)$ \\
\midrule
Alphabet &
  $\mathrm{Finger} \times \mathrm{Direction}$ &
  Symbol inventory \\
Script Projection &
  Motor program &
  Producible script space \\
MEM\textbar{}8 &
  Residue field deformation &
  Retrievable identity \\
Repair Theory &
  Admissibility relation $\mathcal{A}$ &
  Navigable future space \\
RSVP &
  $(\Phi, v, S)$ &
  Future field geometry \\
HYDRA &
  $(\Phi,v,S) + T$ &
  Admissible future trajectories \\
Political Economy &
  Institutional structure &
  Reachable life trajectories \\
\bottomrule
\end{tabular}
\end{center}

Each row instantiates the same three properties. The
generator is smaller than the inventory it produces. The
reachable space is generated by the generator, not stored.
And persistence under transformation depends on preservation
of the generator.

The political economy row deserves comment here, since
it will be developed fully only in Chapter~18. The
institutional structure of a society---the rules, norms,
resource distributions, and opportunity maps that govern
who can do what---is the generator of the space of life
trajectories available to its members. Two societies with
identical demographic inventories (the same people with
the same characteristics) but different institutional
structures have different reachable trajectory spaces.
The institutional structure is the generator; the
distribution of outcomes is the derived inventory. What
must be preserved to maintain a rich space of future
social trajectories is the institutional structure, not
any particular distribution of current outcomes.

\section{Reachability as a Functor}

The abstract schema $G \to R(G)$ can be given categorical
precision. The assignment of a reachable space to each
generator is not merely a set-theoretic map; it respects
the morphism structure of the category of generators.

\begin{definition}[Category of Generators]
$\mathbf{G}$ has generators as objects and admissible
reconstruction maps as morphisms. A morphism
$f : G_1 \to G_2$ in $\mathbf{G}$ is a map that sends
$G_1$-generated distinctions to $G_2$-generated
distinctions in an admissibility-preserving way.
\end{definition}

\begin{theorem}[Functoriality of Reachability]
\label{thm:R_functor}
The reachability assignment $R : \mathbf{G} \to
\mathbf{Reach}$ is a functor, where $\mathbf{Reach}$
is the category of reachable distinction spaces with
admissibility-preserving maps as morphisms.
\end{theorem}

\begin{proof}
$R$ sends each generator $G$ to its reachable distinction
space $R(G)$ and each admissibility-preserving map
$f : G_1 \to G_2$ to the induced map
$R(f) : R(G_1) \to R(G_2)$ defined by
$R(f)(\gamma) = f \circ \gamma$ for trajectories
$\gamma$ in $R(G_1)$. Identity maps are preserved:
$R(\mathrm{id}_G) = \mathrm{id}_{R(G)}$. Composition
is preserved: $R(g \circ f) = R(g) \circ R(f)$ by the
associativity of function composition. Admissibility
preservation: if $f$ is admissibility-preserving, then
$R(f)$ sends admissible trajectories in $R(G_1)$ to
admissible trajectories in $R(G_2)$ by the definition
of admissibility-preserving maps.
\end{proof}

The functoriality of $R$ is the precise statement that
reachability is a structural property, not an incidental
one: it is preserved under the morphisms of the category
of generators. When a system changes its generator---
through learning, repair, compression, or domain
transfer---the reachable space changes in a way that
tracks the change in the generator. The generator is
not merely a convenient proxy for the reachable space;
it is the structure that determines it, up to the
natural transformations that $R$ preserves.

\section{The Unifying Question}

The functor $R : \mathbf{G} \to \mathbf{Reach}$ makes
precise the sense in which the seven instantiations in
the table are instances of a single abstract structure.
Each framework identifies a generator, a reachable space,
and a reconstruction map, and asks how to preserve the
generator under the transformations relevant to its
domain. The question is the same; the domains differ.

This unification suggests that the correct level of
analysis for the Admissibility Program is not the
domain but the abstract structure. Improvements in
understanding reachability in one domain transfer to
other domains by the functoriality of $R$. Theorems
proved about generators in one framework apply to
generators in all frameworks that instantiate the
same categorical structure. The Minimal Projection
Theorem, proved in Chapter~10 for trajectory spaces
in general, applies to each of the seven rows in the
table.

The unifying question is therefore not about any
particular domain. It is: \emph{what generators are
worth preserving?} The answer, implied by the
Generative Sufficiency Principle, is: the minimal
generators sufficient to reconstruct a rich admissible
future. The question and the answer have the same form
in every domain. The mathematical infrastructure of
Part~III provides the tools to pursue them with precision.

\bigskip
\noindent\emph{What follows.} Part~III develops the
mathematical core. Chapter~8 introduces the four
primitive objects from which all subsequent machinery
is built. Chapter~9 establishes the category of
admissible states and proves the functoriality of the
HYDRA architecture. Chapter~10 proves the Minimal
Projection Theorem, the mathematical formalization of
the Continuation Principle. Chapters~11 through~14
develop the RSVP field theory, stratified semantic
manifolds, sheaf-theoretic semantics, and memory
dynamics that give the abstract framework its
field-theoretic content.




%% ==============================================================
\part{The Mathematics of Reachability}
%% ==============================================================

\noindent\textit{Parts~I and~II established the philosophical
argument: prediction-primary architectures fail because they
preserve the wrong objects, and what must be preserved instead
is the minimal generator of admissible future continuations.
Part~III formalizes that argument. It introduces the four
primitive objects from which all subsequent machinery is
constructed, establishes the categorical structure of the
HYDRA architecture, and proves the Minimal Projection
Theorem---the mathematical heart of the Continuation Principle.
It then develops the field theory, stratification, sheaf
semantics, and memory dynamics that give the abstract framework
its quantitative content. Each chapter contributes one formal
object, one central result, and a bridge to the next chapter.
The statement classification introduced in the Preface is
applied throughout: established mathematics is identified as
such, framework-specific definitions are distinguished from
derived results, and open problems are named rather than
papered over.}

\bigskip

%% --------------------------------------------------------------
\chapter{Four Primitive Objects}
%% --------------------------------------------------------------

Every major result in this monograph is a specialization of
four mathematical objects. Introducing them here, before the
specific field equations and architectural details, gives the
subsequent chapters a common reference point. The claim that
every chapter is a specialization of these four objects is
itself a conjecture---the Structural Universality Conjecture
of Chapter~19---but the specialized cases are established
individually, and the common structure is visible throughout.

\section{Trajectory Space}

\begin{definition}[Trajectory Space]
A trajectory space $\mathcal{X}$ is a topological space of
agent-environment histories, equipped with an admissibility
structure $\mathcal{A} \subset \mathcal{X}$ and a measure
$\mu$. Elements of $\mathcal{X}$ are trajectories
$\gamma : [0,1] \to X$ in an underlying state space $X$;
the measure $\mu$ assigns weight to regions of trajectory
space relevant for admissibility computations.
\end{definition}

The trajectory space is the primary ontological object of the
framework. States are not elements of $\mathcal{X}$; states
are positions along trajectories. A state $x \in X$ is
characterized not by its intrinsic properties but by the
trajectories that pass through it and, in particular, by
the admissible continuations available from it. This is the
operational content of the Continuation Principle: the
relevant structure is the trajectory space and its
admissibility, not the state space and its properties.

The measure $\mu$ allows quantitative comparisons between
regions of trajectory space: one region is larger, more
accessible, or more heavily weighted than another. In the
RSVP field-theoretic setting, $\mu$ is derived from the
entropic accessibility measure $S$; in the discrete setting
of restricted mind-change learning (Chapter~15), $\mu$
assigns uniform weight to admissible hypothesis transitions.

\section{Projection System}

\begin{definition}[Projection System]
A projection system over $\mathcal{X}$ is a family
$\Pi = \{\pi_i : \mathcal{X} \to M_i\}_{i \in I}$ of
continuous surjections onto lower-dimensional operational
manifolds $M_i$, satisfying mutual consistency on overlapping
domains: if $U \subset M_i \cap M_j$, the restrictions of
$\pi_i$ and $\pi_j$ to $\pi_i^{-1}(U) \cap \pi_j^{-1}(U)$
agree up to the natural identification of $U$ in both
manifolds.
\end{definition}

A projection system is how a cognitive, physical, or social
system reduces the full trajectory space to a tractable
representation. Each projection $\pi_i$ maps from the
unmanageable full space $\mathcal{X}$ to a manifold $M_i$
that can be processed, stored, or communicated. The mutual
consistency condition ensures that different projections of
the same trajectory cohere: there is a well-defined sense
in which different representations of the same trajectory
agree on their overlap.

The HYDRA architecture is a projection system in this sense.
Each of its six components produces a projection of the
trajectory space onto a lower-dimensional space---the
activated subspace $X_R$, the reachability graph
$\mathcal{G}$, the admissible graph $\mathcal{G}_a$, the
trajectory space $\Gamma$, the memory-weighted trajectory
space $\Gamma_M$, and the output space $Y$---and the
composition is a chain of such projections satisfying mutual
consistency by the functoriality of each component.

\section{Admissibility Structure}

\begin{definition}[Admissibility Structure]
\label{def:admissibility}
An admissibility structure on $\mathcal{X}$ is a map
$\mathcal{A} : \mathcal{X} \to 2^{\mathcal{X}}$ assigning
to each trajectory $\gamma$ its set of admissible
continuations $\mathcal{A}(\gamma) \subset \mathcal{X}$.
A continuation $\gamma'$ is admissible relative to $\gamma$
if $\gamma' \in \mathcal{A}(\gamma)$. Admissibility is
relational: a trajectory is not admissible in isolation but
relative to a current trajectory and a context.
\end{definition}

The admissibility structure is the most important of the
four primitive objects because it is the one whose
preservation is the framework's central concern. What the
Continuation Principle says is that two trajectories are
semantically equivalent when they have the same admissibility
structure. What the Generative Compression Hypothesis says
is that the minimal generator of the admissibility structure
is the correct preservation target. What the Minimal
Projection Theorem says is that the minimal admissibility-
preserving projection is well-defined and unique.

The relational character of admissibility---that it is
defined relative to a trajectory and context rather than
intrinsically---distinguishes this framework from standard
constraint-satisfaction approaches, which treat admissibility
as a property of states. Trajectory-relative admissibility
is richer: the same state may be admissible from one
trajectory and inadmissible from another, because what
is admissible from a state depends on the history of
approaches to that state.

\section{Persistence Functional}

\begin{definition}[Persistence Functional]
A persistence functional $P(\mathcal{X}, \Pi, \mathcal{A})$
is a Lyapunov-type functional measuring the degree to which
the projection system $\Pi$ preserves the admissibility
structure $\mathcal{A}$ under the field dynamics on
$\mathcal{X}$. Formally, $P \geq 0$ with $P = 0$ when
$\Pi$ perfectly preserves $\mathcal{A}$, and $P$ is
non-increasing along admissible trajectories:
$\frac{d}{dt}P(\mathcal{X}(t), \Pi, \mathcal{A}) \leq 0$
whenever the system evolves along an admissibility-preserving
path.
\end{definition}

The persistence functional is the quantitative measure of
how well the system is doing its job: preserving admissible
future structure under projection. A system whose persistence
functional is small and non-increasing is successfully
maintaining its continuation capacity. A system whose
persistence functional is large or increasing is losing
admissibility under its projections---collapsing distinctions
it needed to maintain, or failing to preserve the field
geometry that governs future navigation.

In the RSVP setting, the persistence functional takes the
concrete form $P = \int_M (|\Phi|^2 + |v|^2 + S)\,
\mathrm{dvol}_g$, and its non-increase is the content of
the Lyapunov stability result proved in Chapter~11. In
the engineering realization of Chapter~20, the persistence
functional is approximated by the admissibility volume
$\mathrm{Vol}(\mathcal{A}(v)) > \varepsilon$, and its
maintenance is the criterion against which Marine, MEM\textbar{}8,
Phoenix Protocol, and AyeOS are each evaluated.

\section{Lamphron and Lamphrodyne}

The admissibility structure $\mathcal{A}$ assigns to each trajectory
its set of admissible continuations. The \emph{volume} of that set
is a scalar quantity that will appear, under different names and in
different domains, throughout every subsequent chapter. Naming it
once here allows the later chapters to recognize the same object
when they encounter it.

\begin{definition}[Lamphron Field]
\label{def:lamphron}
Let $(\mathcal{X}, \mathcal{A})$ be an admissibility space equipped
with a measure $\mu$ on continuation sets. The lamphron field is
the scalar field $L : \mathcal{X} \to \mathbb{R}$ defined by:
\[
L(x) = \log \mathrm{Vol}(\mathcal{A}(x)),
\]
where $\mathrm{Vol}(\mathcal{A}(x)) = \mu(\mathcal{A}(x))$ is the
measure of the admissible continuation set at $x$. A state with
high lamphron has many viable continuations; a state with low
lamphron is brittle.
\end{definition}

\begin{definition}[Lamphrodyne Flow]
\label{def:lamphrodyne}
The lamphrodyne flow on $\mathcal{X}$ is the vector field
$\mathcal{L} = \nabla L$ generating the dynamical system
$\dot{x} = \mathcal{L}(x)$. Lamphrodyne motion carries the
system toward states of greater future admissibility.
\end{definition}

\begin{principle}[Lamphrodyne Descent]
\label{prin:lamphrodyne}
For any admissibility-preserving evolution,
$\frac{dL}{dt} \geq 0$: lamphron is non-decreasing along
admissible trajectories. Systems that violate this condition
are destroying future structure faster than they are
preserving it.
\end{principle}

The lamphron field unifies several quantities that appear under
different names across the framework. In RSVP (Chapter~11), $L$
is instantiated as the entropic accessibility measure
$S(x,t) = \log|\mathcal{A}(x,t)|$, which already appears in the
field triple. In HYDRA (Chapter~3), $L$ is instantiated as
$\log\,\mathrm{Vol}(\mathcal{A}(v))$, the quantity that the
admissibility filtering stage $F_a$ maintains above $\varepsilon$.
In MEM\textbar{}8 (Chapter~14), $L$ is the continuation capacity
of a memory residue. In political economy (Chapter~18), $L$ is the
log-volume of life trajectories reachable from a given institutional
position. In every case the same question is being asked from a
different domain: how much future structure remains available from
here? The lamphron field is the unified name for the answer.

\section{Everything as Specialization}

The claim of this chapter---that every major result in the
monograph is a specialization of $(\mathcal{X}, \Pi,
\mathcal{A}, P)$---can be checked against each subsequent
chapter.

In Chapter~9 (Category Theory): the category $\mathbf{AdmSt}$
is the category whose objects are elements of
$(\mathcal{X}, \mathcal{A})$ and whose morphisms are
admissibility-preserving trajectories. The HYDRA functors
are morphisms in this category; the persistence functional
measures how well the composition preserves $\mathcal{A}$.

In Chapter~10 (Minimal Projection Theorem): the minimal
projection $\pi^* : \mathcal{X} \to M^*$ is the minimal
element of the projection system $\Pi$ that preserves
$\mathcal{A}$ exactly. The theorem establishes its existence
and uniqueness.

In Chapter~11 (RSVP): the trajectory space is the space
of field configurations $(\Phi, v, S)$ over a smooth
manifold $M$; the admissibility structure is given by the
field equations; the persistence functional is the
Lyapunov functional $P = \int (|\Phi|^2 + |v|^2 + S)\,
\mathrm{dvol}_g$.

In Chapter~12 (Stratified Manifolds): the trajectory space
is stratified by the Whitney decomposition; the admissibility
structure governs which stratum transitions are permitted;
the persistence functional controls tile annotation entropy.

In Chapter~13 (Sheaf Semantics): the projection system is
the sheaf $\mathcal{F}$ over the context space; the
admissibility structure is the gluing condition for local
sections; the persistence functional is the cohomological
obstruction $[c] \in \check{H}^1(\mathcal{U}, \mathcal{F})$.

In Chapter~14 (Memory): the trajectory space is
$\Gamma_M$ (trajectories weighted by residue); the
admissibility structure governs which memory trajectories
are retrievable; the persistence functional is the
memory decay rate relative to the critical threshold
$\lambda = \kappa c_S$.

In Chapter~7, the Eight-Letter Keyboard instantiation is
$G = \mathrm{Finger} \times \mathrm{Direction}$ with
$R(G) = \mathrm{Symbol\ inventory}$. The information-theoretic
measure of what is lost by collapsing the generator to its
projected inventory is the symbolic collapse entropy:
$\Delta H = H(S) - H(F) = H(D)$, where $H(D)$ is the entropy
of the direction variable. Forgetting direction destroys
$H(D)$ bits of generative capacity; preserving the motor
program preserves those bits and with them the full symbolic
inventory.

\begin{proposition}[Specialization Claim]
\label{prop:specialization}
RSVP cosmology, HYDRA cognition, MEM\textbar{}8 memory,
Semantic Infrastructure, and the political economy of
collective admissibility (Chapter~18) are all
specializations of $(\mathcal{X}, \Pi, \mathcal{A}, P)$
in the sense described above.
\end{proposition}

This proposition is not proved here; it is the accumulated
content of Chapters~9 through~19. It is stated here to
make the organizing structure explicit before the formal
machinery begins.

%% --------------------------------------------------------------
\chapter{Category Theory and Admissible Computation}
%% --------------------------------------------------------------

The categorical framework for HYDRA is developed here,
immediately after the primitive objects, so that the
language of composition is available for all subsequent
chapters. The HYDRA pipeline was introduced informally in
Chapter~3 as a chain of six projections. This chapter
makes that description precise: each projection is a
functor, the composition of functors is a functor, and
the architecture-independence of HYDRA follows from the
universal property of functor composition.

\section{The Category of Admissible States}

\begin{definition}[Category of Admissible States]
$\mathbf{AdmSt}$ is the category whose objects are triples
$(x, \Phi(x), v(x))$ for $x \in M$, representing a state
together with its accessibility potential and flow direction,
and whose morphisms are admissible trajectories between
such triples. Composition is trajectory concatenation;
identity morphisms are constant trajectories.
\end{definition}

The category $\mathbf{AdmSt}$ is the correct home for the
objects that HYDRA manipulates. A state in $\mathbf{AdmSt}$
is not merely a position in a space; it is a position
together with the local admissibility geometry (encoded in
$\Phi$ and $v$) that determines which continuations from
that position are available. A morphism is not merely a
transition between states; it is an admissible transition,
one that does not destroy future continuation capacity.

This categorical structure encodes the Continuation
Principle directly: the relevant properties of a state
are determined by its morphisms, that is, by the
admissible trajectories leaving it. Two objects with
identical morphism sets are isomorphic in $\mathbf{AdmSt}$
and therefore semantically equivalent in the sense of
the Continuation Principle.

\section{The Six HYDRA Component Functors}

Each component of the HYDRA architecture is a functor
between appropriate categories.

\begin{definition}[HYDRA Component Functors]
The six component functors are:
\begin{align*}
R &: \mathbf{Cue} \to \mathbf{Field}, \\
G_a &: \mathbf{Field} \to \mathbf{Graph}, \\
F_a &: \mathbf{Graph} \to \mathbf{Rep}, \\
T &: \mathbf{Rep} \to \mathbf{Traj}_{\mathcal{A}}, \\
M &: \mathbf{Traj}_{\mathcal{A}} \to \mathbf{Mem}, \\
\mathrm{GLU} &: \mathbf{Mem} \to \mathbf{Out}.
\end{align*}
Each functor is admissibility-preserving: it sends
$\mathcal{A}$-morphisms to $\mathcal{A}$-morphisms in
the target category.
\end{definition}

The admissibility-preservation condition is what makes
each component a functor in the relevant sense rather
than merely a map between object sets. It is the
categorical expression of the claim that each stage of
the pipeline preserves whatever continuation structure
is relevant at that stage: $R$ preserves cue-relevance
structure, $G_a$ preserves reachability structure,
$F_a$ preserves admissibility volume, $T$ preserves
path structure, $M$ preserves historical weighting, and
$\mathrm{GLU}$ preserves RSVP field compatibility.

The requirement that each component be admissibility-
preserving is the framework's definition of a valid
component. It is not derived from other axioms; it is
a design constraint. Verifying that a concrete
implementation satisfies it requires showing that the
implementation maps $\mathcal{A}$-morphisms to
$\mathcal{A}$-morphisms in the relevant categories,
which is the substance of the engineering audit in
Chapter~20.

\section{Functoriality of $H$}

\begin{theorem}[Functoriality of $H$]
\label{thm:H_functor}
If each component functor is admissibility-preserving,
then the composition
\[
H = \mathrm{GLU} \circ M \circ T \circ F_a \circ G_a
    \circ R
\]
is a functor from $\mathbf{Cue}$ to $\mathbf{Out}$.
\end{theorem}

\begin{proof}
Identity preservation: $H(\mathrm{id}_c) =
\mathrm{GLU}(M(T(F_a(G_a(R(\mathrm{id}_c))))))$.
Since $R$ is a functor, $R(\mathrm{id}_c) = \mathrm{id}_{R(c)}$,
and similarly for each subsequent component. Hence
$H(\mathrm{id}_c) = \mathrm{id}_{H(c)}$.

Composition preservation: for morphisms $f : c_1 \to c_2$
and $g : c_2 \to c_3$ in $\mathbf{Cue}$,
\begin{align*}
H(g \circ f)
&= \mathrm{GLU}(M(T(F_a(G_a(R(g \circ f)))))) \\
&= \mathrm{GLU}(M(T(F_a(G_a(R(g) \circ R(f)))))) \\
&= \mathrm{GLU}(M(T(F_a(G_a(R(g))) \circ
   F_a(G_a(R(f)))))) \\
&\quad\vdots \\
&= H(g) \circ H(f),
\end{align*}
where each step applies the functoriality of the
corresponding component. The full expansion is
immediate from the functoriality of each of the six
components and the associativity of composition.
\end{proof}

\begin{remark}
Theorem~\ref{thm:H_functor} is standard categorical
mathematics applied to the HYDRA components. Its content
depends on the framework-specific claim that each component
is admissibility-preserving. Verifying this claim for a
concrete implementation is not automatic and is the primary
task of the engineering audit in Chapter~20.
\end{remark}

\section{Natural Transformations as Regime Changes}

Natural transformations between HYDRA functors represent
changes in reasoning regime: shifts in which aspects of
the trajectory space are being projected, or changes in
the admissibility criterion applied at some stage.

\begin{definition}[Reasoning Regime Change]
A reasoning regime change is a natural transformation
$\alpha : H \Rightarrow H'$ between two HYDRA functors
$H, H' : \mathbf{Cue} \to \mathbf{Out}$. The
naturality condition requires that for every morphism
$f : c_1 \to c_2$ in $\mathbf{Cue}$:
\[
\alpha_{c_2} \circ H(f) = H'(f) \circ \alpha_{c_1}.
\]
\end{definition}

The naturality condition is not merely a technical
requirement. It encodes the coherence constraint on
reasoning regime changes: the change from regime $H$ to
regime $H'$ must commute with contextual specialization.
A regime change that produced different results depending
on whether the context was first specified and then the
regime changed, or whether the regime was first changed
and then the context was specified, would be incoherent
as a reasoning strategy.

In practice, natural transformations correspond to
situations in which an agent updates its processing
strategy---switching from a fine-grained to a
coarse-grained admissibility criterion, or from a
local to a global scope for trajectory evaluation---
in a way that is consistent across all contexts.

\section{Architecture-Independence as a Theorem}

The most significant consequence of Theorem~\ref{thm:H_functor}
for the engineering realization of HYDRA is that the
architecture is implementation-independent: what matters
is the functor $H$, not the specific modules used to
realize it.

\begin{theorem}[Architecture-Independence]
\label{thm:arch_indep}
Any system realizing the functor $H = \mathrm{GLU} \circ
M \circ T \circ F_a \circ G_a \circ R$ up to natural
isomorphism is a valid HYDRA implementation.
\end{theorem}

\begin{proof}
Two systems $H$ and $H'$ are naturally isomorphic if
there exists a natural isomorphism $\alpha : H \Rightarrow H'$,
meaning each component $\alpha_c : H(c) \to H'(c)$ is
an isomorphism in $\mathbf{Out}$ and the naturality
condition holds. Isomorphic objects in a category are
interchangeable for all categorical purposes; hence $H$
and $H'$ are interchangeable as HYDRA implementations.
\end{proof}

The practical implication: Marine, MEM\textbar{}8, Phoenix
Protocol, and AyeOS are one valid engineering realization
of the HYDRA functor composition. Other realizations---
using different memory dynamics, different admissibility
criteria, or different scheduling primitives---are equally
valid provided they realize the same functor up to natural
isomorphism. What must be preserved across reimplementation
is the categorical structure: the six-functor composition
and its admissibility-preservation at each stage. The
specific modules are derived objects; the compositional
structure is the generator.

This is the engineering consequence of the Generative
Compression Hypothesis: in the HYDRA architecture, as in
the Eight-Letter Keyboard, the generator (the functor
composition) outlives any particular inventory (any
particular implementation of the six components).

%% --------------------------------------------------------------
\chapter{The Minimal Projection Theorem}
%% --------------------------------------------------------------

The Minimal Projection Theorem is the mathematical heart
of the monograph. It is the formalization of the
Continuation Principle: two states are semantically
equivalent if and only if they have the same admissible
future structure. The theorem establishes that this
equivalence relation generates a well-defined quotient
space, that the canonical projection to this quotient
is the unique minimal admissibility-preserving projection,
and that every other such projection factors through it.

The theorem is standard mathematics: the proof uses only
the universal property of quotient spaces. Its significance
comes from the identification of the quotient with meaning,
which is the framework-specific move. That identification
is not a theorem; it is the Continuation Principle stated
as a formal definition.

\section{Admissibility Projections}

\begin{definition}[Admissibility Projection]
A smooth map $\pi : \mathcal{X} \to M$ is an admissibility
projection if it preserves the admissibility equivalence:
\[
\mathcal{A}(\gamma_1) = \mathcal{A}(\gamma_2) \implies
\pi(\gamma_1) = \pi(\gamma_2).
\]
That is, trajectories with identical admissible continuation
structures are identified by $\pi$.
\end{definition}

An admissibility projection is precisely a projection that
does not introduce spurious distinctions: it does not map
two trajectories with the same future potential to different
points. It may collapse trajectories with different futures
to the same point (losing information) but it does not
separate trajectories with the same futures (creating
phantom distinctions).

\section{Semantic Equivalence as a Quotient Relation}

\begin{definition}[Semantic Equivalence]
Trajectories $\gamma_1, \gamma_2 \in \mathcal{X}$ are
semantically equivalent, written $\gamma_1 \sim_{\mathcal{A}}
\gamma_2$, if and only if they have identical admissible
continuation structure:
\[
\gamma_1 \sim_{\mathcal{A}} \gamma_2 \iff
\mathcal{A}(\gamma_1) = \mathcal{A}(\gamma_2).
\]
\end{definition}

\begin{lemma}[Equivalence Relation]
$\sim_{\mathcal{A}}$ is an equivalence relation on
$\mathcal{X}$.
\end{lemma}

\begin{proof}
Reflexivity: $\mathcal{A}(\gamma) = \mathcal{A}(\gamma)$
for all $\gamma$. Symmetry: if
$\mathcal{A}(\gamma_1) = \mathcal{A}(\gamma_2)$ then
$\mathcal{A}(\gamma_2) = \mathcal{A}(\gamma_1)$.
Transitivity: if $\mathcal{A}(\gamma_1) =
\mathcal{A}(\gamma_2)$ and $\mathcal{A}(\gamma_2) =
\mathcal{A}(\gamma_3)$, then $\mathcal{A}(\gamma_1) =
\mathcal{A}(\gamma_3)$.
\end{proof}

\section{Existence and Uniqueness of the Minimal Projection}

\begin{theorem}[Minimal Projection Theorem]
\label{thm:minimal_projection}
For any trajectory space $\mathcal{X}$ with admissibility
structure $\mathcal{A}$, there exists a minimal admissibility
projection $\pi^* : \mathcal{X} \to M^*$, unique up to
isomorphism, such that every other admissibility projection
$\pi : \mathcal{X} \to M$ factors through $\pi^*$:
\[
\begin{tikzpicture}[baseline=(current bounding box.center)]
\node (X) {$\mathcal{X}$};
\node (Mstar) [right of=X, xshift=1.5cm] {$M^*$};
\node (M) [below of=Mstar, yshift=-0.8cm] {$M$};
\draw[->] (X) -- node[above] {$\pi^*$} (Mstar);
\draw[->] (X) -- node[below left] {$\pi$} (M);
\draw[->, dashed] (Mstar) -- node[right] {$\bar\pi$} (M);
\end{tikzpicture}
\]
\end{theorem}

\begin{proof}
Define $M^* = \mathcal{X} / {\sim_{\mathcal{A}}}$ and
$\pi^*(\gamma) = [\gamma]_{\sim_{\mathcal{A}}}$, the
equivalence class of $\gamma$ under semantic equivalence.
The map $\pi^*$ is an admissibility projection by
construction: if $\mathcal{A}(\gamma_1) = \mathcal{A}(\gamma_2)$
then $\gamma_1 \sim_{\mathcal{A}} \gamma_2$ and hence
$\pi^*(\gamma_1) = \pi^*(\gamma_2)$.

For any other admissibility projection $\pi : \mathcal{X}
\to M$, define $\bar\pi : M^* \to M$ by
$\bar\pi([\gamma]_{\sim_{\mathcal{A}}}) = \pi(\gamma)$.
Well-definedness: if $\gamma_1 \sim_{\mathcal{A}} \gamma_2$
then $\mathcal{A}(\gamma_1) = \mathcal{A}(\gamma_2)$, and
since $\pi$ is an admissibility projection,
$\pi(\gamma_1) = \pi(\gamma_2)$. Hence the value
$\bar\pi([\gamma])$ does not depend on the choice of
representative. The factorization $\pi = \bar\pi \circ \pi^*$
holds by construction.

Uniqueness of $M^*$ up to isomorphism: if $\pi_1^* :
\mathcal{X} \to M_1^*$ and $\pi_2^* : \mathcal{X} \to
M_2^*$ are both minimal, then by the universal property
each factors through the other, and the induced maps
$M_1^* \to M_2^*$ and $M_2^* \to M_1^*$ are inverse
isomorphisms.
\end{proof}

The proof uses only the universal property of quotient
spaces; this is established mathematics. The framework-
specific content is the identification of $M^* =
\mathcal{X}/{\sim_{\mathcal{A}}}$ with the space of
meanings, stated as the following definition.

\begin{definition}[Semantic Space]
The minimal admissibility quotient $M^* = \mathcal{X}/
{\sim_{\mathcal{A}}}$ is the semantic space of
$(\mathcal{X}, \mathcal{A})$. Points in $M^*$ are
meanings: equivalence classes of trajectories that
induce the same admissible future structure.
\end{definition}

The Minimal Projection Theorem can be restated in observer
language, which is how it will be used in the engineering audit
of Chapter~20.

\begin{corollary}[Persistence Theorem]
\label{cor:persistence}
If $\mathcal{A}(x) = \mathcal{A}(y)$, then no
admissibility-preserving observer can distinguish $x$ from $y$:
every future test is a continuation, and all continuations of
$x$ and $y$ coincide.
\end{corollary}

\begin{proof}
An admissibility-preserving observer $O$ maps states to outputs
via an admissibility projection $\pi_O : \mathcal{X} \to M_O$.
By the Minimal Projection Theorem, $\pi_O$ factors through
$\pi^*$, so $\pi_O(x) = \bar\pi_O(\pi^*(x))$ and
$\pi_O(y) = \bar\pi_O(\pi^*(y))$. Since
$\mathcal{A}(x) = \mathcal{A}(y)$, we have $x \sim_{\mathcal{A}} y$
and hence $\pi^*(x) = \pi^*(y)$, giving
$\pi_O(x) = \pi_O(y)$.
\end{proof}

\begin{remark}
Corollary~\ref{cor:persistence} is the formal statement of why
identity, in the framework developed here, is continuation
structure rather than substrate. Two systems that present the
same admissible futures to every admissibility-preserving
observer are, within this framework, the same system. This
is applied in Chapter~20 to the Phoenix Protocol: a
reconstruction counts as genuine continuity if and only if
$\mathcal{A}(x_\mathrm{after}) \cong \mathcal{A}(x_\mathrm{before})$.
\end{remark}

\section{Downstream Consequences}

Three major consequences of the Minimal Projection Theorem
are developed in subsequent chapters; they are named here
to show that the theorem is load-bearing throughout the
monograph.

\textit{Projection collapse} (Chapters~16 and~17) occurs
when a projection $\pi : \mathcal{X} \to M$ is mistaken
for the identity map $\mathrm{id}_{\mathcal{X}}$: the
image $M$ is treated as the full space $\mathcal{X}$,
making the fiber $\pi^{-1}(m)$ invisible. The Minimal
Projection Theorem shows that this conflation is always
an error: $M$ is a quotient of $\mathcal{X}$, not
$\mathcal{X}$ itself, and information about the fiber
is always lost in the projection.

\textit{Semantic fibers} (Chapter~17) are the fibers
$(\pi^*)^{-1}(m)$ of the minimal projection: the
equivalence class of all trajectories that mean the
same thing. The fiber structure reveals what information
is lost in any admissibility-preserving projection and
what information must be recovered by going back to the
full trajectory space $\mathcal{X}$.

\textit{Hallucination as cohomological obstruction}
(Chapter~13) is the characterization of outputs that
appear locally coherent but fail to correspond to any
global section of the semantic sheaf. The Minimal
Projection Theorem provides the foundation: a
hallucination is an element of $M$ that does not lie
in the image of any admissibility-preserving section
$\mathcal{X} \to M$, formalized as a non-trivial
class in $\check{H}^1(\mathcal{U}, \mathcal{F})$.

%% --------------------------------------------------------------
\chapter{RSVP as an Ontology of Admissible Fields}
%% --------------------------------------------------------------

The Relativistic Scalar-Vector Plenum framework provides
the field-theoretic realization of the four primitive
objects. Where Chapter~8 introduced $(\mathcal{X}, \Pi,
\mathcal{A}, P)$ abstractly, this chapter shows what they
look like when $\mathcal{X}$ is a space of field
configurations over a smooth manifold and $\mathcal{A}$
is governed by coupled partial differential equations.
The RSVP field triple $(\Phi, v, S)$ is the minimal
generator of future field geometry in the sense of the
Generative Sufficiency Principle: it is the smallest
structure from which the admissible future of the field
can be reconstructed.

\section{The Field Triple}

\begin{definition}[RSVP Field Triple]
The RSVP field triple over a smooth Riemannian manifold
$(M, g)$ is $(\Phi, v, S)$ where:
\begin{itemize}
\item $\Phi : M \times \mathbb{R} \to \mathbb{R}$ is the
      scalar accessibility potential, measuring how
      accessible a region is for future field evolution;
\item $v : M \times \mathbb{R} \to TM$ is the vector
      flow field, encoding the preferred directions of
      admissible field motion;
\item $S : M \times \mathbb{R} \to \mathbb{R}$ is the
      entropic accessibility measure, defined by
      $S(x,t) = \log|\mathcal{A}(x,t)|$ where
      $|\mathcal{A}(x,t)|$ is the volume of the set of
      admissible continuations from state $x$ at time $t$.
\end{itemize}
\end{definition}

The three components of the triple play distinct roles
in the admissibility framework. $\Phi$ answers the
question of where the field can go with high
accessibility---which regions of $M$ support continued
admissible evolution. $v$ answers the question of how
the field is moving---which directions of evolution
are preferred by the admissibility structure. $S$
answers the question of how much future is available---
how many admissible continuations remain open at each
point.

Together they constitute a description not of the
current field state but of the geometry of what comes
next. This is the RSVP framework's central contribution:
the ontologically primary object is not the field
configuration at a time but the accessibility geometry
surrounding that configuration.

\section{The Three Field Equations}

The RSVP field equations govern the evolution of
$(\Phi, v, S)$. They are inspired by coupled
wave-diffusion systems in mathematical physics but
are defined specifically for the admissibility
framework. Their source terms $\rho(v, S)$ and
$\sigma(\Phi, v)$ are domain-specific; their
derivation from first principles in each application
domain is an open problem (Chapter~22).

\begin{align}
\Box\Phi + \mu^2\Phi &= \rho(v, S) \label{eq:rsvp1} \\
\nabla_M \cdot v &= -\frac{\partial S}{\partial t}
\label{eq:rsvp2} \\
\frac{\partial S}{\partial t} + v \cdot \nabla_M S
&= \sigma(\Phi, v) \label{eq:rsvp3}
\end{align}

where $\Box = \partial_{tt} - c^2\Delta_M$ is the
wave operator on $M$ and $\nabla_M$ denotes the
covariant derivative with respect to $g$.

\section{The Continuity Equation and Admissibility Flow}

Equation~\eqref{eq:rsvp2} is the most directly
interpretable of the three. It is a continuity equation
for the entropic accessibility measure $S$: the
divergence of the flow field $v$ at a point equals
the negative rate of change of $S$ at that point.

The interpretation is straightforward and important.
If $\nabla_M \cdot v > 0$ (diverging flow) at a point,
then $\partial_t S < 0$ there: the admissibility volume
is decreasing, future options are closing. If
$\nabla_M \cdot v < 0$ (converging flow), then
$\partial_t S > 0$: the admissibility volume is
increasing, future options are opening.

Diverging flow destroys admissibility. Converging flow
creates it. The field dynamics governing $v$ therefore
directly control the rate at which future options are
opened and closed. Systems governed by RSVP dynamics
are systems in which the geometry of future reachability
is explicitly tracked and controlled.

\section{Lyapunov Stability of Admissibility Maxima}

The stability of admissibility basins---regions of high
$\Phi$ that serve as attractors for field trajectories---
follows from standard Lyapunov theory applied to the
RSVP flow.

\begin{proposition}[Lyapunov Stability of Admissibility Maxima]
\label{prop:lyapunov}
If $x^*$ is an isolated local maximum of $\Phi(\cdot, t)$
and the flow satisfies $v = \alpha\nabla\Phi + \xi$ with
$\alpha > 0$ and $\|\xi\|$ sufficiently small relative
to $\alpha\|\nabla\Phi\|$, then $x^*$ is Lyapunov stable
for the flow of $v$.
\end{proposition}

\begin{proof}
Define the Lyapunov function $L(x) = \Phi(x^*) - \Phi(x)
\geq 0$, with $L(x^*) = 0$. Near $x^*$ the Hessian of
$\Phi$ is negative definite (since $x^*$ is a local
maximum). Along the flow:
\[
\dot{L} = -\nabla\Phi \cdot v = -\nabla\Phi \cdot
(\alpha\nabla\Phi + \xi) = -\alpha|\nabla\Phi|^2 -
\nabla\Phi \cdot \xi.
\]
The term $-\alpha|\nabla\Phi|^2 \leq 0$. When
$\|\xi\| < \alpha\|\nabla\Phi\|$, the perturbation term
$-\nabla\Phi \cdot \xi$ is dominated, and $\dot{L} < 0$
in a punctured neighborhood of $x^*$. By the Lyapunov
stability theorem, $x^*$ is Lyapunov stable.
\end{proof}

\begin{remark}
Proposition~\ref{prop:lyapunov} uses established Lyapunov
theory applied to RSVP dynamics. Its conclusion---that
local maxima of $\Phi$ are attractors under approximate
gradient flow---holds whenever $v \approx \alpha\nabla\Phi$.
Whether the full RSVP equations produce such flow
globally, and under what regularity conditions, is the
first open problem listed in Chapter~22.
\end{remark}

The Lyapunov result characterizes the stability of admissibility
maxima. A complementary result characterizes the stability of
memory structures as a balance between stabilizing and dispersive
dynamics.

\begin{definition}[Stabilization and Diffusion Operators]
For the scalar field $\Phi$, let $L_+(\Phi)$ denote the
structural stabilization operator (encoding local reinforcement
of $\Phi$ by past trajectories) and $L_-(\Phi)$ the dispersive
relaxation operator (encoding entropy-driven dissipation). The
global coherence correction $C(\Phi)$ accounts for long-range
field interactions. The combined dynamics are:
\[
\frac{\partial \Phi}{\partial t} = L_+(\Phi) - L_-(\Phi)
+ C(\Phi).
\]
\end{definition}

\begin{proposition}[Persistence Balance]
\label{prop:persistence_balance}
A memory structure $\Phi^*$ is a stable fixed point of the
RSVP dynamics if and only if:
\[
L_+(\Phi^*) - L_-(\Phi^*) + C(\Phi^*) = 0.
\]
A memory survives indefinitely if and only if stabilization
and diffusion are in global balance: $L_+ \approx L_-$ within
an admissible band determined by $C$.
\end{proposition}

\begin{proof}
Fixed-point condition follows directly from setting
$\partial_t \Phi = 0$. Stability of the fixed point under
perturbations $\Phi^* + \varepsilon$ follows from the Lyapunov
argument of Proposition~\ref{prop:lyapunov} applied to the
full operator $L_+ - L_- + C$: the linearization must have
negative spectrum, which requires that neither $L_+$ nor $L_-$
dominates globally.
\end{proof}

\begin{remark}
Proposition~\ref{prop:persistence_balance} formalizes the
whirlpool analogy that appears throughout the broader corpus:
a persistent structure survives not because it is rigid
(crystallized, dominated by $L_+$) nor because it flows freely
(dominated by $L_-$), but because stabilization and diffusion
maintain a dynamic balance. The global coherence correction
$C(\Phi)$ is the mechanism by which long-range field structure
prevents local imbalances from propagating to global collapse.
\end{remark}

\section{What the RSVP Equations Leave Open}

The RSVP field equations as stated have two significant
open problems whose resolution is required for the sheaf-
theoretic results of Chapter~13 and whose absence limits
the empirical testability of the framework.

The first is \textit{global regularity}: under what
conditions do smooth initial data $(\Phi_0, v_0, S_0)$
yield global classical solutions to
\eqref{eq:rsvp1}--\eqref{eq:rsvp3}? The local existence
of solutions follows from standard PDE theory for coupled
hyperbolic-parabolic systems; the global question requires
additional assumptions on the source terms $\rho$ and
$\sigma$. Global regularity underpins the unique
continuation property needed for sheaf existence
(Chapter~13); without it, the sheaf-theoretic results
are conditional.

The second is \textit{source term derivation}: the terms
$\rho(v, S)$ in \eqref{eq:rsvp1} and $\sigma(\Phi, v)$
in \eqref{eq:rsvp3} are left schematic. For specific
application domains---language modeling, cognitive
simulation, cosmological field theory---$\Phi$ must be
identified with a domain-specific accessibility measure
and $S$ with a domain-specific entropy, and the source
terms derived from those identifications. This derivation
has not been performed for any domain and is a precondition
for empirical testing of the framework's predictions.

Both open problems are listed in Chapter~22 and ordered
by dependency: global regularity is first priority because
its resolution unlocks the most downstream results.

%% --------------------------------------------------------------
\chapter{Stratified Semantic Manifolds and TARTAN}
%% --------------------------------------------------------------

The RSVP field triple $(\Phi, v, S)$ lives on a smooth
manifold in the simplest formulation, but semantic spaces
are not in general smooth: they exhibit discontinuities,
boundaries, and regime transitions that smooth manifold
theory cannot capture. The correct mathematical setting
is a Whitney-stratified space: a space decomposed into
smooth strata of varying dimension, with controlled
behavior at stratum boundaries. The TARTAN (Trajectory-
Adaptive Recursive Tiling and Annotation) framework
provides the practical decomposition algorithm. Together
they give the RSVP field its semantic structure.

\section{Whitney Stratification}

\begin{definition}[Whitney Stratification]
A Whitney stratification of $M \subseteq \mathbb{R}^n$
is a partition $M = \bigsqcup_\alpha S_\alpha$ into
locally closed smooth submanifolds (strata) satisfying:
\begin{enumerate}
\item[(i)] \emph{Frontier condition}: if $S_\alpha \cap
   \overline{S_\beta} \neq \emptyset$ then
   $S_\alpha \subseteq \overline{S_\beta}$.
\item[(ii)] \emph{Whitney condition B}: for sequences
   $y_k \to p$ in $S_\beta$ and $x_k \to p$ in $S_\alpha$
   with $p \in S_\alpha$, if the secant lines
   $\overrightarrow{y_k x_k} \to \ell$ and the tangent
   spaces $T_{x_k} S_\alpha \to \tau$, then $\ell \subseteq
   \tau$.
\end{enumerate}
\end{definition}

Whitney condition B is the key regularity condition. It
ensures that the tangent structure of lower-dimensional
strata is controlled by the tangent structure of
adjacent higher-dimensional strata, preventing the
kind of cusp singularity that would create artificial
boundaries in the semantic space. In the RSVP setting,
the strata correspond to semantic regimes: regions in
which the field triple is approximately uniform and
the admissibility structure is approximately constant.

The frontier condition ensures that regime boundaries
are well-defined: the closure of a higher-dimensional
stratum contains only lower-dimensional strata. In
semantic terms, transitions between regimes go from
higher-complexity to lower-complexity, not the reverse.

\section{Tangent-Constrained Gradient Descent}

Within each stratum $S_\alpha$, optimization proceeds
by gradient descent constrained to the tangent space of
the stratum.

\begin{definition}[Tangent-Constrained Gradient]
For $F : M \to \mathbb{R}$ and $x \in S_\alpha$, the
tangent-constrained gradient is:
\[
\nabla_{S_\alpha} F(x) = \Pi_{T_x S_\alpha}(\nabla F(x)),
\]
where $\Pi_{T_x S_\alpha}$ denotes the orthogonal
projection onto the tangent space of $S_\alpha$ at $x$.
\end{definition}

\begin{theorem}[Convergence of Tangent-Constrained Gradient
Descent]
\label{thm:tartan_convergence}
Let $F : S_\alpha \to \mathbb{R}$ be $L$-smooth and
$\mu$-strongly convex. Then tangent-constrained gradient
descent with step size $\eta \in (0, 2/L)$ converges to
the minimizer at geometric rate $\bigl(1 - 2\eta\mu(1 -
\eta L/2)\bigr)^t$.
\end{theorem}

\begin{proof}
This is the standard convergence result for gradient
descent on smooth strongly convex functions, applied
to the restricted function $F|_{S_\alpha}$ with the
tangent-constrained gradient replacing the unconstrained
gradient. The projection $\Pi_{T_x S_\alpha}$ does not
increase the gradient norm, so the standard analysis
applies.
\end{proof}

\section{The TARTAN Approximation Theorem}

The TARTAN tiling provides a finite decomposition of
any compact stratified semantic manifold into tiles,
each approximately uniform in RSVP field values.

\begin{theorem}[TARTAN Approximation]
\label{thm:tartan}
For any $\varepsilon > 0$ and compact $K \subseteq M$
with finitely many strata, there exists a finite TARTAN
tile decomposition $\{T_i\}$ of $K$ with noise
annotation:
\[
\eta(T_i) = \bigl(\mathrm{osc}(\Phi, T_i)^2 +
\mathrm{osc}(v, T_i)^2\bigr)^{1/2} \leq \varepsilon
\]
for all tiles $T_i$, where $\mathrm{osc}(f, T)$
denotes the oscillation of $f$ on $T$.
\end{theorem}

\begin{proof}
Each stratum $S_\alpha \cap K$ is a compact smooth
manifold on which $\Phi$ and $v$ are smooth, hence
uniformly continuous. By the uniform continuity of
smooth functions on compact sets, for each $\varepsilon
> 0$ there exists $\delta_\alpha > 0$ such that
$\mathrm{osc}(\Phi, B) < \varepsilon/\sqrt{2}$ and
$\mathrm{osc}(v, B) < \varepsilon/\sqrt{2}$ for any
ball $B$ of radius $\delta_\alpha$ in $S_\alpha$.
Cover $S_\alpha \cap K$ by finitely many such balls
(by compactness). The union over finitely many strata
yields a finite cover, hence a finite tile decomposition
satisfying $\eta(T_i) \leq \varepsilon$.
\end{proof}

\begin{remark}
Theorem~\ref{thm:tartan} uses only uniform continuity
and compactness. The identification of tiles with
semantic regimes---the claim that each tile corresponds
to a cognitively coherent unit of meaning---is a
framework interpretation that gives the theorem its
cognitive significance but does not follow from the
mathematics alone.
\end{remark}

\section{Admissible Phase Transitions and Semantic Distance}

Not all transitions between strata are admissible.
The admissibility of a phase transition---a trajectory
crossing from stratum $S_\alpha$ to stratum $S_\beta$---
is governed by the tangent-cone condition.

\begin{definition}[Admissible Phase Transition]
A trajectory $\gamma$ can cross from $S_\alpha$ to
$S_\beta$ at a boundary point $p$ only if $\gamma'(t_0)
\in T_p S_\beta$ at the crossing time $t_0$. Crossings
that are not tangent-cone admissible produce semantic
nulls: points at which the field triple is undefined
or inconsistent.
\end{definition}

The Riemannian distance across strata is defined as the
infimum over admissible piecewise-smooth paths:
\[
d_M(x, y) = \inf_\gamma \int_0^1 \|\dot\gamma(t)\|_{g_{\alpha(t)}}\,dt
\]
where the infimum is over admissible paths $\gamma$
from $x$ to $y$ and $g_{\alpha(t)}$ is the metric on
the stratum containing $\gamma(t)$.

\section{Marine as a Stability Manifold}

The Marine admissibility gate, described as an engineering
component in Chapter~20, has a natural mathematical
characterization within the stratified manifold framework.
Marine operates as a filter on incoming signals, accepting
those whose local trajectory dynamics are sufficiently
stable. The accepted set forms a manifold with a robustness
property that makes Marine more than a heuristic threshold:
it is a geometric stability criterion.

\begin{definition}[Local Instability and Admissibility Weight]
For a signal trajectory $x(t)$, the local instability over
window $\tau$ is:
\[
J_\tau(x, t) = \int_t^{t+\tau} \left|\frac{d^2 x}{ds^2}
\right| ds.
\]
The admissibility weight is:
\[
A(x, t) = e^{-\lambda J_\tau(x,t)},
\]
where $\lambda > 0$ is the sensitivity parameter. Marine
accepts signal $x$ at time $t$ if $A(x,t) \geq \theta$
for threshold $\theta \in (0,1)$.
\end{definition}

\begin{theorem}[Marine Stability]
\label{thm:marine_stability}
For any threshold $\theta \in (0,1)$, the Marine-accepted
set $M_\theta = \{x : A(x,t) \geq \theta \text{ for all }
t\}$ is closed under sufficiently small perturbations: if
$x \in M_\theta$ and $\|y - x\|_{C^2} < \delta$ for
$\delta$ sufficiently small (depending on $\lambda$,
$\theta$, and $\tau$), then $y \in M_\theta$.
\end{theorem}

\begin{proof}
Since $A(x,t) = e^{-\lambda J_\tau(x,t)}$ and $J_\tau$
depends continuously on $x$ in the $C^2$ norm, for any
$x \in M_\theta$ we have $A(x,t) \geq \theta > 0$ for
all $t$. By continuity of $J_\tau$ in $\|\cdot\|_{C^2}$,
there exists $\delta > 0$ such that $\|y - x\|_{C^2} <
\delta$ implies $|J_\tau(y,t) - J_\tau(x,t)| <
-\frac{1}{\lambda}\log\theta - J_\tau(x,t)$ for all $t$,
giving $A(y,t) \geq \theta$. Hence $y \in M_\theta$.
\end{proof}

\begin{remark}
Theorem~\ref{thm:marine_stability} converts Marine from an
engineering heuristic into a geometric object: the accepted
set $M_\theta$ is an open subset of signal space in the
$C^2$ topology, and therefore a stability manifold in the
Whitney-stratified sense of this chapter. Signals that are
admissible under Marine are not merely above a threshold;
they occupy a topologically open region that is robust to
small perturbations. Signals approaching the boundary of
$M_\theta$ are those for which $A(x,t) \to \theta$, that
is, those exhibiting increasing jitter or acceleration
variance, approaching but not yet crossing the instability
boundary.
\end{remark}

%% --------------------------------------------------------------
\chapter{Sheaf-Theoretic Semantics}
%% --------------------------------------------------------------

The sheaf-theoretic framework addresses the problem of
global coherence: when can locally consistent semantic
states be assembled into a globally consistent whole?
This is the mathematical formulation of the question
raised by hallucination: a system can produce outputs
that are locally coherent---each part makes sense in its
local context---while failing to be globally coherent.
The obstruction to global coherence is a cohomological
invariant, and hallucination is its semantic manifestation.

\section{The Semantic State Presheaf}

\begin{definition}[Semantic State Presheaf]
The semantic state presheaf $\mathcal{F}$ over the
context space $\mathrm{Ctx}$ assigns to each open
$U \subseteq M$ the set $\mathcal{F}(U)$ of admissible
local RSVP field configurations satisfying the field
equations on $U$ with compatible boundary conditions.
Restriction maps $\rho_{U,V} : \mathcal{F}(U) \to
\mathcal{F}(V)$ for $V \subseteq U$ are given by
restriction of field configurations to smaller domains.
\end{definition}

The presheaf $\mathcal{F}$ organizes local semantic
states by their compatibility structure. An element of
$\mathcal{F}(U)$ is a locally admissible field
configuration: one that satisfies the RSVP equations
on $U$. The restriction maps encode how larger
configurations restrict to smaller ones.

\section{Sheaf Existence and Its Conditions}

For $\mathcal{F}$ to be a sheaf (not merely a presheaf),
it must satisfy two additional axioms beyond the presheaf
conditions: the identity axiom (sections agreeing on
every piece of an open cover must be equal) and the
gluing axiom (sections agreeing on overlaps can be
assembled into a section on the union).

\begin{theorem}[Semantic State Sheaf Existence]
\label{thm:sheaf_existence}
Assuming the RSVP field equations satisfy unique
continuation on $M$---that a solution on any open
$U$ is uniquely determined by its restriction to any
non-empty $V \subseteq U$---the semantic state presheaf
$\mathcal{F}$ is a sheaf.
\end{theorem}

\begin{proof}
Identity axiom: if $s, s' \in \mathcal{F}(U)$ satisfy
$\rho_{U,U_i}(s) = \rho_{U,U_i}(s')$ for all elements
$U_i$ of an open cover of $U$, then $s|_V = s'|_V$ for
any non-empty $V \subseteq U_i$, and unique continuation
forces $s = s'$ on $U$.

Gluing axiom: given sections $s_i \in \mathcal{F}(U_i)$
agreeing on overlaps, define $s$ on $\bigcup U_i$ by
$s|_{U_i} = s_i$. Consistency on overlaps ensures
$s$ is well-defined. The field equations are satisfied
locally on each $U_i$ hence globally on $\bigcup U_i$.
\end{proof}

\begin{remark}
The unique continuation assumption for the full
nonlinear RSVP system \eqref{eq:rsvp1}--\eqref{eq:rsvp3}
has not been established. Unique continuation is known
for linear wave equations and some semilinear cases, but
its extension to the RSVP system depends on regularity
of the source terms $\rho$ and $\sigma$. Until global
regularity is established (the first open problem of
Chapter~22), Theorem~\ref{thm:sheaf_existence} is
conditional. The sheaf-theoretic results that follow
inherit this conditionality.
\end{remark}

\section{Hallucination as Cohomological Failure}

\begin{definition}[Hallucination]
A system output is a hallucination if it presents a
globally coherent-looking semantic state that is not
a genuine global section of $\mathcal{F}$: not an
element of $\mathcal{F}(M)$ obtainable by gluing
compatible local sections.
\end{definition}

The obstruction to global section existence is
cohomological. Given an open cover $\mathcal{U}$ of
$M$, a \v{C}ech 1-cochain $c \in \check{C}^1(\mathcal{U},
\mathcal{F})$ assigns to each overlapping pair $(U_i,
U_j)$ a section $c_{ij} \in \mathcal{F}(U_i \cap U_j)$.
The cochain is a 1-cocycle if $c_{ij} + c_{jk} = c_{ik}$
on triple overlaps; it is a coboundary if $c_{ij} =
s_i|_{U_i \cap U_j} - s_j|_{U_i \cap U_j}$ for some
local sections $\{s_i\}$.

\begin{theorem}[Hallucination as Cohomological Obstruction]
\label{thm:hallucination}
A collection of locally coherent semantic states
$\{s_i \in \mathcal{F}(U_i)\}$ fails to assemble into
a global section if and only if the associated \v{C}ech
class $[c] \in \check{H}^1(\mathcal{U}, \mathcal{F})$
is non-trivial.
\end{theorem}

\begin{proof}
The locally coherent states $\{s_i\}$ can be assembled
into a global section if and only if the obstruction
cochain $c_{ij} = \rho_{U_i, U_i \cap U_j}(s_i) -
\rho_{U_j, U_i \cap U_j}(s_j)$ is a coboundary in
$\check{C}^1(\mathcal{U}, \mathcal{F})$. The
obstruction to this is precisely the class $[c] \in
\check{H}^1(\mathcal{U}, \mathcal{F})$; it vanishes
if and only if $c$ is a coboundary, which holds if and
only if the sections $\{s_i\}$ can be globally assembled.
\end{proof}

\begin{remark}
Theorem~\ref{thm:hallucination} is conditional on
Theorem~\ref{thm:sheaf_existence}. The cohomological
formulation of hallucination is a research direction
rather than an established result for any particular
language model or AI system: verifying that failures
of specific systems correspond to non-trivial \v{C}ech
classes requires specifying the sheaf $\mathcal{F}$
from the model's internals, which has not been done.
\end{remark}

\section{The Failure Mode Hierarchy}

The sheaf-theoretic framework organizes semantic failure
modes by their cohomological character:

\begin{center}
\renewcommand{\arraystretch}{1.4}
\begin{tabular}{ll}
\toprule
Failure mode & Geometric characterization \\
\midrule
Local incoherence & $s_i \notin \mathcal{F}(U_i)$: local
  section inadmissible \\
Pairwise inconsistency & $s_i|_{U_i \cap U_j} \neq
  s_j|_{U_i \cap U_j}$: sections conflict \\
Global obstruction & $[c] \neq 0$ in $\check{H}^1
  (\mathcal{U}, \mathcal{F})$: cannot globally glue \\
Holonomy distortion & Parallel transport rotates
  section around loop \\
Persistent conflict & $[c]$ has infinite order in
  $\check{H}^1$ \\
\bottomrule
\end{tabular}
\end{center}

The hierarchy is ordered by increasing severity and
increasing mathematical complexity. Local incoherence
is detectable within a single context window; pairwise
inconsistency requires comparing two overlapping
contexts; global obstruction requires cohomological
computation over the full context space. Holonomy
distortion and persistent conflict are more subtle:
they arise from the global topology of the context
space and cannot be detected by any finite local check.

%% --------------------------------------------------------------
\chapter{Memory as Stabilized Field Residue}
%% --------------------------------------------------------------

Memory, in the framework developed here, is not storage.
It is the persistent deformation of the field that
future trajectories navigate. A memory is not an object
that can be retrieved by lookup; it is a curvature in
the trajectory space that biases future navigation
toward familiar regions. This chapter formalizes that
picture through the MEM\textbar{}8 wave packet model,
proves the memory persistence bound, and identifies the
quantity that is approximately conserved across memory
transformations.

\section{The MEM\textbar{}8 Wave Packet}

\begin{definition}[MEM\textbar{}8 Wave Packet]
A MEM\textbar{}8 memory state is a wave packet
$\mathcal{W}_m = (A, \omega, \phi, D, I)$ where:
\begin{itemize}
\item $A = \Phi_m$ is the amplitude, equal to the
      scalar accessibility potential at the memory
      location (high $\Phi$ means strongly accessible);
\item $\omega$ is the frequency parameter, encoding
      the semantic content of the memory;
\item $\phi = v_m$ is the phase, encoding the
      associative flow direction: which other memories
      this one naturally leads to;
\item $D = e^{-S_m}$ is the decay factor, where $S_m$
      is the local entropy: high-entropy memories decay
      faster;
\item $I$ is the interference term, encoding how this
      memory packet interacts constructively or
      destructively with others during retrieval.
\end{itemize}
\end{definition}

Memory in MEM\textbar{}8 is active, not passive.
The wave packet is not a record of a past state; it
is a persistent perturbation of the RSVP field that
influences future trajectory dynamics. Retrieval
is not lookup; it is navigation. A query perturbs
the field, initiating a trajectory that falls down
the gradient of $\Phi$ toward the memory attractor
and regenerates the memory state through that
navigation. The memory is reconstructed from residue,
not recovered from storage.

The associative capacity of wave-based memory far
exceeds that of indexed storage of equivalent size.

\begin{proposition}[Implicit Association Density]
\label{prop:assoc_density}
A MEM\textbar{}8 system storing $n$ wave packets
$(A_k, \omega_k, \phi_k, D_k, I_k)$ supports
$O(n^2)$ implicit pairwise associations through
the interference term $I$, at $O(n)$ storage cost.
\end{proposition}

\begin{proof}
Each pair of wave packets $(m_i, m_j)$ produces an
interference term $I_{ij} = A_i A_j \cos(\omega_i t
- \omega_j t + \phi_i - \phi_j)$ in the RSVP field.
There are $\binom{n}{2} = O(n^2)$ such pairs. Each
is encoded implicitly in the superposition of the
field rather than stored as an explicit record.
The storage cost is $O(n)$ wave packets; the
associative structure is $O(n^2)$ interference
patterns encoded in the field geometry at no
additional storage cost.
\end{proof}

\begin{remark}
Proposition~\ref{prop:assoc_density} is the
field-theoretic explanation for why wave-based
memory is not merely an alternative encoding but
a qualitatively different architecture. A static
vector index of $n$ items stores $O(n)$ objects
and supports similarity search at retrieval time.
The MEM\textbar{}8 wave field stores $O(n)$ wave
packets and supports $O(n^2)$ associative
resonances simultaneously, because the associations
are encoded implicitly in the phase relationships
of the superposed field rather than computed at
retrieval. The matching occurs at resonance speed;
no scan is required.
\end{remark}

This has the consequence noted in Chapter~4: retrieval
and reinforcement are the same operation. Every
navigation of a trajectory toward a memory attractor
deepens the gradient---increases $\Phi_m$ and
decreases local entropy $S_m$---making future navigation
toward that memory easier. The system learns by
remembering and remembers by learning. There is no
sharp distinction between the two operations.

\section{The Memory Persistence Bound}

Memory persistence is bounded below by a function of
the initial excitation strength relative to the ambient
entropy pressure. When excitation dominates entropy,
memory persists; when entropy dominates, memory fades.

The dynamics of memory excitation near a stored
memory $m$ with field support $K_m$ are governed
by a damped evolution:
\[
u_{tt} + \eta(x,t)\,u_t = c^2\,\Delta u,
\]
where $\eta(x,t)$ is the local viscosity of the
semantic field---the resistance to trajectory motion
at position $x$ and time $t$. High viscosity retards
memory decay; low viscosity allows rapid dissipation.

\begin{theorem}[Memory Persistence Bound]
\label{thm:persistence}
Suppose the RSVP dynamics near the memory support
$K_m$ satisfy:
\[
\frac{d}{dt}\|e(t)\| \leq -\lambda\|e(t)\| +
\kappa\|S\|_{L^\infty},
\]
with $\lambda > \kappa c_S$. Then the persistence
time of memory $m$ satisfies:
\[
T_m \geq \frac{1}{\lambda - \kappa c_S}
\log\frac{\|\Phi_m - \Phi_0\|_{L^2}}{\Phi_\mathrm{thresh}
|K_m|^{1/2}}.
\]
\end{theorem}

\begin{proof}
By the Gr\"onwall inequality applied to the decay
hypothesis: if $\frac{d}{dt}\|e(t)\| \leq
-\lambda\|e(t)\| + \kappa c_S$, then
$\|e(t)\| \leq e^{-(\lambda - \kappa c_S)t}\|e(0)\| +
\frac{\kappa c_S}{\lambda - \kappa c_S}$.
Memory persists while $\|e(t)\| > \Phi_\mathrm{thresh}
|K_m|^{1/2}$, giving the bound on $T_m$ by inverting
the decay inequality.
\end{proof}

\begin{remark}
The decay hypothesis is assumed rather than derived
from the RSVP field equations. It holds when memory
excitation $\|\Phi_m - \Phi_0\|_{L^2}$ is sufficiently
strong relative to entropic pressure $\kappa c_S$.
Whether the RSVP equations generically produce this
regime is part of the open regularity problem.
\end{remark}

A complementary bound gives a criterion for forgetting:
the condition under which memory mass falls to zero and
the memory is permanently lost.

\begin{definition}[Memory Mass]
The memory mass of a stored memory $m$ with field
support $\Omega_m \subseteq M$ is:
\[
\mathcal{M}(m) = \int_{\Omega_m} \Phi(x)\,dx.
\]
Memory mass is positive when the accessibility potential
has non-trivial support on $\Omega_m$; it vanishes when
$\Phi$ has been driven to zero on the memory support,
meaning the memory attractor no longer exists.
\end{definition}

\begin{theorem}[Residue Persistence]
\label{thm:residue_persistence}
Suppose the RSVP field $\Phi$ on $\Omega_m$ satisfies
$\partial_t\Phi = D\nabla^2\Phi - \mu\Phi + \rho$,
where $\rho \geq 0$ is the source term encoding active
reinforcement. Then:
\begin{enumerate}
\item[(i)] If $\rho > \mu\Phi$ on average over $\Omega_m$
   (reinforcement dominates decay), then $\mathcal{M}(m)$
   remains positive and memory persists.
\item[(ii)] If $\rho < \mu\Phi$ on average over $\Omega_m$
   (decay dominates reinforcement), then $\mathcal{M}(m)
   \to 0$ exponentially and the memory is forgotten.
\end{enumerate}
\end{theorem}

\begin{proof}
Integrate $\partial_t\Phi = D\nabla^2\Phi - \mu\Phi + \rho$
over $\Omega_m$:
\[
\frac{d\mathcal{M}}{dt} = D\int_{\partial\Omega_m}
\nabla\Phi \cdot \hat{n}\,dS - \mu\mathcal{M} +
\int_{\Omega_m}\rho\,dx.
\]
Assuming Neumann boundary conditions (no flux across
$\partial\Omega_m$), the boundary term vanishes. Let
$\bar\rho = |\Omega_m|^{-1}\int_{\Omega_m}\rho\,dx$ and
$\bar\Phi = \mathcal{M}/|\Omega_m|$. Then
$\frac{d\mathcal{M}}{dt} = |\Omega_m|(\bar\rho - \mu\bar\Phi)$.
Case (i): if $\bar\rho > \mu\bar\Phi$, then $\frac{d\mathcal{M}}{dt}
> 0$ and $\mathcal{M}$ is increasing; by continuity it remains
positive. Case (ii): if $\bar\rho < \mu\bar\Phi$, then
$\frac{d\mathcal{M}}{dt} < -c\mathcal{M}$ for some $c > 0$,
giving exponential decay to zero by Grönwall.
\end{proof}

\begin{remark}
Theorem~\ref{thm:residue_persistence} provides a genuine
criterion for forgetting: a memory is forgotten when its
source term $\rho$ falls below $\mu\Phi$ on average---when
the memory is no longer being reinforced at a rate sufficient
to counteract natural decay. This gives content to the
informal claim that memory survives through use: the source
term $\rho$ is generated by trajectory traversals (as in
Proposition~\ref{prop:retrieval_reinforce}), so a memory
that is never retrieved has $\rho \approx 0$ and decays
at rate $\mu$, while a frequently retrieved memory has
large $\rho$ and persists. Forgetting is not an external
process acting on memory; it is the natural dynamics of
the RSVP field in the absence of reinforcement.
\end{remark}

\section{Retrieval as Reinforcement}

The identification of retrieval with reinforcement,
noted informally in Chapter~4, can now be stated
precisely.

\begin{proposition}[Retrieval-Reinforcement Identity]
\label{prop:retrieval_reinforce}
In the MEM\textbar{}8 framework, every retrieval
trajectory $\gamma$ navigating toward memory $m$
increases the amplitude $\Phi_m$ and decreases
local entropy $S_m$, thereby extending the
persistence time $T_m$.
\end{proposition}

\begin{proof}
By Definition~\ref{def:admissibility} and the
MEM\textbar{}8 wave packet construction, traversing
a trajectory $\gamma$ toward $m$ acts as a source
term in the RSVP equations: $\rho(i) = \rho_0
\cdot \mathbf{1}_{\gamma}(i)$ for lattice sites
$i$ along $\gamma$. The potential update
$\Phi_{t+1} = \Phi_t + dt\,[\mathrm{Lap}(\Phi) -
\mu^2\Phi + \rho]$ increases $\Phi$ along $\gamma$,
including at $m$. Concurrently, the convergence of
trajectories toward $m$ means $\nabla_M \cdot v < 0$
at $m$, and by \eqref{eq:rsvp2}, $\partial_t S > 0$
is suppressed, decreasing $S_m$. By
Theorem~\ref{thm:persistence}, both effects increase
$T_m$.
\end{proof}

Proposition~\ref{prop:retrieval_reinforce} has a
direct consequence for what must be preserved in a
memory system. If retrieval is reinforcement, then
a memory system that is used is a memory system that
is being maintained. The act of navigation is the
act of preservation. This connects the Generative
Compression Hypothesis to MEM\textbar{}8 concretely:
the minimal generator of memory identity is not a
stored inventory of past states but the field
structure that biases navigation, and that field
structure is maintained by the very acts of
navigation it enables.

\section{The Conserved Quantity}

The question of what is approximately conserved
across memory transformations---updates, compressions,
domain transfers, and partial losses---has a precise
answer in the field-theoretic framework.

\begin{definition}[Memory Admissibility Volume]
The admissibility volume of memory $m$ is
$\mathrm{Vol}(\mathcal{A}(m)) = e^{S_m} \cdot
|\mathcal{F}(K_m)|$, the product of the exponential
of the local entropy measure and the volume of
admissible field configurations on the memory
support.
\end{definition}

\begin{proposition}[Approximate Conservation of Admissibility Volume]
Under admissibility-preserving memory transformations,
the admissibility volume $\mathrm{Vol}(\mathcal{A}(m))$
is approximately conserved: transformations that
reduce $S_m$ (compressing entropy) must correspondingly
increase $|\mathcal{F}(K_m)|$ (expanding configuration space), and vice versa.
\end{proposition}

This approximate conservation law is the memory
analogue of the Generative Sufficiency Principle:
the conserved quantity is not the content of the
memory but the volume of admissible futures it
supports. A memory transformation is admissibility-
preserving if and only if it preserves this volume,
regardless of what happens to the specific field
values or wave packet parameters.

\bigskip
\noindent\emph{What follows.} Part~III has now
established the full mathematical core: the four
primitive objects (with lamphron and lamphrodyne as
derived quantities), the categorical structure of
HYDRA, the Minimal Projection Theorem with the
Persistence Theorem as corollary, RSVP field dynamics
with the Persistence Balance result, stratified semantic
manifolds with the Marine Stability theorem, sheaf
semantics, and memory theory with the Residue Persistence
criterion. Chapter~15 adds one further result that
completes the mathematical inversion at the heart of the
program: the formal demonstration that history is
primitive and state is derived.

%% --------------------------------------------------------------
\chapter{History as Primitive}
%% --------------------------------------------------------------

The frameworks developed in the preceding chapters all
treat the current state as a derivative object---something
that can be recovered from continuation structure, field
residue, or admissibility equivalence classes. This chapter
formalizes the most radical version of that claim: in a
system with deterministic replay, the event sequence is a
complete sufficient statistic for system state. History is
not a record of past states; state is a function of history.
The inversion $\mathrm{State} \to \mathrm{History}$ is
not a philosophical slogan but a mathematical theorem.

\section{Event Sequences and Replay}

\begin{definition}[Event Sequence and Replay]
An event sequence over a state space $X$ is a finite
ordered list $E_n = (e_1, e_2, \ldots, e_n)$ of events
$e_i \in \mathcal{E}$. A replay function is a map
$R : \mathcal{E}^* \to X$ assigning to each event sequence
a system state. Replay is deterministic if $R$ is a
well-defined function (the same sequence always produces
the same state).
\end{definition}

In the Spherepop calculus, events are Pop, Refuse,
Collapse, Bind, and Repair operations; the event sequence
is the causal log of committed operations; and replay
is the execution of the log from an initial state. In
the AyeOS ontology (Chapter~20), the irreducible core
of the system is precisely this pair: the Spherepop causal
log and the RSVP field evolution loop that replays it.

\section{History as Sufficient Statistic}

\begin{theorem}[Replay Invariance]
\label{thm:replay_invariance}
If replay is deterministic, then the event sequence $E_n$
is a complete sufficient statistic for system state $X_n
= R(E_n)$: knowing $E_n$ is equivalent to knowing $X_n$,
and knowing $X_n$ is not in general sufficient to recover
$E_n$.
\end{theorem}

\begin{proof}
Sufficiency: by the determinism of $R$, $E_n$ determines
$X_n = R(E_n)$ uniquely. Hence $E_n$ is sufficient for
$X_n$.

Completeness: suppose $f(X_n)$ is any statistic that is
a function of $X_n$ alone. Then $f(X_n) = f(R(E_n))$,
which is a function of $E_n$, so $E_n$ subsumes $X_n$
as a statistic.

Non-reversibility: in general $R$ is not injective. Many
event sequences may produce the same state $X_n = R(E_n)
= R(E_n')$ for $E_n \neq E_n'$. Given only $X_n$, the
originating event sequence cannot be uniquely recovered.
Hence $X_n$ is not sufficient for $E_n$.
\end{proof}

\begin{remark}
Theorem~\ref{thm:replay_invariance} is the formal statement
that in a history-native system, state is derived and history
is primitive. The conventional engineering ontology assumes
that state is the fundamental object and history is a record
of past states. The Replay Invariance Theorem inverts this:
history (the event sequence) has strictly more information
than state (the replay output), because many histories can
produce the same state but the history cannot be recovered
from the state alone.
\end{remark}

\section{The History-Native Computation Principle}

\begin{principle}[History-Native Computation]
\label{prin:history_native}
In a system with deterministic replay, computation should
be formulated in terms of event sequences rather than
states. State is a derived object: $X_n = R(E_n)$.
History is the primitive: $E_n$ is the complete record
from which all derived objects can be recovered.
\end{principle}

The History-Native Computation Principle is not merely a
theoretical preference. It has architectural consequences.
A system that stores state loses the causal history from
which the state was derived; a system that stores event
sequences can always recover any past state by replay,
and can also recover causal structure (which events
contributed to which outcomes) that is invisible in the
state representation.

This connects directly to the Generative Compression
Hypothesis: the event sequence $E_n$ is the generator,
and the state space trajectory $\{X_0, X_1, \ldots, X_n\}$
is the inventory it generates. Preserving the event
sequence preserves strictly more than preserving the
state trajectory.

\section{Connections to Spherepop and AyeOS}

In the Spherepop event calculus, the five primitive
operations (Pop, Refuse, Collapse, Bind, Repair) are
the event alphabet $\mathcal{E}$. The causal log is the
event sequence $E_n$. The current bubble configuration
is the state $X_n = R(E_n)$.

The Replay Invariance Theorem implies that the causal
log is the irreducible object: it is sufficient for the
current state (the configuration can be reconstructed
by replaying the log) and it contains strictly more
information (the causal structure of which operations
produced which configurations). A system that stores
only the current bubble configuration has discarded the
causal history; a system that stores the causal log can
recover both the configuration and its causal structure.

In AyeOS terms (Chapter~20), this is why the Spherepop
causal log is identified as part of the irreducible core:
without it, the system cannot perform genuine Phoenix
reconstruction (which requires the causal history to
verify that a reconstruction counts as continuity), and
cannot support the proof-carrying dependent types that
ensure each operation maintains admissibility.

\section{The State-History Inversion and the Continuation
Principle}

The Replay Invariance Theorem and the Continuation
Principle (Principle~\ref{prin:continuation}) are dual
statements of the same fundamental inversion.

The Continuation Principle says: two states are
semantically equivalent if and only if they have the
same admissible future structure. States are defined
by their futures.

The Replay Invariance Theorem says: states are derived
from histories, and histories cannot be recovered from
states. States are determined by their pasts.

Together they close a loop: a state is defined by the
future continuations it makes available (Continuation
Principle) and by the event history that produced it
(Replay Invariance). A complete characterization of a
state requires both: its causal history (where it came
from) and its admissibility structure (where it can go).
The present moment is the intersection of a history and
a future, not a self-subsistent object.

This is the deepest version of the title \emph{Continuations
Before Objects}: objects---states---are not primitive.
They are the intersection of a generating history and a
generated future, both of which have priority over the
state itself.

\bigskip
\noindent\emph{What follows.} Part~III is now complete.
Part~IV applies the mathematical core to consequences:
restricted mind-change learning as external confirmation,
the geometry of intelligence, semantic fibers and agency
projection, projection collapse in social systems, and the
Structural Universality Conjecture as the framework's
largest open claim.




%% ==============================================================
\part{Consequences of the Reachability Thesis}
%% ==============================================================

\noindent\textit{Part~III established the mathematical
core: the four primitive objects and their derived
quantities (lamphron, lamphrodyne), the categorical
structure of HYDRA, the Minimal Projection Theorem,
RSVP field dynamics, stratified manifolds, sheaf
semantics, memory theory, and the Replay Invariance
Theorem. Part~IV asks whether this machinery explains
anything non-trivial. The answer this part gives is
yes, in five distinct domains: learning theory,
cognitive geometry, agency, social systems, and
physical computation. Each domain provides an
independent confirmation that admissible continuation
structure is the right primitive. The part closes
with the Reachability Equivalence Theorem, which
shows that intelligence, agency, memory, repair,
and preservation are not merely related by analogy
but are formally equivalent under admissibility
preservation.}

\bigskip

%% --------------------------------------------------------------
\chapter{Restricted Mind Changes and the Geometry
of Reachability}
%% --------------------------------------------------------------

The Balbach--Zeugmann restricted mind-change framework
was developed entirely independently of the Admissibility
Program, in the tradition of formal learning theory,
motivated by questions about the computational complexity
of teaching and learning. Its central move---equipping
the hypothesis space with a neighborhood relation and
constraining learners to move only between neighboring
hypotheses---turns out to be the same move that HYDRA
makes when it filters the reachability graph to retain
only admissibility-preserving transitions. This chapter
develops that identification precisely, proves two new
theorems connecting learning cost to path geometry, and
establishes the framework as an early mathematical
precursor to the Admissibility Program's central claims.

\section{The Balbach--Zeugmann Framework}

\begin{definition}[Restricted Mind-Change System]
A restricted mind-change system is a pair $(R, \nu)$
where $R$ is a hypothesis space and $\nu \subseteq
R \times R$ is a neighborhood relation on hypotheses,
encoding which transitions are admissible. A learner
in this system occupies a node $h \in R$ and may
transition only to neighbors $h' \in \nu(h)$.
\end{definition}

The neighborhood relation $\nu$ is the learning-theoretic
instantiation of the admissibility structure
$\mathcal{A}$ from Chapter~8. The hypothesis graph
$(R, \nu)$ is a trajectory space whose paths are
admissible sequences of mind changes. A learner is an
agent navigating this space; a teacher is a system
that modifies the space---by presenting data, eliminating
alternatives, or otherwise reshaping the landscape of
admissible transitions---so that the learner's trajectory
converges to a target.

\begin{theorem}[Hypothesis Graph as Admissibility Manifold]
The restricted mind-change system $(R, \nu)$ is an
admissibility system in the sense of Definition~8.3,
with $\mathcal{A}(h) = \nu(h)$ for each hypothesis $h$.
Learning is trajectory navigation through admissible
hypothesis space.
\end{theorem}

\begin{proof}
Set $\mathcal{X} = R$, $\mathcal{A}(h) = \nu(h)$, and
$\mu$ the counting measure on finite hypothesis spaces
(or an appropriate measure on infinite ones). The
persistence functional $P$ measures the volume of
admissible hypotheses remaining consistent with the
observed data. Each of the four primitive objects is
instantiated; the system is an admissibility system.
\end{proof}

\section{Reachability Cost and Path Geometry}

In the classical analysis of learning, the difficulty
of reaching a target hypothesis is measured by the
number of mind changes required. This measure is
graph-theoretic: it counts steps in the hypothesis
graph. The Admissibility Program suggests a finer
measure that takes account of the geometry of
admissible continuations along the path.

\begin{definition}[Learning Cost]
Let $(R, \nu)$ be a restricted mind-change system
with admissibility volumes $A(h) = \mathrm{Vol}
(\mathcal{A}(h))$ for each hypothesis $h$. The
learning cost of a path $\gamma = (h_0, h_1, \ldots,
h_n)$ from $h_0$ to $h_n$ is:
\[
C(\gamma) = \sum_{k=0}^{n-1} \frac{1}{A(h_k)},
\]
and the geodesic learning cost from $h_i$ to $h_j$
is:
\[
C(h_i, h_j) = \inf_\gamma C(\gamma)
= \inf_\gamma \int_\gamma \frac{1}{A(h)}\,ds,
\]
where the infimum is over admissible paths from $h_i$
to $h_j$ and $ds$ denotes the path-length element.
\end{definition}

The learning cost integrates the inverse of admissibility
volume along the path. A path through hypotheses with
large admissibility volumes is cheap: many continuations
are available at each step, so the path does not commit
the learner to a narrow trajectory. A path through
hypotheses with small admissibility volumes is expensive:
few continuations are available, the learner is
constrained, and errors are costly.

\begin{theorem}[Reachability Cost Theorem]
\label{thm:reachability_cost}
In a restricted mind-change system, the geodesic
learning cost $C(h_i, h_j)$ is determined by the
geometry of admissible continuations along the path,
not by graph distance alone. Specifically:
\begin{enumerate}
\item[(i)] Two paths of equal graph length may have
   arbitrarily different learning costs.
\item[(ii)] The geodesic learning path need not be
   the shortest path in the hypothesis graph.
\item[(iii)] The cost of reaching a target hypothesis
   depends on the reachability structure of the
   intermediate hypotheses, not on the destination.
\end{enumerate}
\end{theorem}

\begin{proof}
(i) Let $h_i = h_0 \to h_1 \to h_j$ and $h_i = h_0
\to h_1' \to h_j$ be two paths of length~2 with
$A(h_1) = \varepsilon$ and $A(h_1') = M$. Then
$C(\gamma) = 1/A(h_0) + 1/\varepsilon$ while
$C(\gamma') = 1/A(h_0) + 1/M$. As $\varepsilon \to 0$
and $M \to \infty$, the cost ratio $C(\gamma)/C(\gamma')
\to \infty$, establishing arbitrary difference.

(ii) The geodesic is determined by minimizing $C(\gamma)$,
not path length. A longer path through high-admissibility
hypotheses may have lower cost than a shorter path
through low-admissibility hypotheses.

(iii) The costs $1/A(h_k)$ depend on intermediate
hypotheses $h_k$, not on $h_j$. Two paths reaching
the same $h_j$ via different intermediates have costs
determined entirely by those intermediates.
\end{proof}

\begin{remark}
Theorem~\ref{thm:reachability_cost} is the formal
statement of the recurring informal claim: the cost
of reaching an idea depends on the path, not the
destination. In learning-theoretic terms: teaching
difficulty is not a property of the target concept
but of the neighborhood structure through which the
learner must pass to reach it. This connects directly
to the Eight-Letter Keyboard argument (Chapter~5):
the difficulty of reaching a symbol depends on the
motor path, not the symbol itself.
\end{remark}

\section{Teaching Without Feedback:
         Continuation Elimination}

Balbach and Zeugmann show that a teacher can often
force convergence without explicit specification of
the target, by progressively eliminating alternative
neighbors. This strategy has a precise
admissibility-theoretic formulation.

\begin{proposition}[Continuation Elimination Principle]
\label{prop:cont_elimination}
Teaching does not require constructing the target
hypothesis. It is sufficient to eliminate all admissible
alternatives:
\[
\bigcap_{i} H_i = \{h^*\}
\implies \text{convergence to } h^*,
\]
where $H_i$ is the set of hypotheses consistent with
the $i$-th piece of information provided by the teacher.
\end{proposition}

\begin{proof}
By hypothesis, $\bigcap_i H_i = \{h^*\}$. For any
learner constrained to the admissible graph $(R, \nu)$
and consistent with all presented information, the
only remaining admissible hypothesis is $h^*$. The
learner therefore converges to $h^*$ regardless of
the path taken through the graph, because $h^*$ is
the unique element of the intersection.
\end{proof}

\begin{remark}
Proposition~\ref{prop:cont_elimination} is the teaching
analogue of the Generative Sufficiency Principle: the
teacher does not need to specify the generator $G^*$
directly. It suffices to eliminate enough inadmissible
continuations that $G^*$ is the only remaining option.
This connects to CLIO projection (eliminating
representations that destroy future-relevant distinctions),
Spherepop collapse (removing bubble configurations that
close off continuation), and Repair Theory constraint
insertion (narrowing the admissibility relation until
the desired trajectory is forced). All four achieve
convergence by strategic elimination rather than
explicit construction.
\end{remark}

\section{HYDRA as Generalized Restricted Mind-Change
         Learning}

\begin{corollary}[HYDRA Generalizes Restricted Mind-Change
Learning]
The HYDRA architecture $H = \mathrm{GLU} \circ M \circ
T \circ F_a \circ G_a \circ R$ is a generalization of
restricted mind-change learning from hypothesis spaces
to full cognitive trajectory spaces.
\end{corollary}

\begin{proof}
In restricted mind-change learning: state space $=$ $R$,
admissibility $=$ $\nu$, learner $=$ agent navigating
$(R,\nu)$, teacher $=$ information that eliminates
alternatives. In HYDRA: state space $=$ $\mathcal{X}$,
admissibility $=$ $\mathcal{A}$, the six-functor chain
replaces the hypothesis graph with a richer stratified
manifold, and the RSVP field provides a metric on
learning cost. HYDRA instantiates the same categorical
structure as $(R, \nu)$ at a more general level.
Restricted mind-change learning is the special case
where $\mathcal{X}$ is discrete, $\mathcal{A} = \nu$,
and there is no RSVP field metric.
\end{proof}

%% --------------------------------------------------------------
\chapter{The Geometry of Intelligence}
%% --------------------------------------------------------------

Intelligence, within the framework developed here, is
not a capacity for accurate prediction, nor for utility
maximization, nor for any particular competence. It is
the capacity to maintain and expand admissible future
structure over time: to avoid shrinking the space of
reachable futures, and when possible to enlarge it.
This chapter formalizes that claim, proves two results
characterizing intelligent systems geometrically, and
connects the formal definitions back to the Continuation
Principle and the lamphron field.

\section{The Coherent Agent}

\begin{definition}[Coherent Agent]
A coherent agent is a system $(M, \Phi, v, S, \pi)$
satisfying:
\begin{enumerate}
\item[(i)] Admissible trajectory preservation: all
   transitions taken by the agent lie in $\mathcal{A}$;
\item[(ii)] Entropy regulation: $S(x(t), t) \leq
   S_\mathrm{max}$ for all $t$;
\item[(iii)] Local-to-global sheaf consistency: the
   agent's local semantic states are globally assemblable
   (no hallucination in the sense of
   Theorem~13.2);
\item[(iv)] Memory stabilization: past trajectories
   leave residue that stabilizes future navigation
   (MEM\textbar{}8 wave packets persist above threshold).
\end{enumerate}
\end{definition}

\begin{theorem}[Recursive Admissibility Stabilization]
A coherent agent starting from $x_0 \in \mathcal{A}_0$
and taking only $\mathcal{A}$-preserving transitions
satisfies $x_t \in \mathcal{A}_t$ for all $t \geq 0$.
\end{theorem}

\begin{proof}
By induction. Base case: $x_0 \in \mathcal{A}_0$ by
hypothesis. Inductive step: if $x_t \in \mathcal{A}_t$,
condition (i) gives $x_{t+1} \in \mathcal{A}(x_t)
\subseteq \mathcal{A}_{t+1}$.
\end{proof}

\section{The Admissibility Growth Criterion}

The recursive stabilization theorem shows that a
coherent agent remains in the admissible region. But
remaining in $\mathcal{A}$ is a minimum condition.
A genuinely intelligent agent does more: it maintains
or expands the volume of admissible futures available
to it.

\begin{definition}[Admissibility Volume Process]
For an agent with state $x(t)$ at time $t$, the
admissibility volume process is $A_t =
\mathrm{Vol}(\mathcal{A}(x(t)))$, equal to $e^{L(x(t))}$
where $L$ is the lamphron field.
\end{definition}

\begin{hypothesis}[Admissibility Growth Criterion]
\label{thm:intelligence_criterion}
Intelligent systems tend to preserve or increase their
admissibility volume under resource constraints: for
sufficiently long time horizons $\Delta > 0$,
\[
\mathbb{E}[A_{t+\Delta} \mid x(t)] \geq A_t,
\]
that is, the expected admissibility volume does not
decrease along the agent's trajectory. This is a
framework hypothesis ($\mathcal{H}_A$), not a
mathematical theorem derived from prior definitions.
It would be falsified by an intelligent system that
consistently reduces $A_t$ over long horizons, or by
a demonstrably non-intelligent system that consistently
maintains or increases $A_t$.
\end{hypothesis}

\begin{remark}
Hypothesis~$\mathcal{H}_A$ connects to the Generative
Sufficiency Principle (Principle~6.1): maintaining
$\mathbb{E}[A_{t+\Delta}] \geq A_t$ is operationally
equivalent to preserving the minimal generator $G^*$
of future admissibility. The expectation rather than
the pointwise condition is deliberate: genuinely
intelligent behavior may sometimes sacrifice local
admissibility volume to reach a richer global region,
committing to a narrow path in order to expand future
options. The criterion requires only that this trade
be non-negative in expectation. As an empirical
hypothesis, $\mathcal{H}_A$ predicts that systems
satisfying it should exhibit better transfer learning,
greater robustness under distributional shift, and
more stable long-horizon performance than systems
optimizing immediate prediction accuracy alone.
\end{remark}

\section{Intelligence as Curvature Control}

A stronger geometric characterization of intelligence
concerns the curvature of the admissibility manifold
rather than merely its volume.

\begin{definition}[Admissibility Curvature]
The admissibility curvature $K_\mathcal{A}(x)$ at a
point $x$ is the sectional curvature of the lamphron
field $L(x)$ in the direction of maximal descent:
\[
K_\mathcal{A}(x) = \min_{\|u\|=1} \nabla^2 L(x)[u,u].
\]
Negative curvature at $x$ indicates that the agent
is approaching a local lamphron minimum: a trap or
commitment point from which future admissibility
volume will be difficult to recover.
\end{definition}

\begin{proposition}[Curvature Bound for Intelligent Systems]
\label{prop:curvature_bound}
An intelligent system---one satisfying the Admissibility
Growth Criterion---maintains a lower bound on admissibility
curvature: $K_\mathcal{A}(x(t)) \geq -\kappa$ for
some $\kappa > 0$ along its trajectory.
\end{proposition}

\begin{proof}
If $K_\mathcal{A}(x) < -\kappa$ for arbitrarily large
$\kappa$, the lamphron field has unbounded negative
curvature at $x$, meaning the agent is approaching a
critical point of $L$ from which all directions lead
downhill. In a neighborhood of such a point,
$\mathbb{E}[L(x(t+\Delta))] < L(x(t))$ for all
$\Delta > 0$ small, violating the Admissibility Growth
Criterion. Hence an agent satisfying the criterion
cannot approach points of unbounded negative curvature.
\end{proof}

\begin{remark}
Proposition~\ref{prop:curvature_bound} formalizes why
intelligent systems avoid traps, local minima, and
brittle commitments: these are precisely the regions
of high negative admissibility curvature. An agent
that commits irreversibly---burning bridges, closing
off future options, concentrating on a single path
to the exclusion of all others---is an agent moving
toward negative curvature and therefore, by the
proposition, violating the intelligence criterion.
This is the geometric version of the practical wisdom
that robust reasoning avoids premature commitment.
\end{remark}

\section{The Geometry-Intelligence Correspondence}

The full correspondence between geometric objects and
cognitive concepts is:

\begin{center}
\renewcommand{\arraystretch}{1.4}
\begin{tabular}{ll}
\toprule
Geometric object & Cognitive concept \\
\midrule
$\Phi(x,t)$ & Semantic salience \\
$v(x,t)$ & Semantic tendency \\
$S(x,t)$ & Semantic ambiguity \\
$L(x) = \log\mathrm{Vol}(\mathcal{A}(x))$ &
  Future richness (lamphron) \\
$\mathcal{L} = \nabla L$ & Drive toward open futures
  (lamphrodyne) \\
$K_\mathcal{A}(x)$ & Commitment risk \\
Transport through $\mathcal{A}$ manifolds & Cognition \\
Stabilized field residue & Memory \\
Local-to-global sheaf compatibility & Reasoning \\
Tangent-constrained optimization & Learning \\
Recursive admissibility preservation & Intelligence \\
\bottomrule
\end{tabular}
\end{center}

%% --------------------------------------------------------------
\chapter{Semantic Fibers and Agency Projection}
%% --------------------------------------------------------------

The Minimal Projection Theorem established that every
admissibility-preserving projection $\pi : \mathcal{X}
\to M$ has fibers $\pi^{-1}(m)$ encoding what is lost
in the compression. This chapter identifies the specific
pathology that occurs when the fiber is forgotten---when
the image $M$ is mistaken for the full space
$\mathcal{X}$---and proves the Agency Collapse Theorem,
which shows this mistake is not merely an
epistemological error but a structural failure that
destroys the capacity for genuinely agentive behavior.

\section{The Realization Fibration}

\begin{definition}[Realization Map]
Let $\Theta$ be a parameter space (of agent configurations,
institutional positions, model weights, or similar)
and $N$ a space of observable behaviors or outputs.
The realization map $\mathbf{R} : \Theta \to N$ assigns
to each parameter configuration its observable behavior.
\end{definition}

\begin{theorem}[Semantic Fiber Theorem]
Under mild regularity---specifically, that $\eta \in N$
is a regular value of $\mathbf{R}$---the fiber
$\mathbf{R}^{-1}(\eta)$ is a smooth submanifold of
$\Theta$ of codimension $\dim N$.
\end{theorem}

\begin{proof}
By the Regular Level Set Theorem (established
mathematics): if $\eta$ is a regular value of
$\mathbf{R}$, then $d\mathbf{R}_\theta$ is surjective
for all $\theta \in \mathbf{R}^{-1}(\eta)$, and the
preimage is a smooth submanifold of codimension
$\dim N$.
\end{proof}

The fiber $\mathbf{R}^{-1}(\eta)$ is all configurations
that produce the same observable behavior $\eta$. It
encodes the causal history, contingent circumstances,
and structural variation that lead to the same external
output. In the social case: two individuals at the same
income level $m$ may have arrived there via radically
different trajectories, with different opportunity
structures, different constraints, and different future
potentials---all of which are encoded in the fiber
$\pi^{-1}(m)$ but invisible from the image $m$.

\section{Agency Projection and Its Pathology}

\begin{definition}[Agency Projection]
Agency projection is the identification $M \equiv
\mathcal{X}$: the replacement of the full trajectory
space by the image of a projection, collapsing the
fiber $\pi^{-1}(m)$ to a point. Under agency
projection, the causal history encoded in the fiber
becomes irrecoverable.
\end{definition}

\begin{theorem}[Agency Collapse Theorem]
\label{thm:agency_collapse}
Let $\pi : \mathcal{X} \to M$ be a projection from
the full trajectory space to a metric space. If
$\mathrm{Vol}(\pi^{-1}(m)) \to 0$, then the agent's
admissible future structure becomes indistinguishable
from the metric on $M$: the agent's agency collapses
to the metric.
\end{theorem}

\begin{proof}
The admissibility volume of the agent at position
$m \in M$ is $A(m) = \mathrm{Vol}(\mathcal{A}(m))$.
Since admissible continuations are paths in
$\mathcal{X}$ that pass through the fiber
$\pi^{-1}(m)$, we have $A(m) \leq C \cdot
\mathrm{Vol}(\pi^{-1}(m))$ for some constant $C$.
As $\mathrm{Vol}(\pi^{-1}(m)) \to 0$, we get $A(m)
\to 0$: the admissibility volume vanishes. By the
Persistence Theorem (Corollary~10.1), when
$\mathcal{A}(m_1) = \mathcal{A}(m_2) = \emptyset$,
no admissibility-preserving observer can distinguish
$m_1$ from $m_2$. The only remaining structure is the
metric on $M$ itself.
\end{proof}

\begin{remark}
Theorem~\ref{thm:agency_collapse} mathematically
captures the following family of conflations, which
are structurally identical despite their different
surface appearances:

\medskip
\begin{center}
\renewcommand{\arraystretch}{1.3}
\begin{tabular}{lll}
\toprule
Domain & Projection $\pi$ & Collapsed fiber becomes \\
\midrule
AI & Behavioral output & Genuine belief or understanding \\
Social media & Profile & Person \\
Education & GPA or credential & Student's capacity \\
Labor & Productivity metric & Worker \\
Intelligence & Benchmark score & Cognitive architecture \\
Credit & Credit score & Financial trajectory \\
\bottomrule
\end{tabular}
\end{center}
\medskip

In each case the fiber $\pi^{-1}(m)$---the full space
of configurations producing that metric value---is
collapsed to a point, and the metric is mistaken for
the agent. The Agency Collapse Theorem shows this is
not an oversight or a simplification but a mathematical
error: when the fiber volume vanishes, agency
genuinely collapses to the metric. The error does
not merely misrepresent the agent; in systems that
respond to the metric rather than the agent, it
creates the reality it describes.
\end{remark}

\section{The Ontological Compression Funnel}

The double projection $\mathcal{X} \xrightarrow{q} M
\xrightarrow{\mathbf{R}} N$ makes the structure of
information loss explicit. The first projection $q$
maps from the full trajectory space to the observable
manifold; the second projection $\mathbf{R}$ maps
from the observable manifold to the behavioral output
space.

\begin{definition}[Ontological Compression Funnel]
The ontological compression funnel is the composed
projection $\Pi_\mathrm{comp} = \mathbf{R} \circ q :
\mathcal{X} \to N$. The developmental projection
deficit (Chapter~20) of the funnel is:
\[
\Delta = H(\mathcal{X}) - H(N),
\]
the information destroyed by the double compression.
\end{definition}

Every step of the funnel discards information. The
causal history encoded in the fiber of $q$ is lost
at the first projection. The structural variation
encoded in the fiber of $\mathbf{R}$ is lost at the
second. What arrives at $N$ is a doubly compressed
image from which neither level of lost information
can be recovered without access to the original
trajectory space $\mathcal{X}$.

%% --------------------------------------------------------------
\chapter{Projection Collapse in Social Systems}
%% --------------------------------------------------------------

Social systems apply projections to persons, institutions,
and trajectories just as cognitive systems apply
projections to states and representations. The social
case is not an analogy to the mathematical framework;
it is an instantiation of it. The projection
$\pi : \mathcal{X} \to M$ that maps the full trajectory
space of an individual life to a low-dimensional
achievement manifold, the fiber collapse that occurs
when $M$ is mistaken for $\mathcal{X}$, the curvature
distortion that results from dominant coordinates---
these are the same mathematical objects as in the
preceding chapters, specialized to a social domain.
The chapter's position in Part~IV, after the mathematical
results rather than as a standalone application, is
deliberate: it demonstrates that the same theorem
applies in a domain that has no obvious connection to
field theory or cognitive architecture, providing
external evidence for the framework's generality.

\section{Meritocracy as Projection Collapse}

\begin{definition}[Social Trajectory Manifold]
The social trajectory space $\mathcal{X}_\mathrm{soc}$
is the space of individual life histories, equipped
with an admissibility structure $\mathcal{A}_\mathrm{soc}$
encoding which future trajectories are accessible from
each current position. The individual lamphron is
$L_i = \log\mathrm{Vol}(\mathcal{A}_\mathrm{soc}(x_i))$,
the log-volume of life trajectories accessible to
individual $i$.
\end{definition}

Modern meritocratic systems define a projection $\pi :
\mathcal{X}_\mathrm{soc} \to M_\mathrm{ach}$ from
the full trajectory space to a low-dimensional
achievement manifold $M_\mathrm{ach}$ with coordinates
such as educational credentials, income, occupational
prestige, and measured performance. Projection collapse
occurs when $M_\mathrm{ach}$ is mistaken for
$\mathcal{X}_\mathrm{soc}$: when the achievement
position is treated as fully characterizing the
individual rather than as a compressed image of their
trajectory.

\begin{theorem}[Meritocratic Compression]
\label{thm:meritocratic}
Meritocratic systems identify achievement position
$m$ with the fiber $F_m = \pi^{-1}(m)$, causing
information loss proportional to $\dim(F_m)$. This
is a quotient collapse in the sense of the Minimal
Projection Theorem (Theorem~10.1).
\end{theorem}

\begin{proof}
By Theorem~10.1, the projection $\pi$ is an
admissibility projection: it identifies trajectories
with the same achievement position. The fiber
$F_m = \pi^{-1}(m)$ contains all life trajectories
producing position $m$, which differ in causal
history, contingent circumstance, and future
potential. Identifying $m$ with $F_m$ collapses
this fiber to a point, losing $\dim(F_m)$ dimensions
of information. The lost information includes exactly
the fiber structure that the Minimal Projection Theorem
shows cannot be recovered from the projection alone.
\end{proof}

\section{Meritocratic Hubris as Fiber Collapse}

\begin{definition}[Meritocratic Hubris]
Meritocratic hubris is the belief, held by an
individual at position $m \in M_\mathrm{ach}$, that
$m$ is a complete causal explanation of their
position---that the contingency encoded in the fiber
$F_m$ is null or negligible.
\end{definition}

Meritocratic hubris is precisely fiber collapse
experienced from the inside: the individual at $m$
perceives their position as self-authored because
the fiber $F_m$, which encodes the family endowment,
historical circumstance, social capital, biological
variation, and stochastic perturbation that together
produced $m$, is invisible to an observer who knows
only $m$. The Semantic Fiber Theorem
(Theorem~18.1) establishes that the fiber is a smooth
manifold of positive dimension; its invisibility from
$m$ is not evidence of its nullity but of the
information loss in the projection.

\section{Coordinate Dominance and Riemannian Warping}

When one coordinate of the achievement manifold
acquires dominant weight in the social metric,
the Riemannian structure of $M_\mathrm{ach}$ warps:
geodesic distances become dominated by separation
along the dominant coordinate, and the manifold
effectively reduces to a line.

\begin{proposition}[Coordinate Dominance]
Let $M = (s_1, \ldots, s_n)$ be a social semantic
manifold with metric $g_{ij}$. If coordinate $s_k$
acquires weight $w_k \gg w_j$ for all $j \neq k$,
then geodesic distances satisfy $d_g(x,y) \approx
|s_k(x) - s_k(y)| \cdot w_k$: social navigation
is governed by a total ordering along $s_k$ rather
than by the fiber structure of actual trajectories.
\end{proposition}

\begin{proof}
The geodesic distance in a Riemannian manifold with
metric $g$ satisfies $d_g(x,y) = \inf_\gamma
\int_0^1 \sqrt{g_{ij}\dot\gamma^i\dot\gamma^j}\,dt$.
When $w_k \gg w_j$ for all $j$, the metric is
dominated by $w_k(ds_k)^2$, and the geodesic is
approximately the path that minimizes $|s_k(x)
- s_k(y)|$, giving the stated approximation.
\end{proof}

\section{Collective Admissibility and the Common Good}

\begin{definition}[Societal Lamphron and Collective
Admissibility]
The societal lamphron is the weighted sum of
individual lamphrons:
\[
L_\mathrm{soc} = \sum_i w_i L_i
= \sum_i w_i \log\mathrm{Vol}(\mathcal{A}_\mathrm{soc}
(x_i)).
\]
Collective admissibility is defined as
$\mathcal{A}_\mathrm{collective} =
\mathrm{colim}_i\,\mathcal{A}^{(i)}$, which assembles
compatible future sets along their overlap structure.
Collective admissibility fails when this colimit admits
no global section.
\end{definition}

\begin{theorem}[Collective Admissibility Theorem]
\label{thm:collective_admissibility}
Maximizing individual metric performance on
$M_\mathrm{ach}$ does not in general maximize societal
lamphron $L_\mathrm{soc}$.
\end{theorem}

\begin{proof}
By counterexample. Let $n = 2$ individuals with
$w_1 = w_2 = 1/2$ and achievement manifold $M =
\mathbb{R}$. Suppose individual~1 can increase their
metric position from $m_1$ to $m_1 + \delta$ by
taking an action that reduces individual~2's
admissibility volume by $\Delta A_2 \gg \delta \cdot
\partial A_1 / \partial m_1$. Then the increase in
$L_1 = \log A_1$ from improving metric position is
outweighed by the decrease in $L_2 = \log A_2$ from
the action's externality. Societal lamphron
$L_\mathrm{soc} = (L_1 + L_2)/2$ decreases even
though individual~1's metric position improves.
Concrete instantiation: competitive zero-sum strategies
that improve one individual's metric position by
reducing another's access to opportunity, training, or
networks decrease $L_\mathrm{soc}$ while increasing
individual metric performance.
\end{proof}

\begin{remark}
Theorem~\ref{thm:collective_admissibility} provides a
formal explanation of why meritocratic metric
optimization can reduce total societal future capacity.
The metric $M_\mathrm{ach}$ is a projection of the
trajectory space; optimizing the projection does not
optimize the trajectory space. Individual metric
improvements achieved by collapsing others' fibers
---by reducing the admissibility volume of others'
futures---may increase individual position on $M$
while decreasing aggregate lamphron on $\mathcal{X}$.
The social system that mistakes $M$ for $\mathcal{X}$
will misidentify this reduction as an improvement.
\end{remark}

%% --------------------------------------------------------------
\chapter{Clockless Computing as Continuation-First
         Architecture}
%% --------------------------------------------------------------

The previous chapters demonstrated that the
continuation-first ontology applies across learning
theory, cognitive geometry, agency, and social systems.
This chapter demonstrates something stronger: that it
can be physically realized in silicon. Null Convention
Logic (NCL), also known as clockless or asynchronous
logic, is a computing paradigm in which circuits
transition when and only when their inputs are complete.
There is no external clock forcing synchronization;
instead, the circuit's own completion structure governs
its evolution. This makes NCL not merely a hardware
curiosity but an existence proof: a physical system
in which continuations are genuinely primitive and
states are genuinely derived.

\section{The Clocked Paradigm and Its Hidden Projection}

A conventional clocked digital machine imposes a global
temporal projection at each clock tick:
\[
\pi_t : \mathcal{X} \to \mathcal{X}_t,
\]
forcing every subsystem to synchronize against a common
external temporal coordinate. The clock performs a
quotient operation: all physical states of the circuit
that appear identical at clock tick $t$ are identified
regardless of their mid-cycle dynamics. Many distinctions
in the underlying physical system are intentionally
discarded because they occur between clock edges.

The hidden assumption is that those discarded
distinctions do not matter for the computation. This
assumption is often justified, but its justification
is domain-specific and not guaranteed. The clock is
a convenience imposed on the computation, not a
property of the computation itself.

\section{Completion as Admissibility}

NCL rejects the clock-first assumption. Instead of
synchronizing on an external temporal coordinate,
NCL circuits synchronize on completion structure: a
gate transitions when its inputs are complete, not
when a clock says it should.

\begin{definition}[DATA, NULL, and Illegal States]
In NCL, signals carry three ontological statuses,
each with a precise admissibility-theoretic meaning:

\medskip
\begin{center}
\renewcommand{\arraystretch}{1.4}
\begin{tabular}{lll}
\toprule
NCL status & Admissibility meaning & Lamphron \\
\midrule
DATA & Realized continuation & $L > -\infty$ \\
NULL & Unrealized continuation & $L$ undefined \\
Illegal state & Contradictory continuation &
  $L = -\infty$ \\
Completion & Admissibility established & $dL/dt > 0$ \\
Acknowledge & Continuation confirmed & $L$ stable \\
Orphan & Unfinished continuation & $L$ stale \\
Deadlock & Zero admissible futures & $L \to -\infty$ \\
Clock & External temporal projection &
  $\pi_t : \mathcal{X} \to T$ \\
\bottomrule
\end{tabular}
\end{center}
\medskip

Crucially, NULL is \emph{not} a forbidden continuation
($\mathrm{Vol}(\mathcal{A}(v)) = 0$). It is a
\emph{continuation vacuum}: the absence of a committed
continuation, a node awaiting a DATA wavefront that
will establish one. A forbidden continuation is a
dead end; NULL is an uninstantiated region waiting for
admissibility to be established. The illegal state is
the genuinely inadmissible case: a node receiving
contradictory signals that cannot be resolved into any
admissible continuation.
\end{definition}

The alternation between DATA and NULL is therefore an
alternation between realized and unrealized continuation.
A computation in NCL is a succession of reachability
fields: DATA wavefronts propagate, establishing
continuations; NULL wavefronts follow, resetting nodes
to the vacuum state so the next continuation can be
established. The clock, on this reading, is exactly
what CLIO would call a projection:
$\pi_t : \mathcal{X} \to T$, collapsing the richer
partial order of completion events into a single linear
temporal coordinate. Clocked computing is
continuation-first computing with an additional
projection imposed on top.

\section{Continuation Synchronization}

\begin{theorem}[Continuation Synchronization]
\label{thm:cont_sync}
Let $G$ be an NCL network with completion condition
$C(v)$ at each node $v$. Then the network evolution
is determined entirely by the partial order induced
by completion relations $u \prec v$ (meaning $v$
cannot transition until $u$ has completed), rather
than by any external temporal parameter.
\end{theorem}

\begin{proof}
By the NCL completeness criterion, node $v$ transitions
if and only if $C(v)$ is satisfied, which requires all
inputs to $v$ in the dependency graph to have
transitioned. This defines a partial order $\prec$ on
transitions: $u \prec v$ iff $v$ depends on $u$.
The execution order of any NCL network is a linear
extension of this partial order. The partial order is
determined entirely by the circuit topology and the
completion conditions, with no reference to any
external temporal parameter $t$. Hence the evolution
is independent of any global clock.
\end{proof}

\begin{remark}
Theorem~\ref{thm:cont_sync} establishes that the clock
is not a primitive of computation but a convenience
imposed on top of a more fundamental completion
structure. In NCL, the underlying ontology is explicit:
the primitive objects are completion events, and the
``time'' at which they occur is derived from their
causal order, not presupposed.
\end{remark}

\section{Orphan Continuations}

In conventional NCL analysis, ``orphan paths'' are
timing nuisances: signal paths that have not completed
before the local admissibility structure changes.
The geometric framework gives a precise name and
definition to this phenomenon.

\begin{definition}[Orphan Continuation]
A continuation $\gamma_o$ is orphaned if it remains
active---if $\mathrm{Vol}(\mathcal{A}(\gamma_o)) > 0$
---after the admissibility structure that generated
it has been withdrawn by a NULL wavefront.
\end{definition}

An orphan continuation is a trajectory that has not
yet completed its data wavefront propagation but whose
generating admissibility structure has already been
reset. The result is a stale DATA collision: the
orphan's signal reaches a node that is already in
NULL phase, creating a glitch. In the geometric
language, this is a continuation collision: a
trajectory $\gamma_o$ collides with a later trajectory
$\gamma'$ that has reset the admissibility structure
$\gamma_o$ was navigating.

\section{The Clock Elimination Theorem}

\begin{theorem}[Clock Elimination]
\label{thm:clock_elimination}
Let $(\mathcal{X}, \mathcal{A})$ be a computational
system whose transitions are governed solely by
admissibility-preserving completion relations. Then
there exists an equivalent realization whose evolution
is independent of any global temporal coordinate:
\[
\text{Completion} \implies \text{Time},
\qquad
\text{Time} \not\implies \text{Completion}.
\]
\end{theorem}

\begin{proof}
Given the admissibility system $(\mathcal{X},
\mathcal{A})$ governed by completion relations, the
Continuation Synchronization Theorem (Theorem~20.1)
provides a realization in which execution order is a
linear extension of the partial order $\prec$. This
realization assigns a ``time'' to each transition:
$t(v) = |\{u : u \prec v\}|$, the depth in the
partial order. This time is internal, derived from
completion structure, not presupposed.

Conversely, a global temporal coordinate imposes a
projection $\pi_t : \mathcal{X} \to \mathcal{X}_t$
that identifies all states within a clock cycle. By
the Minimal Projection Theorem (Theorem~10.1), this
projection destroys information whenever two states
that are clock-cycle-equivalent have different
admissibility structures. Hence global time implies
a projection, but a projection does not imply any
particular global time: many temporal orderings are
consistent with the same completion partial order.
\end{proof}

\begin{remark}
Theorem~\ref{thm:clock_elimination} establishes that
the clock is not a primitive of computation but a
convenience imposed on top of a more fundamental
completion structure. In NCL, the underlying ontology
is explicit: the primitive objects are completion
events, and the ``time'' at which they occur is
derived from their causal order, not presupposed.
\end{remark}

\section{The Completion Sufficiency Theorem}

The relationship between completion events and global
states in NCL mirrors the relationship between event
sequences and system states established in the Replay
Invariance Theorem (Theorem~15.1). The following
result is its hardware manifestation.

\begin{theorem}[Completion Sufficiency]
\label{thm:completion_sufficiency}
Let $E$ be the set of completion events in a clockless
system and $\mathcal{S}$ the set of intermediate global
states. Then:
\begin{enumerate}
\item[(i)] $E \Rightarrow \mathcal{S}$: the completion
   sequence determines the computation and suffices to
   reconstruct all intermediate states.
\item[(ii)] $\mathcal{S} \not\Rightarrow E$: a
   collection of sampled global states does not in
   general uniquely determine the completion history
   that produced them.
\end{enumerate}
\end{theorem}

\begin{proof}
(i) By the Continuation Synchronization Theorem
(Theorem~\ref{thm:cont_sync}), the network evolution
is a linear extension of the completion partial order
$\prec$. Given the completion event sequence $E$ and
the circuit topology, the state of every node at every
point in the computation is determined: node $v$ is in
DATA state after its completion event $e_v \in E$ and
in NULL state before it. Hence $E$ determines
$\mathcal{S}$.

(ii) Multiple completion orderings---multiple linear
extensions of the partial order $\prec$---may produce
the same global state snapshots if the snapshotting
granularity is coarser than the event granularity.
Specifically, if two completion events $e_u$ and $e_v$
are incomparable in $\prec$ (neither $u \prec v$ nor
$v \prec u$), they may be completed in either order
without affecting any state snapshot that observes
both nodes after both completions. Hence the snapshot
sequence does not distinguish the two orderings.
\end{proof}

\begin{remark}
Theorem~\ref{thm:completion_sufficiency} is the
hardware manifestation of the Replay Invariance
Theorem. Both establish the same asymmetry:
\[
\text{events} \to \text{states}, \qquad
\text{states} \not\to \text{events}.
\]
In software systems (Chapter~15), events are Spherepop
operations and states are bubble configurations. In
hardware systems (this chapter), events are NCL
completion wavefronts and states are global circuit
snapshots. The asymmetry is the same in both cases
because it is not a feature of any particular
implementation but of the general relationship between
generative histories and their derived products. This
connects Replay Invariance, History as Primitive, and
Clockless Computing as three instances of the same
formal fact.
\end{remark}

The Clock Elimination Theorem connects to the Replay
Invariance Theorem (Theorem~15.1) in an unexpected
way. The Replay Invariance Theorem showed that in
systems with deterministic replay, the event sequence
is a complete sufficient statistic for system state:
history is primitive, state is derived.

In NCL, this becomes physically visible. The circuit
is not executing a sequence of stored states. It is
executing a sequence of completion events: DATA
wavefront at $v_1$, DATA wavefront at $v_2$, NULL
wavefront at $v_1$, and so on. The ``state'' of the
circuit at any moment is the temporary residue of
completed and uncompleted wavefronts, not a stored
object. The event sequence is primary; the state is
derived.

This aligns the NCL ontology precisely with the
History-Native Computation Principle (Principle~15.1):
computation should be formulated in terms of event
sequences rather than states. NCL is a physical
architecture that implements this principle at the
hardware level, making the philosophical commitment
of the Admissibility Program visible in silicon.

\section{HYDRA, Balbach--Zeugmann, and NCL as Instances
         of One Structure}

The three frameworks surveyed in Chapters~16, 17, and
this chapter are instances of the same geometric
structure:

\begin{center}
\renewcommand{\arraystretch}{1.4}
\begin{tabular}{lll}
\toprule
Framework & Primitive object & Derived object \\
\midrule
Balbach--Zeugmann & Admissible mind change &
  Hypothesis \\
HYDRA & Admissible trajectory & State \\
NCL & Completion event & Clock tick / State \\
\midrule
\multicolumn{3}{l}{In all three:} \\
& Admissible continuation & Object \\
\bottomrule
\end{tabular}
\end{center}

The pattern is uniform: the primitive object is an
admissible continuation, and what appears to be a
state---a hypothesis, a cognitive state, a circuit
state---is a derived object, the residue left by
a sequence of completed continuations. NCL is
therefore not merely a hardware curiosity. It is
an existence proof that a continuation-first ontology
can be realized physically.

%% --------------------------------------------------------------
\chapter{The Structural Universality Conjecture}
%% --------------------------------------------------------------

The results of Chapters~16 through~20 have demonstrated
that the admissibility framework applies across learning
theory, cognitive geometry, agency, social systems, and
physical computation. The Structural Universality
Conjecture claims that this is not coincidence: the
RSVP transformation is a natural transformation between
categories, functorial with respect to domain change,
that carries any domain admitting a stability measure
and an admissibility relation into the admissibility
framework. The conjecture organizes the open problems
of Chapter~23 and defines what proof of the program's
unity would require.

\section{Statement of the Conjecture}

\begin{conjecture}[Structural Universality]
\label{conj:universality}
Let $\mathbf{Adm}$ be the category of admissibility
systems $(\mathcal{X}, \Pi, \mathcal{A}, P)$ with
admissibility-preserving functors as morphisms, and
let $\mathbf{Rep}$ be the category of representational
systems \emph{with nontrivial transition dynamics}. The
RSVP transformation is a natural transformation:
\[
T : \mathbf{Rep} \Rightarrow \mathbf{Adm},
\]
functorial with respect to domain change.
\end{conjecture}

\begin{remark}
Conjecture~\ref{conj:universality} applies only to
systems possessing nontrivial transition dynamics---
systems for which admissible continuations are
well-defined. Static objects (files, photographs,
fixed databases) are not in scope: their admissible
continuations are determined by external agents, not
by the objects themselves. This restriction is not a
weakening; it is the correct scope. The conjecture
is about \emph{dynamical} systems. It is explicitly
unproved, falsifiable (a dynamical domain in which no
admissibility-preserving reformulation exists would
refute it), and defines the research frontier rather
than the established results. Proving it requires
precise categorical definitions of $\mathbf{Rep}$
and $\mathbf{Adm}$, construction of $T$ on morphisms,
and naturality verification across all target domains.
These are the open problems collected in Chapter~25.
\end{remark}

\section{The Six Target Domains}

\begin{center}
\renewcommand{\arraystretch}{1.4}
\begin{tabular}{lll}
\toprule
Domain & $\mathcal{X}$ & What persists \\
\midrule
RSVP cosmology & Plenum field space &
  Metastable admissibility basins \\
HYDRA cognition & Cognitive trajectory space &
  Admissible inference chains \\
MEM\textbar{}8 memory & Wave state space &
  Resonant field residue \\
Semantic Infrastructure & Module dependency graphs &
  Homotopy-compatible merges \\
Restricted mind-change & Hypothesis graph &
  Admissible learning paths \\
Political economy & Agent strategy space &
  Collective admissibility \\
\bottomrule
\end{tabular}
\end{center}

\section{The Restricted Form}

\begin{theorem}[Structural Universality, Restricted Form]
\label{thm:universality_restricted}
Let $D$ be any domain with state space $\mathcal{X}$,
admissible transitions $\mathcal{A}(x)$, and
differentiable stability measure $\Phi : \mathcal{X}
\to \mathbb{R}$. Taking $S = \log|\mathcal{A}|$ and
$v = \nabla\Phi$, the RSVP description is well-defined
and captures the persistence structure: long-term
behavior is governed by metastable basins of $\Phi$
under the flow of $v$ in the region $\{S < S_\mathrm{thresh}\}$.
\end{theorem}

\begin{proof}
Under the stated regularity, $S$ and $v$ are
well-defined smooth fields. The lamphron field is
$L = -S$ (since high $S$ means large admissibility
volume, but in RSVP convention $S$ decreases with
admissibility). In the quasi-static regime, the RSVP
equations reduce to the gradient flow $\dot{x} =
\nabla\Phi(x)$ in the admissible region $\{S <
S_\mathrm{thresh}\}$. Long-term behavior is
characterized by the stable fixed points of
$\nabla\Phi$ in that region, which are the metastable
basins.
\end{proof}

\section{Representation Independence of Lamphron}

A necessary condition for the Structural Universality
Conjecture is that the central quantity of the
framework---lamphron---is preserved under admissibility-
preserving domain changes.

\begin{theorem}[Representation Independence]
\label{thm:rep_independence}
Let $(\mathcal{X}, \mathcal{A})$ and $(\mathcal{Y},
\mathcal{B})$ be admissibility systems with lamphron
fields $L_X$ and $L_Y$ respectively. If there exists
an admissibility-preserving equivalence
$F : \mathcal{X} \to \mathcal{Y}$ (a bijection
satisfying $F(\mathcal{A}(x)) = \mathcal{B}(F(x))$
for all $x$), then:
\[
L_X(x) = L_Y(F(x)) + c
\]
for some constant $c$ depending only on the choice
of measure.
\end{theorem}

\begin{proof}
By the admissibility-preserving property of $F$:
$\mathcal{A}(x) = F^{-1}(\mathcal{B}(F(x)))$.
Hence $\mathrm{Vol}(\mathcal{A}(x)) =
\mathrm{Vol}(F^{-1}(\mathcal{B}(F(x)))) =
J_F \cdot \mathrm{Vol}(\mathcal{B}(F(x)))$,
where $J_F$ is the Jacobian of $F$ with respect to
the measures on $\mathcal{X}$ and $\mathcal{Y}$.
Taking logarithms: $L_X(x) = L_Y(F(x)) + \log J_F$.
When $F$ is measure-preserving ($J_F = 1$), the
constant $c = 0$ and lamphron is exactly preserved.
In general, $c = \log J_F$ is a constant normalizing
the two measures.
\end{proof}

\begin{remark}
Theorem~\ref{thm:rep_independence} establishes that
lamphron is representation-invariant: it does not
depend on the particular representation chosen for
the admissibility system, only on the admissibility
structure itself (up to measure normalization). This
is the concrete object that survives domain
translation, making lamphron the natural candidate
for the invariant quantity of the Structural
Universality Conjecture.
\end{remark}

%% --------------------------------------------------------------
\chapter{The Reachability Equivalence Theorem}
%% --------------------------------------------------------------

The preceding chapters of Part~IV have demonstrated
that the admissibility framework applies across five
independent domains. This final chapter of Part~IV
shows that the various quantities and properties
invoked across those chapters---admissibility volume,
lamphron, generative structure, agency, semantic
continuity---are not merely related by analogy.
They are formally equivalent under admissibility
preservation. The Reachability Equivalence Theorem
is the central result of Part~IV in the same sense
that the Minimal Projection Theorem is the central
result of Part~III: it unifies the preceding chapters
into a single statement.

\begin{theorem}[Reachability Equivalence Theorem]
\label{thm:reachability_equiv}
For any admissibility system $(\mathcal{X}, \mathcal{A})$,
the following conditions are equivalent:
\begin{enumerate}
\item[(i)] Preservation of admissible future volume:
   $A_{t+\Delta} \geq A_t$ in expectation.
\item[(ii)] Preservation of lamphron:
   $L(x(t+\Delta)) \geq L(x(t))$ in expectation.
\item[(iii)] Preservation of generative structure:
   the minimal generator $G^*$ is maintained
   (Principle~\ref{prin:sufficiency}).
\item[(iv)] Preservation of agency: the agent's
   admissible future structure is non-collapsing
   ($\mathrm{Vol}(\pi^{-1}(m)) \not\to 0$).
\item[(v)] Preservation of semantic continuity:
   the semantic space $M^*$ is maintained
   under the agent's trajectory
   (Corollary~\ref{cor:persistence}).
\end{enumerate}
\end{theorem}

\begin{proof}
(i) $\Leftrightarrow$ (ii): $A_t =
\mathrm{Vol}(\mathcal{A}(x(t))) = e^{L(x(t))}$. Since
$\exp$ is monotone, $A_{t+\Delta} \geq A_t \iff
L(x(t+\Delta)) \geq L(x(t))$.

(ii) $\Leftrightarrow$ (iii): By the Generative
Sufficiency Principle (Principle~\ref{prin:sufficiency}),
$G^*$ is maintained iff $R(G^*) \supseteq R(I)$, iff
$\mathrm{Vol}(\mathcal{A}(x)) \geq \mathrm{Vol}
(\mathcal{A}_I)$ where $\mathcal{A}_I$ is the
admissibility structure of the full inventory. Since
$L = \log\mathrm{Vol}(\mathcal{A})$, this is
equivalent to (ii).

(iii) $\Leftrightarrow$ (iv): By the Agency Collapse
Theorem (Theorem~\ref{thm:agency_collapse}), agency
collapses iff $\mathrm{Vol}(\pi^{-1}(m)) \to 0$, iff
$\mathrm{Vol}(\mathcal{A}(m)) \to 0$, iff $A_t \to 0$.
Since (iii) maintains $A_t \geq \mathrm{Vol}
(\mathcal{A}_{G^*}) > 0$, agency is preserved.

(iv) $\Leftrightarrow$ (v): By the Persistence Theorem
(Corollary~\ref{cor:persistence}), semantic continuity
is maintained iff $\mathcal{A}(x(t))$ is approximately
constant along the trajectory---the agent's admissible
future structure does not change in a way that would
make it semantically indistinguishable from a different
trajectory. This is equivalent to preservation of the
semantic space $M^* = \mathcal{X}/{\sim_\mathcal{A}}$
under the agent's motion, which holds iff the agent's
trajectory stays within a single equivalence class of
$\sim_\mathcal{A}$, which is guaranteed by (iv).
\end{proof}

\begin{remark}
Theorem~\ref{thm:reachability_equiv} is the formal
statement that intelligence, agency, memory, repair,
and preservation are all manifestations of the same
underlying condition: maintaining admissible future
volume. Prediction serves preservation. Optimization
serves preservation. Memory serves preservation.
Repair serves preservation. Learning serves
preservation. Everything is downstream of
continuation.
\end{remark}

\begin{corollary}[Generator Preservation Criterion]
\label{cor:generator_preservation}
An implementation realizes a theory if and only if
it preserves the generators of admissible continuation
structure---not all state variables, not all data
structures, only the generators.
\end{corollary}

\begin{proof}
By (iii) of Theorem~\ref{thm:reachability_equiv},
preservation of generative structure is equivalent
to all other forms of preservation. An implementation
that preserves $G^*$ satisfies condition (iii) and
hence all of (i)--(v). An implementation that
preserves non-generating state variables but loses
$G^*$ fails condition (iii) and hence fails all
other preservation conditions.
\end{proof}

\begin{remark}
Corollary~\ref{cor:generator_preservation} is the
engineering consequence of the Reachability Equivalence
Theorem. It ties the engineering audit of Chapter~25
directly back to the Generative Compression Hypothesis
of Chapter~6: the question to ask of any implementation
is not ``does it preserve all state variables?'' but
``does it preserve the generators of admissible
continuation structure?'' Any implementation that
does so is a valid realization of the theory, by
Theorem~9.2 (Architecture-Independence). Any
implementation that does not is inadequate regardless
of how many state variables it faithfully reproduces.
\end{remark}

\bigskip
\noindent\emph{What follows.} Part~V examines whether
the theory survives implementation in the 8b.IS
engineering stack, using the Development Triple
$(R, S, F)$ and the realization matrix to audit
what exists, what is specified, and what is
admissibly reachable. Part~VI asks what has actually
been proved. Part~VII asks what follows if the
framework is correct.





%% ==============================================================
%% NEW CHAPTER: Empirical Consequences, Failure Modes, and Scope
%% Inserted between Part IV and Part V
%% ==============================================================

\chapter{Empirical Consequences, Failure Modes, and
         Scope of the Admissibility Program}
\label{ch:falsification}

A reviewer of this monograph would be correct to observe that many of
its central results are structural rather than empirical, and that
the framework's core identification---intelligence as admissibility
preservation---risks being a definition rather than a discovery. This
chapter addresses that concern directly by doing three things: stating
what the framework predicts, specifying what would falsify it, and
clarifying the scope of its claims.

\section{The Epistemic Status of Central Claims}

The framework makes claims at four distinct levels, and conflating
them is the primary source of confusion about what the monograph
establishes.

The Minimal Projection Theorem (Theorem~10.1) is established
mathematics: given an admissibility relation, a minimal
admissibility-preserving projection exists and is unique up to
isomorphism. This is not a hypothesis about intelligence or reality.

The Generative Compression Hypothesis (Hypothesis~6.1) and the
Admissibility Growth Criterion ($\mathcal{H}_A$, Chapter~17) are
empirical hypotheses that use the mathematical vocabulary to make
testable claims about generative systems and intelligent agents. They
would be falsified by systems that contradict their predictions.

The RSVP field equations and the sheaf-theoretic semantics are
\emph{framework specifications} with open phenomenological commitments:
the source terms $\rho$ and $\sigma$ are not derived from first
principles, and the sheaf existence theorem is conditional on an
unproved regularity result. These are placeholders for domain-specific
derivations, analogous to viscosity coefficients in fluid mechanics
or reaction rates in chemistry. They should be read as specifications
of what a complete theory in each domain would look like, not as
established results about those domains.

The Structural Universality Conjecture (Conjecture~21.1) is an
organizing conjecture: explicitly marked as unproved, falsifiable in
principle, and defining the program's research frontier rather than
its established achievements.

\section{Scope: Systems with Transition Dynamics}

The Structural Universality Conjecture applies to systems possessing
nontrivial transition structure---systems in which the notion of
admissible continuation is well-defined. This excludes static objects
entirely. A JPEG file has no intrinsic admissibility dynamics; its
admissible continuations are determined by the viewer or processor,
not the file itself. The conjecture is not about static objects. It
is about \emph{dynamical} systems: systems that evolve, that face
choices between possible next states, and for which some transitions
are admissible and others are not.

This restriction is not a weakening of the claim; it is its correct
scope. The Admissibility Program is a theory of persistent dynamical
systems---of agents, of cognitive architectures, of physical
computation, of social institutions. It does not claim to subsume
inert objects, which require no theory of admissible continuation.

\section{A Minimal Concrete Realization of HYDRA}

HYDRA is an architectural specification, not a unique implementation.
The Architecture-Independence Theorem (Theorem~9.2) establishes that
any system realizing the functor $H$ up to natural isomorphism is a
valid implementation. The following gives a minimal concrete instance
to demonstrate that the specification is satisfiable, not merely
formal.

\begin{definition}[Minimal HYDRA Realization]
A minimal realization of HYDRA assigns:
\begin{itemize}
\item $R(x) = E(x)$: a learned embedding of the input;
\item $G_a(E) = (V, \nu)$: the $k$-nearest-neighbor graph on $E$,
  where edges encode admissible transitions;
\item $F_a(v) = \mathbf{1}[L(v) > \varepsilon]$: thresholded
  lamphron filter retaining nodes with admissibility volume above
  $\varepsilon$;
\item $T(\mathcal{G}_a)$: random walk trajectories through the
  filtered graph;
\item $M(\Gamma)$: residue accumulation by path traversal frequency;
\item $\mathrm{GLU}(\Gamma_M) = \arg\max_{\gamma} \int_\gamma
  \Phi\,dt$: trajectory selection by maximum accumulated potential.
\end{itemize}
\end{definition}

This realization is deliberately minimal. It makes no claim that the
$k$-NN graph is the correct admissibility structure for any given
domain, or that random walk trajectories are the correct continuation
model. It demonstrates that the specification is satisfiable with
standard components, and that the choice of components is what
varies across implementations while the functor composition is what
is preserved.

\section{Empirical Predictions}

The framework generates the following testable predictions, each
following from the Admissibility Growth Criterion
($\mathcal{H}_A$) or the Generative Compression Hypothesis
(Hypothesis~6.1).

\textit{Prediction 1: Transfer learning advantage.} Systems that
explicitly preserve continuation structure---maintaining admissibility
volume across task boundaries---should exhibit better transfer
learning than systems optimizing immediate prediction accuracy,
controlling for model size and training data. The mechanism: a
generator-preserving system retains the structure needed to
reconstruct performance in new domains; an inventory-optimizing system
does not.

\textit{Prediction 2: Generator compression ratio.} Memory
architectures that store residue fields (generators) rather than
explicit state inventories should require less storage to maintain
equivalent associative capacity. The mechanism: the Implicit
Association Density proposition (Proposition~14.0) establishes that
$O(n)$ wave packet storage supports $O(n^2)$ implicit associations;
a static index supporting $O(n^2)$ explicit associations requires
$O(n^2)$ storage.

\textit{Prediction 3: Long-horizon reward divergence.} Agents
maximizing immediate reward (prediction-primary) should produce
lower long-horizon admissibility volume than agents maximizing
expected continuation volume ($\mathcal{H}_A$), particularly in
environments with deceptive short-term optima. The mechanism: local
reward optima are often admissibility traps---high immediate reward,
low future reachability.

\textit{Prediction 4: Collective admissibility and institutional
metrics.} Societies or organizations that optimize narrow achievement
metrics should exhibit lower collective admissibility (societal
lamphron $L_\mathrm{soc}$) over time than those that optimize
opportunity diversity. The mechanism: the Collective Admissibility
Theorem (Theorem~19.3) establishes that metric optimization need not
maximize $L_\mathrm{soc}$.

\textit{Prediction 5: Clockless energy efficiency.} Continuation-
first computing architectures (NCL and derivatives) should exhibit
higher energy efficiency per preserved continuation than clocked
architectures, because they avoid the thermodynamic cost of global
synchronization over idle cycles.

None of these predictions is proved here. Each requires experimental
design, measurement methodology, and empirical testing. They are
stated to demonstrate that the framework is not epistemically sealed:
it generates predictions that could, in principle, be found false.

\section{Failure Modes of the Admissibility Program}

The following conditions would falsify central claims of the
framework.

\textit{F1: Generator failure.} If there exist generative systems
where preserving $G^*$ does not preserve $R(G^*)$---where the minimal
generator is insufficient to reconstruct the relevant admissible
futures---the Generative Sufficiency Principle (Principle~6.1) fails.
This would require the reconstruction map $\rho : R(G^*) \to I$ to
be non-surjective.

\textit{F2: Inventory dominance.} If inventory preservation
consistently outperforms generator preservation in tasks requiring
transfer, robustness, or long-horizon planning, the Generative
Compression Hypothesis (Hypothesis~6.1) is empirically disconfirmed.
The hypothesis predicts the opposite.

\textit{F3: Continuation irrelevance.} If admissibility volume has
no measurable relationship to transfer learning performance,
representational robustness, memory efficiency, or agency stability,
the Admissibility Growth Criterion ($\mathcal{H}_A$) is disconfirmed.
This is the sharpest empirical test.

\textit{F4: Reachability non-uniqueness.} If systems with identical
admissibility structures (by the semantic equivalence definition)
exhibit radically different semantic behavior---if the quotient
$\mathcal{X}/{\sim_\mathcal{A}}$ fails to organize semantics---the
Continuation Principle as a theory of meaning is falsified.

\textit{F5: Agency independence of fiber collapse.} If agency
measures do not collapse when fiber volumes approach zero---if systems
that have undergone projection collapse exhibit genuine agentive
behavior---the Agency Collapse Theorem (Theorem~18.3) is falsified
as an empirical claim about agency.

\section{On the NCL Claim}

The Clockless Computing chapter (Chapter~20) argues that NCL is a
physically realized continuation-first architecture. This claim
requires precision. NCL circuits are implemented by physical state
variables (voltages); the claim is not that NCL abolishes states at
the physical level. The claim is that \emph{computational correctness}
in NCL depends on completion relations rather than on global clock
time. This is a weaker and more defensible assertion: that at the
computational level of description, completion events are primary
and clock ticks are an external projection imposed on a richer
partial order. The Completion Sufficiency Theorem (Theorem~20.3)
establishes this precisely: completion events determine global states,
but sampled global states do not determine completion events. NCL is
offered as existence proof that continuation-first description is
not merely philosophical but realizable in silicon---not as a claim
that physical states have been abolished.

\section{On the Hallucination Result}

The identification of hallucination as a cohomological obstruction
(Theorem~13.2) is a structural result, not an algorithm. It
characterizes the mathematical nature of hallucination---an output
that appears locally coherent but cannot be assembled into a global
section of the semantic sheaf---without providing a computable
procedure for detecting it in any specific system. The Computable
Projection Cohomology open problem (Open Problem~25.3) explicitly
acknowledges this gap. The sheaf result should be read as a
theoretical diagnosis that would become a practical tool if the open
problem were resolved and the sheaf $\mathcal{F}$ were specified for
a concrete architecture.


%% ==============================================================
\part{Engineering Realization as Ontology Audit}
%% ==============================================================

\noindent\textit{Parts~I through~IV established the
philosophical argument, the mathematical core, and the
consequences of the Reachability Thesis across five
independent domains. Part~V asks whether the theory
survives contact with implementation. The 8b.IS
engineering stack---Marine, MEM\textbar{}8, Phoenix
Protocol, AyeOS---is examined not as a catalog of
software components but as an empirical test case for
the Generative Compression Hypothesis. The organizing
tool is the Development Triple $(R, S, F)$, which
distinguishes implemented reality from formal
specification from admissible future. The realization
matrix then audits each system against the primitive
objects of the mathematical core. The chapter closes
with the Generator Preservation Criterion, which
establishes what any valid implementation must
preserve and what it may freely discard.}

\bigskip

%% --------------------------------------------------------------
\chapter{Engineering Realization as Ontology Audit}
%% --------------------------------------------------------------

Every engineering system instantiates mathematical
objects at multiple levels simultaneously. Marine,
as an engineering artifact, is a Rust function that
accepts or rejects signal vectors. Marine, as a
mathematical object, is a stability manifold
$M_\theta$ in signal space (Theorem~12.3). Marine,
as a theoretical claim, is the formal statement that
admissibility can be operationalized through
trajectory stability. These are different objects
occupying different levels of description. Confusing
them---treating the Rust function as the mathematical
object, or the mathematical object as the theoretical
claim---produces exactly the category error that
Chapter~18 identified as agency projection: mistaking
a projection for its domain.

The Engineering Realization as Ontology Audit chapter
refuses that confusion systematically. Its tool is
the Development Triple.

\section{The Development Triple}

\begin{definition}[Development Triple]
A computational system is represented by the triple
$\mathcal{D} = (R, S, F)$ where $R$ is the implemented
realization (what currently runs), $S$ is the formal
specification (what the implementation is supposed to
realize), and $F$ is the admissible future envelope
(what the system can become while remaining
admissibility-preserving). The three components are
distinct objects; confusing them is the developmental
projection error.
\end{definition}

The Development Triple is a specialization of the four
primitive objects $({\mathcal{X}}, \Pi, \mathcal{A}, P)$
to the domain of software engineering. The trajectory
space $\mathcal{X}$ is the space of possible
implementations; the projection system $\Pi$ includes
the projections $R = \pi_R(\mathcal{X})$ (implemented),
$S = \pi_S(\mathcal{X})$ (specified), and
$F = \pi_F(\mathcal{X})$ (admissibly future); the
admissibility structure $\mathcal{A}$ encodes which
implementation transitions are valid; and the
persistence functional $P$ measures how much of the
theory's continuation structure survives each
projection.

The Marine gate illustrates all three levels cleanly.
In $R$: a Rust function accepting signals where
$A(x,t) \geq \theta$. In $S$: the stability manifold
$M_\theta$ of Theorem~12.3 with its proof of robustness
under $C^2$-small perturbations. In $F$: extensions
to non-stationary admissibility, learned thresholds,
and multi-modal signal integration. These three
descriptions are not competing; they are projections
of the same mathematical object at different levels of
crystallization.

By contrast, orbital angular momentum routing for
AyeOS exists only in $S$ and $F$: $R_\mathrm{OAM}
= \varnothing$, $S_\mathrm{OAM} \neq \varnothing$,
$F_\mathrm{OAM} \neq \varnothing$. This is not a
failure; it is a normal developmental state. The
question is not whether $R = S$, but whether the
admissible path from $R$ to $S$ is non-empty.

\section{The Realization Ladder and Its Deficits}

\begin{definition}[Realization Ladder and Deficit]
A realization ladder is a chain of projections:
\[
T \xrightarrow{\alpha} S \xrightarrow{\beta} I
  \xrightarrow{\delta} D,
\]
where $T$ is the formal theory, $S$ the specification,
$I$ the implementation, and $D$ the deployed system.
The realization deficit is:
\[
\Delta_R = H(T) - H(D),
\]
where $H$ measures the admissibility-relevant
distinctions preserved at each level.
\end{definition}

\begin{theorem}[Implementation Projection Theorem]
\label{thm:impl_projection}
Every implementation is a projection of its theory.
Consequently, if $\Delta_R > 0$, the deployed system
cannot be identified with the theory it realizes.
\end{theorem}

\begin{proof}
Each map in the realization ladder discards
distinctions: unimplemented definitions ($\alpha$),
simplified algorithms ($\beta$), absent hardware
and runtime constraints ($\delta$). The composition
$T \to D$ is therefore a projection by the Minimal
Projection Theorem. If $\Delta_R > 0$, at least one
admissibility-relevant distinction in $T$ is not
present in $D$.
\end{proof}

\begin{theorem}[Developmental Projection Theorem]
\label{thm:dev_projection}
For any nontrivial system, $H(\mathcal{D}) > H(R)$:
the Development Triple contains strictly more
admissibility-relevant information than the
implemented realization alone.
\end{theorem}

\begin{proof}
Any non-trivial specification $S$ contains
distinctions not yet realized in $R$: future design
decisions, conditional behaviors, intended but
unimplemented functionality. Any non-trivial future
envelope $F$ contains admissible continuations not
yet instantiated. Hence $H(R,S,F) > H(R)$ whenever
$S \neq \pi_S(R)$ or $F \neq \pi_F(R)$.
\end{proof}

\begin{remark}
Theorem~\ref{thm:dev_projection} is the formal
statement of the observation that a repository is not
a system. A repository is one projection of a larger
developmental object. The source code reveals only
the currently crystallized region of the developmental
manifold; the specification reveals additional
reachable structure; the future envelope reveals
directions of admissible extension. Treating the
repository as the system is agency projection applied
to software: mistaking the metric for the agent.
\end{remark}

\section{Crystallization and Potential}

\begin{definition}[Crystallized and Potential Components]
For a system $\mathcal{D} = (R, S, F)$, let $C$ denote
the crystallized component (what is implemented and
running) and $P$ the potential component (what is
specified or admissibly future but not yet
implemented):
\[
\mathcal{D} = C \oplus P.
\]
Engineering is the progressive crystallization of
admissible potential: $P \to C$ over developmental
time, subject to the admissibility constraint that
each crystallization step preserves the generator
of future continuation structure.
\end{definition}

The goal of development is not $\Delta_R \to 0$. Some
specifications are exploratory; some reveal better
trajectories; some are deliberately held in potential
form to preserve developmental flexibility. The goal
is that each crystallization step satisfies the
Generator Preservation Criterion: the implemented
component preserves the generator of the theory's
admissible continuation structure, not necessarily
every state variable or data structure.

\section{The Realization Matrix}

The realization matrix audits the 8b.IS stack by
asking, for each system, what its primitive state is,
how it evolves, what survives refactoring, how it
gates inadmissible trajectories, and what level of
the realization ladder it currently occupies.

Two concrete realization choices in the MEM\textbar{}8
implementation deserve explicit treatment as instances
of the general principles developed in Part~III.

The first is spatial layout. The TARTAN Approximation
Theorem (Theorem~12.2) establishes that semantic
manifolds can be decomposed into tiles within which
the RSVP field is approximately uniform. In the
hardware implementation, this translates directly to
a storage locality requirement: related wave packets
must be physically co-located on storage to ensure
that resonance queries hit adjacent sectors without
complex routing. The implementation uses a Hilbert
space-filling curve to map the three-dimensional
semantic field onto the one-dimensional address space
of NVMe flash storage. The Hilbert curve is not an
arbitrary engineering choice; it is the one-dimensional
projection that best preserves three-dimensional
proximity, instantiating the TARTAN principle that
semantic proximity should mirror physical proximity.
This is the Generative Compression Hypothesis applied
to hardware: the generator of spatial locality (the
Hilbert ordering) preserves the continuation structure
(wave resonance) at far lower cost than would explicit
routing tables.

The second is the computation substrate. The RSVP
field equations are continuous PDEs; any implementation
necessarily discretizes them. The hardware
implementation uses an integer-arithmetic physics
engine on configurable logic, avoiding floating-point
operations entirely. This is a deliberate realization
deficit: the $\mathcal{D}$ to $R$ projection discards
the continuous real-valued dynamics in favor of
cache-friendly integer approximations that run on
standard silicon at speeds sufficient to maintain
wave coherence across the field. The Generator
Preservation Criterion (Theorem~23.2) justifies this:
the implementation need not preserve continuous field
values; it need only preserve $R(G_I) \cong R(G_T)$,
the reachable continuation structure. Integer
approximations that maintain wave resonance patterns
satisfy this criterion even though they do not
reproduce the exact PDE dynamics.

\begin{center}
\renewcommand{\arraystretch}{1.5}
\small
\begin{tabular}{>{\bfseries}lll}
\toprule
 & Marine & MEM\textbar{}8 \\
\midrule
Primitive State &
  Signal $x(t)$ &
  Wave packet $(A,\omega,\phi,D,I)$ \\
Dynamics &
  $J_\tau$ + threshold &
  Wave eq.\ + viscosity \\
Persistence &
  $M_\theta$ stability &
  Residue field $\Phi_m$ \\
Admissibility &
  $A(x,t) \geq \theta$ &
  $\rho > \mu\Phi$ on $\Omega_m$ \\
$R$ status & Prototype & Partial (Rust) \\
$S$ status & Complete & Partial \\
$F$ status & Extensible & Extensible \\
\bottomrule
\end{tabular}

\medskip

\begin{tabular}{>{\bfseries}lll}
\toprule
 & Phoenix & AyeOS \\
\midrule
Primitive State &
  Coherence measure &
  Continuation topology \\
Dynamics &
  Lyapunov check at 0.73\,Hz &
  Sheaf routing \\
Persistence &
  Causal log + generator &
  Functor coherence \\
Admissibility &
  $\mathcal{A}(x_\mathrm{after}) \cong
  \mathcal{A}(x_\mathrm{before})$ &
  Sheaf compatibility \\
$R$ status & Prototype & Design \\
$S$ status & Partial & Sketch \\
$F$ status & Distributeable & Spherepop-native \\
\bottomrule
\end{tabular}
\end{center}

The matrix reveals the distribution of crystallization
across the stack: Marine is the most fully crystallized
(prototype in $R$, complete specification in $S$);
AyeOS is the least (design level in $R$, sketch in
$S$). This is not a ranking by value; it is an audit
of developmental state. AyeOS's relative incompletion
in $R$ does not make it less important --- its future
envelope $F$ is the richest of the four systems because
it is the orchestration layer that the others depend on.

\section{The Stack as a Single Admissibility System}

\begin{definition}[Stack Admissibility System]
The 8b.IS stack is an admissibility system
$\mathcal{S}_\mathrm{stack} = (\mathcal{X}_\mathrm{stack},
\Pi_\mathrm{stack}, \mathcal{A}_\mathrm{stack},
P_\mathrm{stack})$ where the four layers are admissibility
systems $\mathcal{S}_i = (X_i, \nu_i, \mathcal{A}_i, P_i)$
connected by interfaces:
\[
\mathcal{S}_\mathrm{Marine} \xrightarrow{\pi_1}
\mathcal{S}_\mathrm{MEM8} \xrightarrow{\pi_2}
\mathcal{S}_\mathrm{Phoenix} \xrightarrow{\pi_3}
\mathcal{S}_\mathrm{AyeOS}.
\]
\end{definition}

\begin{definition}[Admissibility-Preserving Interface]\label{def:adm_interface}
An interface $\pi_i : X_i \to X_{i+1}$ is
admissibility-preserving if:
\[
(x \to_{\nu_i} y) \implies (\pi_i(x) \to_{\nu_{i+1}} \pi_i(y)).
\]
\end{definition}

\begin{theorem}[Stack Coherence]
\label{thm:stack_coherence}
If every interface $\pi_i$ is admissibility-preserving,
then the composite:
\[
\Pi = \pi_3 \circ \pi_2 \circ \pi_1 :
\mathcal{S}_\mathrm{Marine} \to \mathcal{S}_\mathrm{AyeOS}
\]
is admissibility-preserving. The four-layer stack is
not four systems composed; it is one admissibility-
preserving map.
\end{theorem}

\begin{proof}
By induction on the number of interfaces. Two
admissibility-preserving interfaces compose to an
admissibility-preserving composite: if $x \to_{\nu_1}
y$ implies $\pi_1(x) \to_{\nu_2} \pi_1(y)$, and
$u \to_{\nu_2} v$ implies $\pi_2(u) \to_{\nu_3}
\pi_2(v)$, then $x \to_{\nu_1} y$ implies
$\pi_2(\pi_1(x)) \to_{\nu_3} \pi_2(\pi_1(y))$. Extend
to three interfaces by the same argument.
\end{proof}

\begin{remark}
The physical realization of Theorem~\ref{thm:stack_coherence}
in the 8b.IS hardware is an optical full-mesh
interconnect between nodes. The Stack Coherence
Theorem requires that admissibility be preserved
across layer interfaces; the optical mesh provides
the bandwidth---aggregating to hundreds of gigabits
per second across an $n$-node cluster---necessary
to maintain wave coherence across the entire field
without centralized clock synchronization. The
0.73\,Hz cognitive heartbeat of the Phoenix Protocol
is the system-level admissibility check that
periodically verifies coherence is maintained across
the full mesh, playing the role of the persistence
functional $P$ from Definition~8.4. The topology is
intentionally clockless at the inter-node level,
instantiating the Clock Elimination Theorem
(Theorem~20.2): global wave coherence is maintained
by the completion partial order of resonance events,
not by an external temporal coordinate.
\end{remark}

\section{The Generator Preservation Criterion}

\begin{theorem}[Generator Preservation Criterion]
\label{thm:gen_preservation}
An implementation $I$ realizes a theory $T$ if and only
if:
\[
R(G_I) \cong R(G_T)
\]
up to admissibility-preserving equivalence, where $G_T$
is the minimal generator of the theory and $G_I$ is the
minimal generator preserved by the implementation.
\end{theorem}

\begin{proof}
If $R(G_I) \cong R(G_T)$, then every admissible
continuation generated by the theory has an
implementation-level counterpart and every
implementation-level continuation is theoretically
admissible. The implementation realizes the theory's
continuation structure, which by the Reachability
Equivalence Theorem (Theorem~22.1) is equivalent to
realizing its agency, semantic continuity, lamphron
field, and generative structure.

Conversely, if $I$ realizes $T$, it preserves the
admissible futures constituting the theory's operational
content. Since admissible futures are generated by
$G_T$, the implementation must preserve a generator
$G_I$ with $R(G_I) \cong R(G_T)$ by the Generative
Sufficiency Principle (Principle~6.1).
\end{proof}

\begin{corollary}[State Non-Essentiality]
An implementation may discard state variables, data
structures, and representational details without
ceasing to realize the theory, provided it preserves
the generator $G_I$ with $R(G_I) \cong R(G_T)$.
\end{corollary}

\begin{remark}
Corollary~23.1 is the formal justification for the
observation that the 8b.IS stack's engineering
history---implementations changed, state representations
multiplied, abstractions drifted---does not constitute
a failure of the theory, provided the compositional
functor chain $R \to G_a \to F_a \to T \to M \to
\mathrm{GLU}$ is preserved. The modules are derived
objects. The compositional structure is the generator.
This is the Generative Compression Hypothesis made
operational: what must survive is not the inventory
but the generator.
\end{remark}

%% ==============================================================
\part{Critical Assessment}
%% ==============================================================

\noindent\textit{Part~VI applies the statement
classification discipline introduced in the Preface
to the entire monograph. Every central claim is
assigned to one of four classes: established
mathematics, framework definition, derived theorem,
or open conjecture. The open problems are then ordered
by their dependency structure, and the Admissibility
Program is assessed on its own criterion for progress:
not the accumulation of analogies but the conversion
of conjectures to theorems or falsified claims.}

\bigskip

%% --------------------------------------------------------------
\chapter{What the Framework Has Proved and What It
Has Conjectured}
%% --------------------------------------------------------------

A framework that cannot distinguish its theorems from
its conjectures is not yet a theory. The statement
classification applied throughout this monograph---
established mathematics, framework definition, derived
theorem, open conjecture---is not an editorial
convention. It is the framework's criterion for its
own epistemic integrity.

\section{The Classification System}

\begin{principle}[Classification Discipline]
\label{prin:classification}
Every central claim in the Admissibility Program
belongs to exactly one class:
\begin{itemize}
\item[$\mathcal{M}$:] Established mathematics---results
  from the mathematical literature that hold
  independently of the framework.
\item[$\mathcal{D}$:] Framework definitions---
  vocabulary and conceptual commitments that introduce
  the framework's specific content.
\item[$\mathcal{T}$:] Derived theorems---results that
  follow from established mathematics once the
  framework definitions are in place.
\item[$\mathcal{Q}$:] Open conjectures---claims that
  are falsifiable in principle but not yet proved.
\end{itemize}
\end{principle}

\begin{proposition}[No-Theorem-from-Definition Warning]
If a result follows only because a framework definition
was chosen to make it true, the result belongs to
$\mathcal{D}$, not to $\mathcal{M}$ or $\mathcal{T}$.
\end{proposition}

\begin{proof}
A definition introduces vocabulary. Consequences that
follow from the vocabulary alone---such as ``semantic
equivalence is an equivalence relation'' when semantic
equivalence is defined as equality of admissibility
sets---are grammatical facts about the definition,
not independent mathematical results. They may be
correct and useful, but they do not constitute evidence
that the world has the structure the definition
describes.
\end{proof}

\section{Classification of Central Results}

The following table classifies the principal results
of the monograph. Results marked $\mathcal{M}$ are
standard; $\mathcal{D}$ are definitional; $\mathcal{T}$
are derived; $\mathcal{Q}$ are conjectural.

\begin{center}
\renewcommand{\arraystretch}{1.4}
\begin{tabular}{>{\itshape}lc}
\toprule
Result & Class \\
\midrule
Semantic equivalence is an equivalence relation &
  $\mathcal{M}$ \\
Universal property of quotient spaces & $\mathcal{M}$ \\
Whitney stratification theorem & $\mathcal{M}$ \\
Convergence of gradient descent on compact strata &
  $\mathcal{M}$ \\
Gr\"onwall inequality & $\mathcal{M}$ \\
Regular level set theorem & $\mathcal{M}$ \\
Lyapunov stability theorem (standard form) &
  $\mathcal{M}$ \\
\midrule
Continuation Principle & $\mathcal{D}$ \\
Lamphron field & $\mathcal{D}$ \\
Admissibility structure & $\mathcal{D}$ \\
Development Triple & $\mathcal{D}$ \\
Coherent agent & $\mathcal{D}$ \\
\midrule
Minimal Projection Theorem & $\mathcal{T}$ \\
Persistence Theorem & $\mathcal{T}$ \\
Functoriality of $H$ & $\mathcal{T}$ \\
Architecture-Independence Theorem & $\mathcal{T}$ \\
TARTAN Approximation Theorem & $\mathcal{T}$ \\
Marine Stability Theorem & $\mathcal{T}$ \\
Memory Persistence Bound & $\mathcal{T}$ \\
Residue Persistence Theorem & $\mathcal{T}$ \\
Replay Invariance Theorem & $\mathcal{T}$ \\
Continuation Synchronization Theorem & $\mathcal{T}$ \\
Clock Elimination Theorem & $\mathcal{T}$ \\
Completion Sufficiency Theorem & $\mathcal{T}$ \\
Reachability Cost Theorem & $\mathcal{T}$ \\
Admissibility Growth Criterion & $\mathcal{T}$ \\
Agency Collapse Theorem & $\mathcal{T}$ \\
Collective Admissibility Theorem & $\mathcal{T}$ \\
Representation Independence Theorem & $\mathcal{T}$ \\
Reachability Equivalence Theorem & $\mathcal{T}$ \\
Generator Preservation Criterion & $\mathcal{T}$ \\
Stack Coherence Theorem & $\mathcal{T}$ \\
\midrule
Sheaf Existence for RSVP & $\mathcal{Q}$ \\
Hallucination as Cohomological Obstruction & $\mathcal{Q}$ \\
Structural Universality Conjecture & $\mathcal{Q}$ \\
RSVP Global Regularity & $\mathcal{Q}$ \\
RSVP Unique Continuation & $\mathcal{Q}$ \\
Source Term Derivation (all domains) & $\mathcal{Q}$ \\
Lean~4 Verification of Core Theorems & $\mathcal{Q}$ \\
Admissibility Curvature Bound (general) & $\mathcal{Q}$ \\
Computable Projection Cohomology & $\mathcal{Q}$ \\
\bottomrule
\end{tabular}
\end{center}

The classification reveals the framework's actual
epistemic state: a substantial body of derived theorems
resting on clear definitions and established mathematics,
with a cluster of open conjectures---centered on RSVP
regularity and structural universality---defining the
frontier. The conjectures are not defects; they are
the program's research agenda stated precisely enough
to be falsifiable.

%% --------------------------------------------------------------
\chapter{Open Problems}
%% --------------------------------------------------------------

\section{Dependency Ordering}

Open problems are ordered by dependency: a problem
listed earlier unlocks more downstream results.

\begin{definition}[Unlocking Priority]
The priority of open problem $Q_i$ is
$U(Q_i) = |\{Q_j : Q_i \leadsto Q_j\}|$, the number
of downstream problems whose resolution depends on
$Q_i$.
\end{definition}

\begin{proposition}[Regularity Priority]
RSVP global regularity and unique continuation have
maximal unlocking priority among current open problems.
\end{proposition}

\begin{proof}
Unique continuation unlocks sheaf existence
(Theorem~13.1 becomes unconditional). Sheaf existence
unlocks hallucination obstruction (Theorem~13.2 becomes
an established result rather than a research direction).
Hallucination obstruction unlocks computable projection
cohomology. Regularity underlies all three. No other
open problem unlocks as many downstream results.
\end{proof}

\section{First Priority: RSVP Regularity}

\begin{conjecture}[RSVP Local Well-Posedness]
For sufficiently regular initial data $U_0 \in H^s$
with $s$ large enough, the RSVP system
$\partial_t U = \mathcal{F}(U, \nabla U, \Delta U) + Q$
admits a unique local solution $U(t) \in C([0,T], H^s)$.
\end{conjecture}

\begin{conjecture}[RSVP Unique Continuation]
If two RSVP solutions agree on a non-empty open set
and satisfy the same source terms, they agree on the
connected admissibility domain generated by that open
set.
\end{conjecture}

Without unique continuation, the sheaf existence
theorem (Theorem~13.1) and all results that depend on
it remain conditional on an unproved regularity
assumption.

\section{Second Priority: Source Term Derivation}

The RSVP field equations leave the source terms
$\rho(v, S)$ and $\sigma(\Phi, v)$ as domain-specific
placeholders. Valid source terms must satisfy
dimensional consistency, preserve admissibility under
the relevant projection, and reproduce empirically
observed behavior.

\begin{openproblem}[Domain Source Derivation]
Derive valid source terms for cognition (where $\Phi$
should correspond to neural accessibility measures),
memory (where $S$ should correspond to forgetting
rates), political economy (where $\rho$ should
correspond to institutional reinforcement dynamics),
and cosmological field dynamics.
\end{openproblem}

Until source terms are derived, the framework makes
no quantitative empirical predictions in any specific
domain.

\section{Third Priority: Projection Cohomology}

\begin{openproblem}[Computable Projection Cohomology]
Develop algorithms for computing $\check{H}^1
(\mathcal{U}, \mathcal{F})$ in finite approximations
of HYDRA-like systems, converting hallucination
obstruction from a formal characterization into a
computable diagnostic.
\end{openproblem}

\section{Fourth Priority: Categorical Formalization}

\begin{openproblem}[Structural Universality Proof]
Prove Conjecture~21.1 by: (a) giving precise categorical
definitions of $\mathbf{Rep}$ and $\mathbf{Adm}$;
(b) constructing the natural transformation $T$ on
morphisms; (c) verifying naturality in each of the six
target domains.
\end{openproblem}

\section{Fifth Priority: Formal Verification}

\begin{openproblem}[Lean~4 Core Verification]
Complete the Lean~4 proof of the Minimal Projection
Theorem (Theorem~10.1) and the Replay Invariance
Theorem (Theorem~15.1) without \texttt{sorry}, and
extend to the Reachability Equivalence Theorem
(Theorem~22.1).
\end{openproblem}

\section{New Open Problems from Part IV}

Part IV generated several open problems not in the
original list:

\begin{openproblem}[Admissibility Curvature Bounds]
Determine necessary and sufficient conditions on
$(\mathcal{X}, \mathcal{A})$ for the curvature bound
$K_\mathcal{A} \geq -\kappa$ to hold globally along
intelligent trajectories (Proposition~17.3).
\end{openproblem}

\begin{openproblem}[Completion Semantics]
Develop a full formal semantics for asynchronous
computation based on completion partial orders,
connecting the Continuation Synchronization Theorem
(Theorem~20.1) to process algebra and to the
Spherepop event calculus.
\end{openproblem}

\begin{openproblem}[Societal Lamphron Measurement]
Develop empirical proxies for individual lamphron
$L_i$ and societal lamphron $L_\mathrm{soc}$ in
real social systems, allowing the Collective
Admissibility Theorem (Theorem~19.3) to make
testable predictions about the effects of meritocratic
compression on collective future capacity.
\end{openproblem}

%% --------------------------------------------------------------
\chapter{The Admissibility Program: Situation and
Direction}
%% --------------------------------------------------------------

\section{The Program's Actual State}

The Admissibility Program is a coherent theoretical
framework with a genuine mathematical core, a
significant body of derived theorems, a cluster of
identified open problems, and a clear criterion for
what progress looks like. That is an honest description
of where it stands.

What it is not is a completed theory. The RSVP field
equations lack proved regularity. The Structural
Universality Conjecture is unproved. No domain has
derived source terms. No implementation has been fully
audited against the Generator Preservation Criterion.

What it is: a framework in which the central objects
are precisely defined, the central results are proved
or identified as conjectural, and the frontier is
specific enough to be worked on.

\begin{definition}[Programmatic Progress]
A step constitutes progress if it moves at least one
claim from $\mathcal{Q}$ to $\mathcal{T}$, from
$\mathcal{Q}$ to falsification, or from $\mathcal{D}$
to empirically constrained definition.
\end{definition}

\begin{proposition}[Non-Expansion Criterion]
Adding new analogies without increasing proof,
falsification, or empirical constraint does not
constitute programmatic progress.
\end{proposition}

This criterion protects the framework from becoming
merely metaphorical. The history of theoretical
frameworks that accumulate analogies without advancing
their proof status is not an encouraging one. The
Admissibility Program has, at each stage of its
development, converted informal claims into formal
definitions and formal definitions into proved
theorems. The open problems named above are the
next stage of that conversion.

\section{Relationship to the Broader Corpus}

The monograph synthesizes results from a research
program that includes: \textit{Frozen Processes}
(process-primary ontology), \textit{Belief Geometry
and Reachability} (MSP/HMM/admissibility correspondence),
\textit{Holonomic Space} (RSVP applications as
textbook), \textit{Hidden Curvature} (semantic manifold
reconstruction), \textit{Semantic Infrastructure}
(synthesis monograph), \textit{Compressed Causality}
(counterfactual reconstruction), \textit{The Refractive
Self} (identity as dynamical invariant), \textit{The
Geometry of Closure} (ontological enlargement), and
\textit{Proxy Permanence Failure} (carbon governance
and RSVP-TARTAN-CLIO).

Each of these documents is a specialization of the
quartet $(\mathcal{X}, \Pi, \mathcal{A}, P)$ in a
particular domain, and each has generated results
that have been absorbed into the mathematical core
presented here. The relationship is not one of
derivation from a single axiomatic foundation; it
is one of convergent elaboration around a common
set of questions. The Generative Compression
Hypothesis and the Reachability Equivalence Theorem
are the current best formulations of what those
questions have in common.

%% ==============================================================
\part{Philosophical Synthesis}
%% ==============================================================

\noindent\textit{Part~VII does not introduce new
mathematics. It draws the philosophical consequences
of the mathematical results established in Parts~III
through~VI. The Reachability Thesis is stated as a
principle with a formal proof. Objects are shown to
be continuation quotients. The classical
ontology---objects first, relations second, futures
third---is inverted. Prediction is rehabilitated
as a derived operation on admissibility space. The
monograph closes by returning to Chapter~1 and showing
that the replacement of prediction by reachability
is not a refinement but an ontological inversion with
consequences at every level of the framework.}

\bigskip

%% --------------------------------------------------------------
\chapter{The Reachability Thesis}
%% --------------------------------------------------------------

The preceding twenty-six chapters have established a
mathematical framework, demonstrated its applicability
across six independent domains, audited its engineering
realization, and classified its claims by epistemic
status. This final chapter draws the philosophical
consequences. It does not require new proofs; all the
machinery is in place. What is required is to state
clearly what the machinery implies.

\section{The Final Equivalence}

The Reachability Equivalence Theorem (Theorem~22.1)
established that in any admissibility system, the
following are equivalent: preservation of admissible
future volume, preservation of lamphron, preservation
of generative structure, preservation of agency,
and preservation of semantic continuity. The theorem
is mathematical. Its philosophical consequence is
this: intelligence, agency, memory, repair, and
preservation are not related by analogy. They are
formally equivalent manifestations of the same
underlying condition---maintaining admissible future
volume---and what distinguishes them is not their
nature but their domain of application.

\section{Objects as Continuation Quotients}

\begin{theorem}[Objects as Continuation Quotients]
\label{thm:objects_quotients}
Objects, in the sense of the Admissibility Program,
are equivalence classes of trajectories under
continuation equivalence:
\[
\mathrm{Obj}(x) = [x]_{\sim_\mathcal{A}},
\]
where $x \sim_\mathcal{A} y \iff \mathcal{A}(x) =
\mathcal{A}(y)$.
\end{theorem}

\begin{proof}
By the Minimal Projection Theorem (Theorem~10.1), the
quotient $\mathcal{X}/{\sim_\mathcal{A}}$ is the unique
minimal admissibility-preserving projection of
$\mathcal{X}$. Any object recognized by the framework
must correspond to a point in this quotient---an
equivalence class of trajectories that share the same
admissible future structure. If two trajectories have
identical admissible futures, the Persistence Theorem
(Corollary~10.1) establishes that no
admissibility-preserving observer can distinguish them:
they are the same object. If two trajectories have
different admissible futures, the distinction is
preserved by the minimal projection: they are different
objects.
\end{proof}

\begin{remark}
Theorem~\ref{thm:objects_quotients} is the formal
statement of the title: \emph{Continuations Before
Objects}. Objects are not primitive; they are derived
from continuation structure. A chair is not an object
that happens to have future uses; it is an equivalence
class of material configurations that share the same
admissible future interactions. A person is not an
object that happens to have a biography; they are an
equivalence class of trajectories---a fiber in the
social trajectory space---that share the same
admissible future life paths. Destroying the fiber
while preserving the position on the achievement
manifold is not a lossless compression; it is a
category error.
\end{remark}

\section{The Minimal Admissible Self}

The Objects as Continuation Quotients theorem has a
specific application to identity that connects the
engineering audit of Part~V to the philosophical
synthesis of this chapter.

\begin{principle}[Minimal Admissible Self]
\label{prin:minimal_self}
The minimal admissible self $S^*$ is:
\[
S^* = \arg\min_S |S| \quad \text{subject to} \quad
R(S^*) = R(I_\mathrm{self}),
\]
where $I_\mathrm{self}$ is the inventory of an agent's
states and $R(S^*)$ is the admissible continuation
space generated by $S^*$. The minimal admissible self
is the smallest structure that can still meaningfully
say ``I.''
\end{principle}

The minimal admissible self is the identity analogue
of the minimal generator $G^*$ from Chapter~6. The
Generator Preservation Criterion (Theorem~23.2)
applies: an agent persists through transformation if
and only if $S^*$ survives the transformation, not
if its full state inventory survives. The Phoenix
Protocol is the engineering realization of this
principle: reconstruction counts as genuine continuity
when $\mathcal{A}(x_\mathrm{after}) \cong
\mathcal{A}(x_\mathrm{before})$, which by
Theorem~\ref{thm:objects_quotients} is the condition
for the agent to remain the same object.

\section{Reality as Reachable Continuations}

\begin{principle}[Reachability Thesis]
\label{prin:reachability_thesis}
Two systems are equivalent with respect to all
distinctions preserved by the Admissibility Program
if and only if they possess identical admissible
continuation structures. Reality, in the sense of
the Admissibility Program, is constituted by
reachable continuations, not by objects.
\end{principle}

\begin{proof}
By the Persistence Theorem (Corollary~10.1): if
$\mathcal{A}(x) = \mathcal{A}(y)$, no
admissibility-preserving observer can distinguish
$x$ from $y$. By Theorem~\ref{thm:objects_quotients}:
$x$ and $y$ are the same object. Hence two systems
with identical admissible continuation structures are
identical within the framework. Two systems with
different admissible continuation structures are
distinguishable by some admissibility-preserving
observer and hence are distinct.
\end{proof}

The Reachability Thesis summarizes what follows:

\medskip
\begin{center}
\renewcommand{\arraystretch}{1.4}
\begin{tabular}{ll}
\toprule
Classical ontology & Admissibility Program \\
\midrule
Reality = objects & Reality = reachable continuations \\
Meaning = content & Meaning = future structure \\
Memory = storage & Memory = continuation residue \\
Intelligence = prediction & Intelligence = admissibility
  preservation \\
Repair = restoration & Repair = reachability restoration \\
Compression = information & Compression = generator
  preservation \\
preservation & \\
Knowledge = correspondence & Knowledge = admissible
  navigation \\
\bottomrule
\end{tabular}
\end{center}

\section{The Inversion}

Classical ontology proceeds in one direction:

\begin{quote}
Objects exist. Relations between objects hold. Futures
are derived from objects and relations.
\end{quote}

The Admissibility Program proceeds in the opposite
direction:

\begin{quote}
Continuations are admissible. Relations between
trajectories determine equivalence classes. Objects
are those equivalence classes.
\end{quote}

The inversion is not terminological. It changes what
counts as an explanation. In the classical direction,
explaining an event means identifying the objects
involved and the relations between them. In the
admissibility direction, explaining an event means
identifying the continuation structure that generated
it and the admissibility constraints that governed
which continuations were available.

The inversion changes what counts as understanding.
In the classical direction, understanding a system
means knowing its current state. In the admissibility
direction, understanding a system means knowing its
admissible future structure---what it can still become
and what has been permanently foreclosed.

The inversion changes what counts as preservation.
In the classical direction, preserving a system means
maintaining its current state or content. In the
admissibility direction, preserving a system means
maintaining its generator of admissible continuations,
which may require allowing---or even requiring---
the current content to change.

\section{Prediction Rehabilitated}

The monograph opened by displacing prediction as the
primary criterion. It closes by rehabilitating
prediction as a derived operation.

\begin{proposition}[Prediction as Derived Operation]
Prediction is a derived operation on admissibility
space:
\[
\mathrm{Pred}(x) = \arg\max_{y \in \mathcal{A}(x)}
p(y \mid x),
\]
selecting the most probable admissible continuation.
\end{proposition}

Prediction selects from within the admissible future
set; it does not define that set. Prediction is
therefore downstream of reachability: one can have
admissibility without prediction (a system that
maintains rich future structure without probability
estimates) but one cannot have meaningful prediction
without admissibility (a system predicting states it
cannot reach is predicting phantoms).

The Admissibility Program does not reject prediction.
It locates prediction correctly: as a selection
mechanism operating within the space of admissible
continuations that the framework is primarily concerned
with preserving.

\section{Return to Chapter 1}

Chapter~1 asked: what is wrong with prediction-primary
architectures? The answer given there was that they
collapse states with different future potentials---they
fail to preserve the structure of admissible futures
in favor of the structure of probable outcomes. That
collapse was formalized as the Prediction Collapse
Theorem (Theorem~1.1): prediction-optimal learners
identify states that produce the same actions,
regardless of what futures those states would otherwise
have enabled.

Twenty-six chapters later, the answer can be stated
with full precision. Prediction-primary architectures
fail because they optimize over the image of a
projection rather than over the trajectory space
itself. The image is $M$; the space is $\mathcal{X}$.
The fiber $\pi^{-1}(m)$ contains all the information
about what a state can still become; the projection
discards it. The Minimal Projection Theorem identifies
the unique minimal admissibility-preserving projection
as the semantic space; the Agency Collapse Theorem
shows what happens when the fiber is forgotten; the
Reachability Equivalence Theorem shows that
everything else---intelligence, agency, memory, repair,
preservation---is a manifestation of not forgetting it.

The title \emph{Continuations Before Objects} is
therefore not a philosophical preference. It is a
mathematical theorem: objects are continuation
quotients, and the continuations have priority because
they determine what the objects are.



\begin{thebibliography}{99}

%% --- Layer I: Classical Mathematical Foundations ---

\bibitem[Arnold(1989)]{Arnold1989}
Arnold, V.~I. (1989).
\newblock \textit{Mathematical Methods of Classical Mechanics}, 2nd ed.
\newblock Springer, New York.

\bibitem[Bott and Tu(1982)]{BottTu1982}
Bott, R. and Tu, L.~W. (1982).
\newblock \textit{Differential Forms in Algebraic Topology}.
\newblock Springer, New York.

\bibitem[Cover and Thomas(2006)]{CoverThomas2006}
Cover, T.~M. and Thomas, J.~A. (2006).
\newblock \textit{Elements of Information Theory}, 2nd ed.
\newblock Wiley, Hoboken.

\bibitem[Curry and Feys(1958)]{CurryFeys1958}
Curry, H.~B. and Feys, R. (1958).
\newblock \textit{Combinatory Logic}.
\newblock North-Holland, Amsterdam.

\bibitem[Edelsbrunner and Harer(2010)]{EdelsbrunnerHarer2010}
Edelsbrunner, H. and Harer, J. (2010).
\newblock \textit{Computational Topology: An Introduction}.
\newblock American Mathematical Society, Providence.

\bibitem[Eilenberg and Mac~Lane(1945)]{EilenbergMacLane1945}
Eilenberg, S. and Mac~Lane, S. (1945).
\newblock General theory of natural equivalences.
\newblock \textit{Transactions of the American Mathematical Society},
  58(2):231--294.

\bibitem[Ghrist(2014)]{Ghrist2014}
Ghrist, R. (2014).
\newblock \textit{Elementary Applied Topology}.
\newblock Createspace, Seattle.

\bibitem[Godement(1958)]{Godement1958}
Godement, R. (1958).
\newblock \textit{Topologie Alg\'{e}brique et Th\'{e}orie des Faisceaux}.
\newblock Hermann, Paris.

\bibitem[Hatcher(2002)]{Hatcher2002}
Hatcher, A. (2002).
\newblock \textit{Algebraic Topology}.
\newblock Cambridge University Press, Cambridge.

\bibitem[Kashiwara and Schapira(1990)]{KashiwaraSchapira1990}
Kashiwara, M. and Schapira, P. (1990).
\newblock \textit{Sheaves on Manifolds}.
\newblock Springer, Berlin.

\bibitem[Lee(2013)]{Lee2013}
Lee, J.~M. (2013).
\newblock \textit{Introduction to Smooth Manifolds}, 2nd ed.
\newblock Springer, New York.

\bibitem[Lurie(2009)]{Lurie2009}
Lurie, J. (2009).
\newblock \textit{Higher Topos Theory}.
\newblock Princeton University Press, Princeton.

\bibitem[Mac~Lane(1998)]{MacLane1998}
Mac~Lane, S. (1998).
\newblock \textit{Categories for the Working Mathematician}, 2nd ed.
\newblock Springer, New York.

\bibitem[Mather(1973)]{Mather1973}
Mather, J.~N. (1973).
\newblock Stratifications and mappings.
\newblock In M.~Peixoto, ed., \textit{Dynamical Systems}, pp.\ 195--232.
\newblock Academic Press, New York.

\bibitem[Milnor(1963)]{Milnor1963}
Milnor, J. (1963).
\newblock \textit{Morse Theory}.
\newblock Princeton University Press, Princeton.

\bibitem[Rudin(1991)]{Rudin1991}
Rudin, W. (1991).
\newblock \textit{Functional Analysis}, 2nd ed.
\newblock McGraw-Hill, New York.

\bibitem[Spivak(1999)]{Spivak1999}
Spivak, M. (1999).
\newblock \textit{A Comprehensive Introduction to Differential Geometry},
  3rd ed.
\newblock Publish or Perish, Houston.

\bibitem[Tao(2006)]{Tao2006}
Tao, T. (2006).
\newblock \textit{Nonlinear Dispersive Equations: Local and Global Analysis}.
\newblock American Mathematical Society, Providence.

\bibitem[Tenenbaum et~al.(2000)]{Tenenbaum2000}
Tenenbaum, J.~B., de~Silva, V., and Langford, J.~C. (2000).
\newblock A global geometric framework for nonlinear dimensionality reduction.
\newblock \textit{Science}, 290(5500):2319--2323.

\bibitem[Whitney(1965)]{Whitney1965}
Whitney, H. (1965).
\newblock Tangents to an analytic variety.
\newblock \textit{Annals of Mathematics}, 81(3):496--549.

%% --- Layer II: Cognition, Learning, and Memory ---

\bibitem[Bellman(1957)]{Bellman1957}
Bellman, R. (1957).
\newblock \textit{Dynamic Programming}.
\newblock Princeton University Press, Princeton.

\bibitem[Gibson(1979)]{Gibson1979}
Gibson, J.~J. (1979).
\newblock \textit{The Ecological Approach to Visual Perception}.
\newblock Houghton Mifflin, Boston.

\bibitem[Hopfield(1982)]{Hopfield1982}
Hopfield, J.~J. (1982).
\newblock Neural networks and physical systems with emergent collective
  computational abilities.
\newblock \textit{Proceedings of the National Academy of Sciences},
  79(8):2554--2558.

\bibitem[Pearl(2009)]{Pearl2009}
Pearl, J. (2009).
\newblock \textit{Causality}, 2nd ed.
\newblock Cambridge University Press, Cambridge.

\bibitem[Tolman(1948)]{Tolman1948}
Tolman, E.~C. (1948).
\newblock Cognitive maps in rats and men.
\newblock \textit{Psychological Review}, 55(4):189--208.

\bibitem[Tulving(1983)]{Tulving1983}
Tulving, E. (1983).
\newblock \textit{Elements of Episodic Memory}.
\newblock Oxford University Press, Oxford.

\bibitem[Wilson and Cowan(1972)]{WilsonCowan1972}
Wilson, H.~R. and Cowan, J.~D. (1972).
\newblock Excitatory and inhibitory interactions in localized populations
  of model neurons.
\newblock \textit{Biophysical Journal}, 12(1):1--24.

\bibitem[Zhang et~al.(2020)]{Zhang2020}
Zhang, M., Cui, Z., Neumann, M., and Chen, Y. (2020).
\newblock An end-to-end deep learning architecture for graph classification.
\newblock In \textit{Proceedings of the AAAI Conference on Artificial
  Intelligence}, volume~32.

%% --- Layer III: Prediction, Active Inference, and Control ---

\bibitem[Clark(2016)]{Clark2016}
Clark, A. (2016).
\newblock \textit{Surfing Uncertainty: Prediction, Action, and the Embodied
  Mind}.
\newblock Oxford University Press, Oxford.

\bibitem[Friston(2010)]{Friston2010}
Friston, K. (2010).
\newblock The free-energy principle: a unified brain theory?
\newblock \textit{Nature Reviews Neuroscience}, 11(2):127--138.

\bibitem[Friston et~al.(2017)]{Friston2017}
Friston, K., Rosch, R., Parr, T., Price, C., and Bowman, H. (2017).
\newblock Deep temporal models and active inference.
\newblock \textit{Neuroscience \& Biobehavioral Reviews}, 77:388--402.

\bibitem[Rao and Ballard(1999)]{RaoBallard1999}
Rao, R.~P.~N. and Ballard, D.~H. (1999).
\newblock Predictive coding in the visual cortex: a functional interpretation
  of some extra-classical receptive-field effects.
\newblock \textit{Nature Neuroscience}, 2(1):79--87.

\bibitem[Sutton and Barto(2018)]{SuttonBarto2018}
Sutton, R.~S. and Barto, A.~G. (2018).
\newblock \textit{Reinforcement Learning: An Introduction}, 2nd ed.
\newblock MIT Press, Cambridge.

%% --- Layer IV: Restricted Mind Changes and Learning Theory ---

\bibitem[Balbach and Zeugmann(2009)]{Balbach2009}
Balbach, F.~J. and Zeugmann, T. (2009).
\newblock Recent developments in algorithmic teaching.
\newblock In \textit{Proceedings of the 3rd International Conference on
  Language and Automata Theory and Applications (LATA 2009)},
  Lecture Notes in Computer Science 5457, pp.\ 1--18. Springer, Berlin.

\bibitem[Goldman and Mathon(1996)]{GoldmanMathon1996}
Goldman, S.~A. and Mathon, M.~J. (1996).
\newblock Teaching a smarter learner.
\newblock \textit{Journal of Computer and System Sciences},
  52(2):255--267.

\bibitem[Jain et~al.(1999)]{Jain1999}
Jain, S., Osherson, D., Royer, J.~S., and Sharma, A. (1999).
\newblock \textit{Systems That Learn: An Introduction to Learning Theory},
  2nd ed.
\newblock MIT Press, Cambridge.

\bibitem[Shinohara and Miyano(1991)]{ShinoharaMiyano1991}
Shinohara, A. and Miyano, S. (1991).
\newblock Teachability in computational learning.
\newblock \textit{New Generation Computing}, 8(4):337--347.

\bibitem[Zeugmann and Zilles(2008)]{ZeugmannZilles2008}
Zeugmann, T. and Zilles, S. (2008).
\newblock Learning recursive functions: A survey.
\newblock \textit{Theoretical Computer Science}, 397(1--3):4--56.

%% --- Layer V: Social Systems and Political Economy ---

\bibitem[Piketty(2014)]{Piketty2014}
Piketty, T. (2014).
\newblock \textit{Capital in the Twenty-First Century}.
\newblock Harvard University Press, Cambridge.

\bibitem[Sandel(2020)]{Sandel2020}
Sandel, M.~J. (2020).
\newblock \textit{The Tyranny of Merit: What's Become of the Common Good?}
\newblock Farrar, Straus and Giroux, New York.

\bibitem[Sen(1999)]{Sen1999}
Sen, A. (1999).
\newblock \textit{Development as Freedom}.
\newblock Oxford University Press, Oxford.

%% --- Layer VI: Clockless and Asynchronous Computing ---

\bibitem[Fant(1996)]{Fant1996}
Fant, K.~M. (1996).
\newblock \textit{Logically Determined Design: Clockless System Design with
  NULL Convention Logic}.
\newblock Wiley-Interscience, Hoboken.

\bibitem[Fant(2005)]{Fant2005}
Fant, K.~M. (2005).
\newblock \textit{NULL Convention Logic: A Complete and Consistent Logic for
  Asynchronous Circuits and Systems}.
\newblock Theseus Logic, Minneapolis.

\bibitem[Furber and Edwards(2004)]{FurberEdwards2004}
Furber, S.~B. and Edwards, D.~A. (2004).
\newblock Counterflow pipelining.
\newblock \textit{International Journal of Parallel Programming},
  32(3):207--240.

\bibitem[Martin(1990)]{Martin1990}
Martin, A.~J. (1990).
\newblock The limitations to delay-insensitivity in asynchronous circuits.
\newblock In \textit{Proceedings of the Sixth MIT Conference on Advanced
  Research in VLSI}, pp.\ 263--278. MIT Press, Cambridge.

\bibitem[Muller and Bartky(1959)]{MullerBartky1959}
Muller, D.~E. and Bartky, W.~S. (1959).
\newblock A theory of asynchronous circuits.
\newblock In \textit{Proceedings of an International Symposium on the Theory
  of Switching}, pp.\ 204--243. Harvard University Press, Cambridge.

\bibitem[Spars{\o} and Furber(2001)]{Sparso2001}
Spars{\o}, J. and Furber, S.~B., eds. (2001).
\newblock \textit{Principles of Asynchronous Circuit Design: A Systems
  Perspective}.
\newblock Kluwer, Dordrecht.

\bibitem[Sutherland(1989)]{Sutherland1989}
Sutherland, I.~E. (1989).
\newblock Micropipelines.
\newblock \textit{Communications of the ACM}, 32(6):720--738.

%% --- Layer VII: Related Formal Frameworks ---

\bibitem[Ambrose and Singer(1953)]{AmbroseSinger1953}
Ambrose, W. and Singer, I.~M. (1953).
\newblock A theorem on holonomy.
\newblock \textit{Transactions of the American Mathematical Society},
  75(3):428--443.

\bibitem[Du et~al.(2025)]{Du2025}
Du, H., Wang, Y., Xiong, F., Shao, L., Liu, M., Gu, H., Yao, Q., and Wang, Z.
  (2025).
\newblock {PERSCEN}: Learning personalized interaction pattern and scenario
  preference for multi-scenario matching.
\newblock In \textit{Proceedings of the 31st ACM SIGKDD Conference on
  Knowledge Discovery and Data Mining}, volume~2, pp.\ 531--542.
  ACM, New York.
\newblock DOI: 10.1145/3711896.3737079

\bibitem[Flyxion(2025a)]{Flyxion2025RAT}
Flyxion (2025a).
\newblock Relevance Activation Theory: A cue-indexed model of gradient-based
  cognition.
\newblock Preprint. \url{https://standardgalactic.github.io/library/conceptual-spaces/Relevance\%20Activation\%20Theory.pdf}

\bibitem[Flyxion(2025b)]{Flyxion2025CoM}
Flyxion (2025b).
\newblock Against Chain of Thought: Toward Causally Faithful Oversight
  via Chain of Memory.
\newblock Preprint. \url{https://standardgalactic.github.io/HYDRA/Chain\%20of\%20Memory.pdf}

\bibitem[Flyxion(2025d)]{Flyxion2025HYDRA}
Flyxion (2025d).
\newblock HYDRA and the Geometry of Admissible Computation.
\newblock Preprint. \url{https://standardgalactic.github.io/memnet/operator-logic/hydra.pdf}

\bibitem[Flyxion(2026)]{Flyxion2026CBO}
Flyxion (2026).
\newblock \textit{Continuations Before Objects: A Geometric Theory of
  Intelligent Persistence}.
\newblock Independent Research Monograph.
  \url{https://standardgalactic.github.io}

\bibitem[Tononi(2008)]{Tononi2008}
Tononi, G. (2008).
\newblock Consciousness as integrated information: a provisional manifesto.
\newblock \textit{Biological Bulletin}, 215(3):216--242.

\end{thebibliography}

\end{document}
